%\documentclass[aps,prl,singlecolumn,superscriptaddress]{revtex4-2}


\documentclass[a4paper,11pt]{article}
\usepackage{./jheppub}
%\setcounter{secnumdepth}{3}
% Useful packages
\usepackage{tikz-cd}
\usepackage{amsmath}
\usepackage{array,booktabs}
\usepackage{hyperref}
%%\usepackage{showlabels}
\usepackage{amsthm}
\newtheorem{theorem}{Theorem}
\newtheorem{lemma}[theorem]{Lemma}
\newtheorem{corollary}[theorem]{Corollary}
\newtheorem{proposition}[theorem]{Proposition}
\theoremstyle{definition}
\newtheorem{definition}{Definition}[section]
\newtheorem{prop}{Conjecture}
\newtheorem{cor}{Corollary}
%-------------------packages-------------------------
\usepackage{enumerate}
\usepackage[utf8]{inputenc}
\usepackage{amssymb,amsbsy,mathtools,amsmath}
\usepackage{braket}
\usepackage{graphicx}
\usepackage{xcolor}
%\usepackage{multirow}
\usepackage{hyperref}
%\usepackage{caption}
%\usepackage{subcaption}
%\mathcal{A}ptionsetup{compatibility=false}
\usepackage{natbib}
%\usepackage{dcolumn}% Align table columns on decimal point
\usepackage{bm}% bold math
\usepackage{soul}
%\usepackage{showlabels}
%\newtheorem{proposition}{Proposition}
\newcommand{\E}{\mathbb{E}}
\newcommand{\dd}{\,\mathrm{d}}
%\newcommand{\mathcal M}{\mathcal M}
\newcommand{\cH}{\mathcal{H}}
\newcommand{\KLL}{D_{\mathrm{KL}}}
\newcommand{\Mop}{\mathsf{M}}

%Includes "References" in the table of contents
\usepackage[nottoc]{tocbibind}

\usepackage{bbm}
\usepackage{amsthm}
%\newtheorem{theorem}{Theorem}
%\newtheorem{lemma}[theorem]{Lemma}
%\newtheorem{corollary}[theorem]{Corollary}
%\theoremstyle{definition}
%\newtheorem{definition}{Definition}[section]
\usepackage{amssymb}
\newcommand{\id}{\mathbb{I}}
\newcommand{\Tr}{\text{Tr}}
\newcommand{\Div}{\operatorname{div}}
\usepackage{amsmath}
%\usepackage{tikz}
\usepackage{braket}
%\usepackage{quantikz}
\usepackage{graphicx}
\usepackage{caption}
\newcommand{\mA}{\mathcal{A}}% no longer used in the body: algebras are \mathcal N\subset\mathcal M, nets \mathcal A(O)
\newcommand{\mB}{\mathcal{B}}% no longer used in the body
\newcommand{\mP}{\mathcal{P}}
\newcommand{\mR}{\mathcal{R}}
\newcommand{\nn}{\nonumber}
\newcommand{\mN}{\mathcal{N}}
\newcommand{\mL}{\mathcal{L}}
\newcommand{\lb}{\left(}
\newcommand{\rb}{\right)}
\newcommand{\mH}{\mathcal{H}}
\usepackage{subcaption}
\newcommand{\tOO}{\tilde{\mathcal{O}}}
\newcommand{\mO}{\mathcal{O}}
\newcommand{\mK}{\mathcal{K}}
\newcommand{\tH}{\tilde{H}}
\newcommand{\ep}{\epsilon}
\newcommand{\mbR}{\mathbb{R}}
\newcommand{\mM}{\mathcal{M}}
\newcommand{\mT}{\mathcal{T}}
\newcommand{\NL}[1]{\textcolor{blue}{#1}}
\newcommand{\p}{\partial}
\newcommand{\tr}{\text{tr}}
\newcommand{\tT}{\tilde{T}}
\newcommand{\Ad}{\operatorname{Ad}}
\newcommand{\Dom}{\operatorname{Dom}}
\newcommand{\mE}{\mathcal{E}}
\newcommand{\relcomm}[2]{#1'\cap #2}

\newcommand{\SO}[1]{\textcolor{purple}{#1}}
%\newcommand{\MM}[1]{\textcolor{red}{#1}}
\newcommand{\KL}[1]{\textcolor{brown}{#1}}
\newcommand{\YC}[1]{\textcolor{teal}{#1}}
\newcommand{\cL}{\mathcal{L}}
\setlength{\emergencystretch}{2em}
\begin{document}

% Title
\title{Subsystem Flows: From Quantum Error Correction to Modular Chaos}

\author{Yidong Chen$^1$,}
\affiliation{$^1$Department of Mathematics, University of Illinois at Urbana-Champaign, Illinois, IL 61801, USA}
\author{Nima Lashkari$^2$,}
\author{Kwing Lam Leung$^2$}
\affiliation{$^2$Department of Physics and Astronomy, Purdue University, West Lafayette, IN 47907, USA}
\emailAdd{yidongc2@illinois.edu}
\emailAdd{nima@purdue.edu}
\emailAdd{leung60@purdue.edu}
\abstract{
An algebraic coarse-graining retains a subalgebra of observables $\mathcal N\subset\mathcal M$ and erases the rest. For dynamics generated by a faithful normal reference state, we relate autonomy of the retained observables to exact quantum error correction of this restriction. Takesaki's theorem identifies modular autonomy with a reference-preserving normal conditional expectation, which defines an exactly recoverable code. We express this condition as triviality of a compressed subsystem flow. When the reference vector is cyclic and separating for both algebras, this comparison is the unitary cocycle $D_{N\subset M}(t)=\Delta_M^{-it}\Delta_N^{it}$. We show that its infinitesimal expectation is the sum of two nonnegative relative-entropy losses, for the inclusion and its commutant dual.

Motivated by classical ergodic theory and state-generated Hamiltonian flows, we organize modular evolution through the history and memory of a subsystem: the smallest autonomous algebra containing it and the largest autonomous algebra inside it. Common cyclicity implies generation, and standardness with respect to the reference gives trivial memory for inclusions of factors. Half-sided modular inclusions supply positive translations and exact modular growth with exponent $2\pi$, giving the translation energy an interpretation in terms of recoverability. At imaginary modular time, $L^2$-nuclearity produces a positive inclusion Hamiltonian with discrete spectrum and a finite partition function. We establish monotonicity under ordered collar enlargements and an intrinsic real structure, recover the cylinder spectrum for concentric vacuum balls, and formulate a GOE conjecture for the averaged, connected spectral form factor of deformed vacuum inclusions after unfolding at high energy.
}


% \abstract{An algebraic coarse-graining retains a subalgebra of macro-observables $\mathcal N\subset \mathcal M$ and erases the rest. When a dynamical law preserves the subalgebra, it induces a closed effective dynamics on the coarse-grained observables. We call a subalgebra with trivial center a {\it subsystem} (subfactor). 
% Given a state on $\mathcal M$, its modular flow preserves a subsystem $\mathcal N$ if and only if it is the {\it subsystem flow} $D_{N\subset M}(t)=\Delta_M^{-it}\Delta_N^{it}$ that acts trivially on every coarse-grained observable. Such states have no correlations between the degrees of freedom in $\mathcal N$ and those in the relative commutant. This factorization is precisely the condition for exact Quantum Error Correction (QEC): the subalgebra $\mathcal N$ is correctable with respect to the restriction to $\mathcal N$: (erasure of the relative commutant). Subsystem flow measures the failure of correctability and the leakage of information from $\mathcal N$ into the relative commutant degrees of freedom (environment).

% In this work, we first motivate the notion of subsystem flow using the classical ergodic theory, with Hamiltonian flows playing the role of modular flows. Then, we generalize to quantum subsystem flows, translating key notions of the ergodic hierarchy such as ergodicity, mixing, generation, forgetfulness, and K-systems into the language of QEC. At the opposite end of the hierarchy from correctability are chaotic modular flows of infinitely entangled states that ensure any information initially localized in $B$ mixes strongly and is eventually lost. We argue that this requires the inclusion of type III$_1$ factors similar to those of subregions of quantum field theory. Half-sided modular inclusions provide examples of maximally chaotic modular flows. We argue that the inclusions with $L^2$-nuclearity property can be used to define a modular spectral form factor.}

% When the dynamics preserves the subalgebra, it induces a closed effective dynamics on the coarse-grained observables. 
% For modular dynamics, this occurs if and only if the subalgebra flow $D_{N\subset M}(t)=\Delta_M^{-it}\Delta_N^{it}$ acts trivially on every coarse-grained observable. 

% , which compares the modular flow of a faithful normal state on an algebra $\mathcal M$ with that of its restriction to a subalgebra $\mathcal N$; its classical avatar is the Poisson flow generated by the relative Hamiltonian $-\log(\rho/\rho_0)$, whose expectation is the Kullback--Leibler divergence. By Takesaki's theorem the subsystem flow is trivial exactly when a state-preserving conditional expectation exists, and we identify this with exact quantum error correction of the erasure of the complement of $\mathcal N$. In general, the generator of the subsystem flow evaluated on a state is the sum of two Petz sufficiency gaps, so the subsystem flow measures the failure of correctability. 

\maketitle

% %\tableofcontents

% \section{Introduction}\label{sec:intro}

% In this paper, 

% Invariant subsystem is QEC code, the relative commutant is 

% \NL{What can we say about holography?}

% In the Schrödinger picture, a general coarse-graining of a finite quantum system is described by a quantum channel (a trace-preserving Completely Positive (CP) linear map) acting on density matrices as $\rho\to \Phi(\rho)$. We say 


% subalgebra of macro-observables $\mathcal N\subset \mathcal M$ and discards the information needed to distinguish states that agree on $\mathcal N$. The passage from a reversible microscopic description to an apparently irreversible macroscopic one is often attributed to the information lost in this way. Ergodic theory is primarily concerned with the emergence of this irreversibility and information loss. At the opposite end of the ergodic hierarchy, one can imagine 


% This raises two related questions: 
% \begin{itemize}
%     \item Can the discarded information be recovered? 
%     \item Do the retained observables evolve autonomously?
% \end{itemize} 
% In this paper, we use the operator algebraic theory of Quantum Error Correction (QEC) to relate these two questions for dynamics generated by a quantum state; i.e., modular dynamics. The central object is the {\it subsystem flow}, which compares the modular dynamics of an algebra
% with the intrinsic modular dynamics of a subalgebra. Its triviality characterizes modular autonomy and the existence of a reference-state-preserving encoding that is exactly correctable against the error of restriction to the subalgebra. Its infinitesimal generator measures two
% losses of relative entropy under restriction.


% In standard quantum information theory, in the Schr\"odinger picture, a noise map is modeled as a quantum channel: $\rho_B\to\Phi(\rho_B)$, a completely positive (CP) trace-preserving map. A simple example of such a noise map is
% \begin{equation}
%     \Phi(\rho_B)=\text{tr}_R\left(U_{BR}(\rho_B\otimes|0\rangle\langle0|_R)U_{BR}^\dagger\right).
% \end{equation}
% Here $R$ describes the environment degrees of freedom, initially in a fixed state $|0\rangle_R$ uncorrelated with the system, and $U_{BR}$ describes their interaction. The isometric encoding
% \begin{equation}
%     |\psi_i\rangle=U_{BR}\bigl(|i\rangle_N\otimes|0\rangle_R\bigr)\in\mathcal H_{BR}
% \end{equation}
% generates an encoded subspace, with erasure of $R$ as the noise. The Stinespring dilation shows that any quantum channel can be realized by an isometry followed by erasure of auxiliary environment degrees of freedom. It does not, by itself, imply that the resulting noise is correctable on the encoded subspace: a recovery must act on the surviving subsystem, whereas reversing $U_{BR}$ requires access to the full system. We write $BR$ for the composite system.

% In the Heisenberg picture, there is a natural embedding of the algebra of $B$ inside that of $BR$, i.e. $\iota(\mathcal N)\subset\mathcal M$, given by $\iota(b)=b\otimes1_R$. Writing $\alpha_U(a)=U_{BR}^\dagger aU_{BR}$, we have
% \begin{align}
%     \forall b\in\mathcal N:\qquad \iota_U(b)&=\alpha_U\circ\iota(b)=U_{BR}^\dagger(b\otimes1_R)U_{BR},\\
%     \text{tr}\bigl(\Phi(\rho_B)b\bigr)&=\text{tr}\bigl(\rho_{BR}\iota_U(b)\bigr),
%     \qquad \rho_{BR}=\rho_B\otimes|0\rangle\langle0|_R.
% \end{align}
% The Heisenberg dual of the channel includes the expectation over the initial environment state:
% \begin{equation}
%     \Phi^\dagger=P_0\circ\iota_U,\qquad
%     P_0(b\otimes r)=b\langle0|r|0\rangle_R,\qquad
%     P_0\circ\iota=\operatorname{id}_N.
% \end{equation}
% The composition $\mathcal E_0=\iota\circ P_0$ is a normal conditional expectation onto $\iota(\mathcal N)$. Its predual encodes states as $\rho_B\mapsto\rho_B\otimes|0\rangle\langle0|_R$, a family on which erasure of $R$ is exactly correctable. This gives a simple example of the expectation-based codes considered below. When both $\mathcal N$ and $\mathcal M$ have trivial center, we call $\mathcal N\subset\mathcal M$ a {\it subsystem} inclusion, or subfactor. The resulting QEC code is called a {\it subsystem code}.


% In the finite-dimensional tensor-factor setting, $\mathcal E(b\otimes r)=b\,\text{tr}(\rho_Rr)$ for a fixed environment density matrix $\rho_R$, and the code consists of states $\rho_B\otimes\rho_R$. Thus a faithful reference density matrix $\rho_{BR}$ is preserved precisely when it is a product, $\rho_{BR}=\rho_B\otimes\rho_R$. For a finite-dimensional subalgebra with a center, there is instead a sector decomposition $\mathcal H_{BR}=\oplus_i\mH_B^i\otimes\mH_R^i$, and the states preserved by the expectation have the block-product form $\rho_{BR}=\oplus_i p_i\rho_B^{(i)}\otimes\rho_R^{(i)}$, with $\rho_R^{(i)}$ fixed in each sector. These statements describe the code states, rather than arbitrary states of the ambient algebra. In the tensor-factor case, preservation of the reference state is equivalent to
% \begin{align}
%     \forall b\in\mathcal N,\ t\in\mathbb R:\qquad
%     \rho_{BR}^{it}(b\otimes1_R)\rho_{BR}^{-it}
%     =(\rho_B^{it}b\rho_B^{-it})\otimes1_R.
% \end{align}
% We refer to this as the {\it Takesaki condition}. The corresponding comparison of the two matrix evolutions is implemented by
% \begin{align}
%     &d_{N\subset M}(t)=\rho_{BR}^{-it}(\rho_B^{it}\otimes1_R),\nn\\
%     &d_{N\subset M}(t)(b\otimes1_R)d_{N\subset M}(t)^\dagger=b\otimes1_R.
% \end{align}
% We use a lowercase $d$ here to distinguish this matrix implementer from the modular operator comparison defined below.

% In physics, the initial state of an open quantum system can contain correlations between $B$ and $R$. Without a specified assignment of subsystem states to global states, its reduced evolution need not define a CP map on all states of $B$. For a faithful correlated $\rho_{BR}$ in the tensor-factor setting, the full modular evolution does not preserve $\mathcal N$. At time $t$, the comparison above moves the retained algebra to
% \begin{align}
%     \mathcal N_t:=d_{N\subset M}(t)\mathcal N d_{N\subset M}(t)^\dagger
%     =\rho_{BR}^{-it}\mathcal N\rho_{BR}^{it},
% \end{align}
% where we identify $\mathcal N$ with $\iota(\mathcal N)$. In the expectation-based code, $\mathcal N_t=\mathcal N$ for all modular time: the retained observables evolve autonomously. More generally, their modular orbit can spread through the global algebra, and its history and memory record which observables can be generated or remain accessible. For a specified noise channel and code, the information-disturbance relation identifies exact correctability with independence of the complementary output from the logical input (Section~\ref{subsec:QEC-KL}). This is a decoupling condition, rather than a requirement that the environment state be maximally mixed.

% The purpose of this paper is to show that, for dynamics generated by a state, correctability relative to that state is a statement about modular autonomy, and that the object which measures its failure is the \emph{subsystem flow}. Finite-dimensional modular dynamics has conserved observables and recurrences, which obstruct strict mixing of all connected correlations. Local type III$_1$ factors provide the infinite-dimensional setting relevant to quantum field theory. The type of the algebra alone, however, does not determine the ergodic properties of a chosen state or inclusion, nor does it assert that every state is entangled across an arbitrary pair of commuting subalgebras.

% We formulate the ergodic theory of a many-body system in terms of a net of subsystems and their dynamics. Ergodicity and mixing describe the global modular dynamics, while generation, forgetfulness, and half-sided nesting describe its action on a chosen subsystem. Together they organize the ways information can be retained or lost under modular evolution. This is what we mean by many-body ergodicity. We argue that this picture has an avatar in classical ergodic theory and in Hamiltonian flows on a phase space.

% For a general inclusion $\mathcal N\subset\mathcal M$, we define erasure as the restriction of (normal) states,
% \begin{equation}
%     \mathcal T:\mathcal M_*\to\mathcal N_*,
%     \qquad \mathcal T(\varphi)=\varphi|_N.
% \end{equation}
% If $\mathcal M\simeq\mathcal N\,\overline\otimes\,\mathcal R$, this is ordinary erasure of the factor $\mathcal R$, which is then the relative commutant $\mathcal R=\mathcal N'\cap\mathcal M$ of the inclusion\footnote{For a general inclusion, the relative commutant need not account for all discarded observables; see Section~\ref{subsec:QEC-KL}.}.

% A normal conditional expectation $\mathcal E:\mathcal M\to\mathcal N$ defines an encoding by its predual,
% \begin{equation}
%     \mathcal E_*(\varphi_N)=\varphi_N\circ\mathcal E,
%     \qquad \mathcal T\circ\mathcal E_*=\operatorname{id}_{\mathcal N_*}.
% \end{equation}
% The states preserved by $\mathcal E$ form the code,
% \begin{equation}
%     \mathcal C_{\mathcal E}
%     =\{\varphi\in\mathcal M_*:\varphi\text{ is a normal state and }\varphi\circ\mathcal E=\varphi\}.
% \end{equation}
% On this family, $\mathcal E_*\circ\mathcal T=\operatorname{id}_{\mathcal C_{\mathcal E}}$, so the same predual map serves as an encoding before erasure and a recovery map after erasure. We study codes for which $\mathcal E$ preserves a specified faithful normal reference state $\omega$.


% Both classically and quantum mechanically there are two qualitatively different coarse-grainings, which we call \emph{subspace} and \emph{subsystem}. The intuition for a subspace is a one-particle system in which the particle is either in one set of states or in another: a disjoint union of two sets with $d_1$ and $d_2$ states has $d_1+d_2$ states, whereas a pair of distinguishable particles with these state spaces has $d_1d_2$ states. States add for alternatives and multiply for subsystems; the corresponding quantum constructions are direct sums and tensor products. Many-body systems and field theories come with a net of subsystems labeled by spatial regions, and it is the subsystem case that concerns us.

% \paragraph{The classical avatar.}
% If $\rho$ is a positive probability density and $\rho_0$ its coarse-graining, the relative phase $(\rho/\rho_0)^{it}$ is a one-parameter unitary group whose generator has expectation equal to the Kullback--Leibler divergence $S(\rho\|\rho_0)$, the entropy cost of the erasure. Its conjugation action on the commutative observable algebra is trivial. A dynamical comparison becomes available once phase space carries a Poisson bracket: under the smoothness and completeness assumptions used below, the Hamiltonian $K_\rho=-\log\rho$ generates a flow that preserves the state and obeys an integration-by-parts analogue of the KMS condition. Composing the inverse coarse-grained motion with the full motion gives a \emph{relative transformation}, the classical subsystem flow. Its infinitesimal Hamiltonian is $K_\rho-K_{\rho_0}=-\log(\rho/\rho_0)$, whose expectation is $-S(\rho\|\rho_0)$. In general the comparison is a cocycle, rather than the autonomous flow of this Hamiltonian difference. It is trivial on the retained observables precisely when their full evolution agrees with the reduced state-generated evolution; for the connected product symplectic spaces considered here, this happens exactly for product densities. Classically, however, every measurable coarse-graining admits a conditional expectation preserving the reference probability measure, whose predual defines a correctable family of extensions. This holds independently of Poisson autonomy for a different density. The Kullback--Leibler divergence is the sufficiency gap for the restriction of the state and the reference to the retained observables.

% \paragraph{The quantum subsystem flow.}
% In the quantum theory the state-generated flow is the modular automorphism group $\sigma^{\omega;M}_t=\operatorname{Ad}\Delta_{\omega;M}^{it}$ of a faithful normal state $\omega$ on a von Neumann algebra $\mathcal M$, the unique flow for which $\omega$ is KMS with the corresponding time convention. We write $K_M=-\log\Delta_{\omega;M}$, so that evolution generated by $K_M$ is $\sigma^{\omega;M}_{-t}$. For an inclusion $\mathcal N\subset\mathcal M$ with a common cyclic and separating vector representing $\omega$, we define
% \begin{equation}
%  D_{N\subset M}(t)=\Delta_{\omega;M}^{-it}\Delta_{\omega;N}^{it},
%  \qquad
%  \operatorname{Ad}D_{N\subset M}(t)(b)=\sigma^{\omega;M}_{-t}\bigl(\sigma^{\omega;N}_{t}(b)\bigr),
%  \label{eq:intro-subsystem-flow}
% \end{equation}
% which evolves backward by the intrinsic $K_N$ evolution and then forward by the outer $K_M$ evolution. As in the classical comparison, $D(t)$ is generally a cocycle. In the algebraic quantum field theory literature this is the characteristic function of the inclusion~\cite{revolutionizing}; the name \emph{subsystem flow} emphasizes its origin in the comparison of state-generated dynamics. It is analytic in the open strip $0<\operatorname{Im}z<1/2$, with contraction boundary values. The estimate follows from the quadratic-form order $\Delta_N\ge\Delta_M$ of modular operators under inclusion, which in turn follows from the Tomita operator of $\mathcal M$ extending that of $\mathcal N$. In particular, the closure of $D(i\alpha)=\Delta_M^\alpha\Delta_N^{-\alpha}$ is a contraction for $0\le\alpha\le1/2$. This order relation has no counterpart for the classical Poisson Hamiltonians. These analytic properties underlie the modular chaos bound used in holography~\cite{deBoerLamprou2020}, and subsystem flows enter the algebraic reconstruction of bulk spacetime from the net of boundary subalgebras at $G_N=0$.

% \paragraph{Takesaki's theorem as an error-correction statement.}
% The decisive difference between the classical and quantum settings is Takesaki's theorem~\cite{Takesaki1972}: a normal conditional expectation onto $\mathcal N$ preserving a faithful normal state $\omega$ exists if and only if $\mathcal N$ is invariant under the modular flow of $\omega$. Its predual encodes logical states on $\mathcal N$ into states on $\mathcal M$, and the states preserved by the expectation form a family that can be recovered after restriction to $\mathcal N$. Takesaki's theorem thus says that \emph{the subalgebras correctable in this reference-preserving sense are the modularly autonomous ones}. For proper subregion inclusions with a common cyclic reference vector, such expectations do not exist, and the subsystem flow is nontrivial. Four results organize the paper.

% First, \emph{correctability is triviality of the subsystem flow} (Theorem~\ref{thm:correctable}). To include cases in which $\Omega$ is not cyclic for $\mathcal N$ on the full Hilbert space, let $\mathcal H_N=\overline{\mathcal N\Omega}$ and let $V:\mathcal H_N\hookrightarrow\mathcal H$ be the inclusion. The intrinsic modular operator of $\mathcal N$ acts on $\mathcal H_N$, and we use the compressed subsystem flow
% \begin{equation}
%     C_{N\subset M}(t)=V^*\Delta_{\omega;M}^{-it}V\Delta_{\omega;N}^{it}.
% \end{equation}
% Modular autonomy of $\mathcal N$, the existence of an $\omega$-preserving expectation, and the identity $C_{N\subset M}(t)=1_{\mathcal H_N}$ for every $t$ are equivalent. When they hold, the recovery map is the predual of the expectation. If $\Omega$ is also cyclic for $\mathcal N$, then $C=D$; for a proper inclusion, common cyclicity excludes an $\omega$-preserving expectation. These statements concern a specified channel and a specified family of states, and a single preserved state does not by itself establish a code.

% Second, \emph{the subsystem flow is a generating function for the failure of correctability} (Proposition~\ref{prop:sufficiency}). For a unit vector $\Psi$ in the stated modular-generator domains, with the four relative entropies finite,
% \begin{align}
%  -i\,\partial_t\langle\Psi|D_{N\subset M}(t)|\Psi\rangle\big|_{t=0}
%  &=\langle\Psi|(K_M-K_N)|\Psi\rangle\nn\\
%  &=\bigl[S_M(\Psi\|\Omega)-S_N(\Psi\|\Omega)\bigr]\nn\\
%  &\quad+\bigl[S_{N'}(\Psi\|\Omega)-S_{M'}(\Psi\|\Omega)\bigr]\ \geq\ 0,
%  \label{eq:intro-entropy}
% \end{align}
% where each entropy uses the states induced by the vectors on the indicated algebra. Under the faithfulness and finiteness hypotheses of Petz's theorem on the equality case of monotonicity, each bracket vanishes exactly when the corresponding restriction is \emph{sufficient} for the pair $(\Psi,\Omega)$, that is, when both states can be recovered from their restrictions by one channel. This pairwise sufficiency differs from recovery of the entire family $\mathcal C_{\mathcal E}$. In particular, $D(t)\Omega=\Omega$, so the expectation of the generator vanishes on $\Omega$ even for a proper inclusion without an $\omega$-preserving expectation. The classical state--reference comparison has a single sufficiency gap.

% Third, \emph{the ergodic hierarchy organizes the autonomous algebras and their codes} (Section~\ref{sec:QEC}). The history algebra generated by the modular orbit of $\mathcal N$ and the memory algebra given by the intersection of all its translates are the smallest $\omega$-expected algebra containing $\mathcal N$ and the largest $\omega$-expected algebra inside it. The memory algebra consists of observables accessible from the retained algebra at every modular time; its elements need not be pointwise invariant. History and memory obey an exact duality: the relative commutant of the history of $\mathcal N$ is the memory of $\mathcal N'\cap\mathcal M$. Thus a generating inclusion in a factor forces its relative commutant to have trivial memory. Common cyclicity of the reference vector makes the inclusion generating and excludes an expectation preserving that reference, but does not by itself imply mixing or a trivial memory algebra for $\mathcal N$. A half-sided modular inclusion adds one-sided nesting. Borchers' positive translation generator $P$ satisfies $K_M-K_N=2\pi P$, and the physical modular evolution transports it as
% \begin{equation}
%     \Delta_M^{-it}P\Delta_M^{it}=e^{2\pi t}P
% \end{equation}
% in the orientation used below. This exact scaling realizes the maximal growth rate associated with the modular chaos bound. Mixing, a trivial tail, and uniqueness of the translation-invariant reference vector require their own hypotheses. This is the sense of maximal modular growth used here.

% Fourth, \emph{nuclearity attaches a spectrum to the inclusion} (Section~\ref{sec:nuclearity}). For inclusions satisfying Buchholz, D'Antoni and Longo's $L^2$-nuclearity condition~\cite{BDL2007}, the subsystem flow at imaginary modular time $i/4$ is trace class. Nested regions with a positive collar supply the geometric setting in which we study this condition. Writing $X=D(i/4)$, we define the canonical positive trace-class operator $\rho_{N\subset M}=X^*X$ and the modular Hamiltonian of the inclusion $H_{N\subset M}=-\log\rho_{N\subset M}$ on $(\ker\rho_{N\subset M})^\perp$. The reference vector is a ground state, and the spectrum on this support is discrete. The operator $\rho_{N\subset M}$ need not have unit trace. Its partition function and spectral form factor satisfy
% \begin{equation}
%     Z_{N\subset M}(\beta)=\operatorname{Tr}\rho_{N\subset M}^{\beta},
%     \qquad
%     F_{N\subset M}(\beta,t)=|Z_{N\subset M}(\beta+it)|^2,
% \end{equation}
% with $Z_{N\subset M}(1/2)=\|X\|_1$ equal to the nuclearity index and both functions finite for $\beta\ge1/2$. Along the line $\operatorname{Im}z=1/4$, $D(z)^*D(z)$ is unitarily conjugate to $\rho_{N\subset M}$ by the intrinsic modular flow, so its spectrum is unchanged. In chiral conformal nets, the singular values of the vacuum interval comparison are related to the spectrum of $L_0$, tying the inclusion spectral form factor to conformal characters. A ramp and plateau in a general theory remain a proposal rather than a consequence of nuclearity. The standard half-sided inclusions with type III$_1$ relative commutant considered below are not split and hence fail $L^2$-nuclearity. This excludes a finite nuclearity index, but failure of $L^2$-nuclearity alone does not determine convergence of the partition function at every other $\beta$. The nuclear and half-sided regimes give complementary descriptions of the subsystem flow, with real modular time governing transport and imaginary modular time probing the spectrum; they are not ordered solely by correctability.

% \paragraph{Relation to previous work.}
% The identification of a state-preserving conditional expectation with an exactly correctable operator-algebraic code, through Takesaki's theorem, is due to Faulkner~\cite{Faulkner2020}, who used it to characterize the exactly correctable subalgebras of a holographic code; conditional expectations onto subalgebras play the same role in~\cite{FaulknerSperanza2024}. Theorem~\ref{thm:correctable} restates this identification with the compressed subsystem flow as a third equivalent condition, and Section~\ref{subsec:QEC-KL} puts it in Knill--Laflamme form. The subsystem flow itself is Borchers' characteristic function of an inclusion~\cite{revolutionizing}; the half-sided structure and its translation group are due to Wiesbrock and Borchers~\cite{Wiesbrock1993,revolutionizing}, the modular chaos bound to de Boer and Lamprou~\cite{deBoerLamprou2020}, and $L^2$-nuclearity to Buchholz, D'Antoni and Longo~\cite{BDL1990,BDL2007}. What we add is the use of the subsystem flow as the quantitative witness of the failure of correctability: the decomposition of its generator into two sufficiency gaps (Proposition~\ref{prop:sufficiency}) and the identification of the two gaps with the two sides of the information--disturbance relation (Section~\ref{subsec:QEC-KL}); the history--memory duality and its decoupling corollary (Lemma~\ref{lem:history-memory}, Corollary~\ref{cor:generating-forgetful}); the reading of the Borchers--Wiesbrock covariance relation as exponential transport of the generator, that is, of the sufficiency gaps (Section~\ref{subsec:geometry-modular-ergodic}); the modular Hamiltonian of an inclusion and its spectral form factor (Section~\ref{sec:nuclearity}); and the classical avatar of the whole construction in Hamiltonian ergodic theory (Section~\ref{sec:classical}), which isolates what the quantum theory adds.
% % TODO (authors): please check that the attribution to \cite{FaulknerSperanza2024} matches that paper's content, and
% % consider citing here the holographic QEC literature (Almheiri--Dong--Harlow; Harlow) and the work on emergent
% % modular and half-sided dynamics in holography (Leutheusser--Liu), none of which is in the current bibliography.

% Throughout, we keep apart several conditions that are often identified. The split property is statistical independence of $\mathcal N$ and $\mathcal M'$; conormality concerns $\mathcal N$ and $\mathcal N'\cap\mathcal M$; finite index measures the relative size of the algebras; mixing is decay of connected correlations of the state on $\mathcal M$. None of these is a definition of another, and we say in each case what is and is not implied for correctability.

% The paper is organized as follows. Section~\ref{sec:classical} develops subsystem flows and classical error correction in Hamiltonian ergodic theory. Section~\ref{section:basics} develops modular dynamics, quantum coarse-grainings as erasure channels, the subsystem flow, its entropy interpretation, and the algebraic obstructions to conditional expectations. Section~\ref{sec:QEC} contains the dictionary with quantum error correction and the modular ergodic hierarchy. Section~\ref{subsec:geometry-modular-ergodic} treats half-sided modular inclusions and maximal modular chaos. Section~\ref{sec:nuclearity} treats nuclearity, the modular Hamiltonian of an inclusion, and the modular spectral form factor. The correlation coefficients of classical subspaces and subsystems, the finite-dimensional coarse-grainings, and the compressed subsystem flow for reference vectors that are not cyclic for the subalgebra are collected in the appendices. The geometric subsystem flows of wedges and balls in quantum field theory, their conformal perturbation theory, and nonperturbative bounds on them appear in the companion paper~\cite{companionQFT}.
\section{Introduction}\label{sec:intro}

An algebraic coarse-graining retains a subalgebra of observables $\mathcal N\subset\mathcal M$ and discards the information needed to distinguish states that agree on $\mathcal N$. An effective dynamics closes on the retained observables when their evolution preserves $\mathcal N$. An effective description is recoverable when restriction to $\mathcal N$ retains enough information to reconstruct a specified family of states on $\mathcal M$. These are the questions of autonomy and quantum error correction. In this paper, we relate them for the dynamics generated by a faithful normal reference state. The central object is the \emph{subsystem flow}, which compares the modular evolution of the full algebra with the intrinsic modular evolution of the retained algebra. Its triviality characterizes a reference-preserving class of exactly correctable codes; its infinitesimal expectation measures losses of distinguishability under restriction.

We take a subsystem to be an inclusion of factors. This includes tensor factors of finite quantum systems and the local algebras of quantum field theory, where a spatial tensor factorization need not be available. The restriction of normal states,
\begin{equation}
 \mathcal T:\mathcal M_*\longrightarrow\mathcal N_*,
 \qquad \mathcal T(\varphi)=\varphi|_{\mathcal N},
\end{equation}
is the erasure channel in either setting. The modular flow of a faithful normal state $\omega$ supplies a dynamics compatible with the algebraic description, while its restriction $\omega|_{\mathcal N}$ supplies an intrinsic dynamics of the subsystem. In a matrix algebra, the modular flow of a Gibbs state is Hamiltonian evolution in thermal units. For general von Neumann algebras, realizing a prescribed dynamics as a modular flow imposes additional restrictions, reviewed in Section~\ref{sec:modular-vs-dynamics}. We work with the modular dynamics of the chosen state throughout.

\paragraph{Modular autonomy as quantum error correction.}
A normal conditional expectation $\mathcal E:\mathcal M\to\mathcal N$ specifies an encoding of subsystem states into global states. Its predual acts as
\begin{equation}
 \mathcal E_*(\varphi_N)=\varphi_N\circ\mathcal E,
 \qquad
 \mathcal T\circ\mathcal E_*=\operatorname{id}_{\mathcal N_*}.
\end{equation}
The normal states satisfying $\varphi\circ\mathcal E=\varphi$ form the code $\mathcal C_{\mathcal E}$. On this family, $\mathcal E_*\circ\mathcal T=\operatorname{id}_{\mathcal C_{\mathcal E}}$, so the same map encodes the logical states and recovers the corresponding code states after erasure. For $\mathcal M=\mathcal N\overline\otimes\mathcal R$ and $\mathcal E=\operatorname{id}_{\mathcal N}\otimes\chi$, the code consists of the product states $\varphi_N\otimes\chi$, with a fixed normal state $\chi$ on the erased factor. Requiring $\omega\circ\mathcal E=\omega$ asks that the reference state belong to this recoverable family.

Takesaki's theorem states that such an expectation exists precisely when $\mathcal N$ is invariant under the modular flow of $\omega$~\cite{Takesaki1972}. Its interpretation in terms of exact recovery has been developed in operator-algebraic quantum error correction~\cite{Faulkner2020}. We express this condition through a comparison of modular evolutions. In the GNS representation of $(\mathcal M,\omega)$, let $\Omega$ be the cyclic and separating vector, let $\mathcal H_N=\overline{\mathcal N\Omega}$, and let $V:\mathcal H_N\hookrightarrow\mathcal H$ be the inclusion isometry. We use the convention $\sigma_t^{\omega;M}=\operatorname{Ad}\Delta_M^{it}$. The modular operator $\Delta_N$ of the restricted state acts on $\mathcal H_N$. The \emph{compressed subsystem flow} is
\begin{equation}
 C_{N\subset M}(t)=V^*\Delta_M^{-it}V\Delta_N^{it}.
\end{equation}
Theorem~\ref{thm:correctable} identifies three equivalent conditions: modular autonomy of $\mathcal N$, an $\omega$-preserving normal conditional expectation onto $\mathcal N$, and $C_{N\subset M}(t)=1$ for all $t$. This is exact correctability for the stated restriction channel and reference-preserving code. Proposition~\ref{prop:KL} gives the corresponding Jones-projection compression criterion. For conormal inclusions, $\mathcal M=\mathcal N\vee(\mathcal N'\cap\mathcal M)$, it also takes the familiar Knill--Laflamme form: the erased relative commutant carries no information about the logical input (Corollary~\ref{cor:KL-conormal}).

If $\Omega$ is cyclic and separating for both algebras, then $V=1$ and the comparison is the unitary cocycle
\begin{equation}
 D_{N\subset M}(t)=\Delta_M^{-it}\Delta_N^{it},
 \qquad K_C=-\log\Delta_C.
 \label{eq:intro-subsystem-flow}
\end{equation}
Its adjoint action first applies the intrinsic modular evolution and then undoes the ambient modular evolution. In particular,
$\operatorname{Ad}D_{N\subset M}(t)(\mathcal N)=\sigma^{\omega;M}_{-t}(\mathcal N)$.
This is Borchers' characteristic function of the inclusion~\cite{revolutionizing}; the name subsystem flow emphasizes the comparison of intrinsic and ambient motion. For a proper inclusion, common cyclicity excludes an expectation preserving $\omega$, so this unitary comparison is nontrivial. The compressed formulation is what also includes the exactly correctable case.

\paragraph{The subsystem flow as a sufficiency witness.}
The comparison has a direct information-theoretic meaning. Let $\Omega$ and $\Psi$ be unit vectors cyclic and separating for both algebras, inducing states $\omega$ and $\varphi$. When the four relative entropies below are finite and $\Psi\in\operatorname{Dom}K_M\cap\operatorname{Dom}K_N$, Proposition~\ref{prop:sufficiency} gives
\begin{align}
 -i\left.\partial_t\langle\Psi|D_{N\subset M}(t)|\Psi\rangle\right|_{t=0}
 &=\langle\Psi|(K_M-K_N)|\Psi\rangle\nonumber\\
 &=\bigl[S_M(\varphi\|\omega)-S_N(\varphi\|\omega)\bigr]\nonumber\\
 &\quad+\bigl[S_{N'}(\varphi\|\omega)-S_{M'}(\varphi\|\omega)\bigr].
 \label{eq:intro-entropy}
\end{align}
The first term is the loss of distinguishability under restriction from $\mathcal M$ to $\mathcal N$; the second is the corresponding loss for the dual inclusion $\mathcal M'\subset\mathcal N'$. Both are nonnegative. By Petz's theorem, each vanishes precisely when the corresponding restriction is sufficient for the pair $(\varphi,\omega)$, meaning that one channel recovers both states~\cite{Petz1986,Petz1988}. The generator therefore measures two obstructions to recovery. This is a comparison at fixed state, rather than an entropy-production law. Pairwise sufficiency also differs from correctability of an entire code: the expression vanishes on the reference vector even when the inclusion is proper.

The analytic continuation of the same operator contains further information. Modular monotonicity gives the quadratic-form order $\Delta_N\ge\Delta_M$, and the subsystem flow extends to a contraction on the strip $0\le\operatorname{Im}z\le1/2$~\cite{revolutionizing}. The real-time derivative in \eqref{eq:intro-entropy}, the imaginary-time spectrum discussed below, and the boundary value $D(i/2)=J_MJ_N$ are different parts of this one comparison operator.

\paragraph{History, memory, and modular geometry.}
To describe how access to a subsystem changes under modular evolution, we associate to it the algebras
\begin{equation}
 \mathcal N_{\rm hist}=\bigvee_{t\in\mathbb R}\sigma_t^\omega(\mathcal N),
 \qquad
 \mathcal N_{\rm mem}=\bigcap_{t\in\mathbb R}\sigma_t^\omega(\mathcal N).
\end{equation}
The history contains all observables generated from the subsystem by the dynamics. The memory consists of observables accessible from the subsystem at every modular time; its elements need not be pointwise invariant. These are respectively the smallest modularly invariant algebra containing $\mathcal N$ and the largest one contained in it. Takesaki's theorem therefore identifies them with the corresponding smallest and largest reference-preserving correctable algebras.

History and memory obey a commutant duality (Lemma~\ref{lem:history-memory}). Writing $\mathcal R=\mathcal N'\cap\mathcal M$, one consequence is
\begin{equation}
 (\mathcal N_{\rm hist})'\cap\mathcal M=\mathcal R_{\rm mem}.
\end{equation}
Thus a subsystem whose modular orbit generates a factor $\mathcal M$ has a relative commutant with no permanent memory. We prove that common cyclicity gives generation. If the inclusion is standard, so that $\Omega$ is also cyclic and separating for $\mathcal R$, the subsystem's own memory is trivial as well (Lemma~\ref{lem:standard-forgetful}). These statements organize persistent access to observables. A rate of decay of connected correlations is a further dynamical question.

Half-sided modular inclusions add an ordering to this orbit: $\sigma_t^\omega(\mathcal N)\subseteq\mathcal N$ for $t\le0$. The Wiesbrock--Borchers theorem supplies a positive translation generator $P$ with
\begin{equation}
 K_M-K_N=2\pi P,
 \qquad
 \Delta_M^{-it}P\Delta_M^{it}=e^{2\pi t}P,
 \qquad P\ge0,
\end{equation}
with the generator identity understood on its natural domain~\cite{Wiesbrock1993,Borchers1992,ArakiZsido2005}. The subsystem-flow generator is therefore transported at the maximal modular growth rate $2\pi$ associated with the modular chaos bound~\cite{deBoerLamprou2020}. Equation~\eqref{eq:intro-entropy} identifies its translation energy with the sum of the two sufficiency gaps. Standard half-sided inclusions of factors also have trivial tail and hence are modular $K$-systems. Standardness further leads to the positive-energy M\"obius structure of a conformal line~\cite{Wiesbrock1993b,GuidoLongoWiesbrock1998}. This gives the known reconstruction of geometry from modular inclusions an interpretation in terms of recoverability.

\paragraph{Imaginary modular time and spectral chaos.}
The trace-class property of $X=D(i/4)$ is the $L^2$-nuclearity condition of Buchholz, D'Antoni and Longo~\cite{BDL2007}. It defines a positive operator and an inclusion Hamiltonian,
\begin{equation}
 \rho_{N\subset M}=X^*X,
 \qquad
 H_{N\subset M}=-\log\rho_{N\subset M}
 \quad\hbox{on }(\ker\rho_{N\subset M})^\perp.
\end{equation}
This Hamiltonian is nonnegative, annihilates $\Omega$, and has discrete spectrum under nuclearity. The operator $\rho_{N\subset M}$ need not have unit trace. Taking traces on the support above, its partition function and spectral form factor are
\begin{equation}
 Z_{N\subset M}(\beta)=\operatorname{Tr}e^{-\beta H_{N\subset M}},
 \qquad
 \mathcal F_{N\subset M}(\beta,t)=|Z_{N\subset M}(\beta+it)|^2,
\end{equation}
with $Z_{N\subset M}(1/2)=\|X\|_1$ equal to the nuclearity index. Both quantities are finite for $\beta\ge1/2$. We prove that ordered enlargement of the outer region at fixed inner algebra and state decreases the finite partition functions. We also establish the exact real structure $J_NH_{N\subset M}J_N=H_{N\subset M}$, with the same conjugation for every outer algebra when $\mathcal N$ and $\Omega$ are fixed.

Concentric vacuum balls in a CFT provide an exact benchmark:
\begin{equation}
 H_{B_r\subset B_R}\simeq2\log(R/r)\,H_{\rm cyl},
 \qquad r<R,
\end{equation}
where $\simeq$ denotes unitary equivalence and $H_{\rm cyl}$ is the dimensionless cylinder Hamiltonian with vacuum energy zero~\cite{companionQFT}. The collar rescales both the inverse temperature and time in the cylinder partition function. The chiral version recovers the conformal-character identity of~\cite{BDL2007}. These identities transfer the full spectral structure, including degeneracies and conformal descendant ladders.

Spectral universality requires an additional hypothesis. The time-averaged form factor is fixed exactly by the energies and their multiplicities, $\sum_E d_E^2e^{-2\beta E}$; its existence is not evidence for chaotic correlations. Conjecture~\ref{conj:modular-SFF} instead concerns the connected, unfolded form factor in high-energy windows for a specified ensemble of deformed vacuum inclusions in an interacting CFT. With the common symmetries resolved, the intrinsic conjugation singles out the Gaussian orthogonal ensemble (GOE) as the candidate class in a real sector. The physical input is chaotic statistics of unprotected primary dimensions; their transfer to the inclusion spectrum is the conjecture. Nuclearity ensures that the spectral objects exist, while universality depends on how changes of shape couple the levels. Proper half-sided inclusions lie in a different spectral regime: $H_{N\subset M}=2P$ has no finite positive-temperature partition function, even though its real-time modular growth is fixed exactly.

\paragraph{The classical comparison and the QFT setting.}
The classical construction motivates the comparison of intrinsic and ambient motion. On the symplectic phase spaces considered in Section~\ref{sec:classical}, a smooth positive density $\rho$ determines a Hamiltonian $K_\rho=-\log\rho$. When the flow is complete, it preserves the state and satisfies a classical KMS identity. Comparing this motion with that of a coarse-grained density $\rho_0$ gives a cocycle whose infinitesimal Hamiltonian is $-\log(\rho/\rho_0)$; its negative expectation is the Kullback--Leibler divergence. Every classical coarse algebra admits an expectation preserving the reference probability measure, independently of whether the state-generated Poisson dynamics is autonomous on that algebra. This separates the classical recovery construction from the additional modular constraint supplied by Takesaki's theorem.

Quantum field theory supplies a natural setting for the distinction between exact recovery, modular growth, and spectral correlations. Reeh--Schlieder cyclicity makes proper common-cyclic inclusions ubiquitous, while split inclusions and half-sided nesting impose different additional structures. Holography brings both aspects into focus: conditional expectations describe exact operator-algebraic reconstruction~\cite{Faulkner2020}, and half-sided modular inclusions have been used to construct emergent bulk translations and times~\cite{LeutheusserLiu2023}. The entropy identity above provides an information-theoretic reading of those translation generators. Our aim is to organize these connections around the subsystem flow, combining the recovery criterion, the history--memory structure, and the spectral theory of the imaginary-time comparison.

The paper is organized as follows. Section~\ref{sec:classical} develops the classical construction. Section~\ref{section:basics} reviews the operator-algebraic background, including the distinction between modular and arbitrary dynamics and the obstructions to conditional expectations. Section~\ref{sec:QEC} establishes the error-correction dictionary. Section~5 develops the sufficiency identity, history and memory, and the modular ergodic hierarchy with its geometric examples. Section~\ref{sec:nuclearity} constructs the inclusion Hamiltonian and formulates the spectral universality conjecture. The appendices collect supplementary classical and operator-algebraic material. Explicit vacuum subsystem flows for wedges and balls, their deformations, and further bounds are developed in the companion paper~\cite{companionQFT}.


\section{Subsystem flows in classical ergodic theory}\label{sec:classical}

%\subsection{Coarse-gainings and Classical net of subregion algebras}
In this Section, we review notions from classical ergodic theory to motivate and build intuition about subsystem flows. 
We begin by reviewing the measure-theoretic definition of a classical dynamical system in Section~\ref{subsec:classical-definitions}, the subspaces and subsystems of observables in Section~\ref{subsec:classical-region-subsystem}, and their role in classical ergodic theory in Section~\ref{subsec:classical-ergodic-hierarchy}. In Section~\ref{subsec:classical-state-flow}, we argue that when the configuration space $X$ is equipped with a Poisson bracket (as suggested by the phase space intuition from physics) {\it the Poisson flow of a state} is a classical avatar of modular dynamics associated to every state. Similar to quantum modular flow, the Poisson flow of a state generates an automorphism of the algebra leaving the state invariant, and satisfies a classical avatar of the KMS condition \footnote{Recall that in the quantum case, these two properties uniquely fix the modular flow of a state~\cite{Takesaki1970}.}. However, crucially for us, as opposed to the quantum modular Hamiltonian, the Poisson Hamiltonian of a state does not satisfy a monotonicity condition under coarse-graining as in the quantum case \footnote{Recall that in the quantum case, this monotonicity originates from the fact that the Tomita operator of the algebra is an extension of its restriction to the subalgebra. In the classical case, there is no analog of the Tomita operator.}. This monotonicity is at the heart of non-commutative results on the reconstruction of geometry from a net of subsystems, such as half-sided modular inclusions and intersections as we review in Section~\ref{subsec:geometry-modular-ergodic}. Readers who are only interested in the quantum discussion can skip to the next section.

For each coarse-graining, we define a notion of a coarse-graining flow. In the case of subspaces, the {\it subspace flow} has the interpretation of a generating function for entropy fluctuations. However, the focus of this work is {\it subsystem flows} which are richer.
Classical many-body systems and field theory come equipped with a net of subsystems, and the quantum generalization of subsystem flows plays a crucial role in the emergence of geometry from a net of abstract subsystems. 

\subsection{Review of classical dynamical systems}\label{subsec:classical-definitions}
% We will see that it is the failure of conditional expectation to intertwine the microscopic phase action with its coarse-grained counterpart that is of interest to us.  The same relative density that measures this failure generates a cocycle whose expectation is exactly the Kullback--Leibler divergence, and this divergence is exactly the entropy increase produced by coarse-graining.  These results provide a rigorous kinematic account of information loss, but they do not by themselves derive an irreversible physical time evolution.

\paragraph{Algebra of observables:} In classical mechanics, one often begins with a phase space that we take to be an orientable smooth manifold \(X\) that admits a strictly positive volume measure. Points \(x\in X\) represent microstates of the system. But from the algebraic viewpoint, the primary object is not the set of points but the commutative von Neumann algebra of observables (complex measurable functions)
\[
\mathcal M=L^\infty(X,\Sigma,\mu),
\]
where \((X,\Sigma,\mu)\) is a standard probability space. When $X$ is a smooth phase space of finite volume, we take $\mu$ to be the normalized volume measure\footnote{$\Sigma$ is the completion of the Borel $\sigma$-algebra with respect to $\mu$. Elements of $L^\infty(X,\Sigma,\mu)$ are equivalence classes of essentially bounded measurable functions that agree $\mu$-almost everywhere.}.


\paragraph{States:} The normalized measure $\mu$ above is interpreted as a distinguished reference state\footnote{A $\sigma$-finite volume measure of infinite total volume instead defines a distinguished weight. We return to this extension in later sections.}. General normal states correspond to probability measures $\nu$ on $X$ that are absolutely continuous with respect to $\mu$ (denoted $\nu\ll \mu$), meaning that for every measurable set $E\subseteq X$, if $\mu(E)=0$ then $\nu(E)=0$.
There exists a $\mu$-measurable function $\rho_\nu:X\to [0,\infty)$ such that $\nu(E)=\int_E \rho_\nu d\mu$ \footnote{In mathematical terminology, $\rho_\nu$ is called a Radon-Nikodym derivative.}. In other words, every normal state can be viewed as a probability density function $\rho_\nu$ on $X$, with expectation values of an observable $f$ in this state captured by integrals on $X$: 
\begin{align}
\braket{f}_\nu=\int_{x\in X}f(x) d\nu(x)=\int_{x\in X} \rho_\nu(x) f(x) d\mu(x)\ .
\end{align}

\paragraph{GNS Hilbert space:} The GNS Hilbert space of the reference state $\mu$ is
\[
  \cH=L^2(X,\Sigma,\mu),
  \qquad
  \mathcal M=L^\infty(X,\Sigma,\mu).
\]
Observables act by multiplication, and the reference state $\mu$ is represented by the constant unit vector $\ket{1}$. A general normal state with density $\rho$ is represented by
\[
  \ket{\sqrt{\rho}}\in\cH,
  \qquad
  \|\ket{\sqrt{\rho}}\|_2^2=\int_X\rho\,\dd\mu=1.
\]
The GNS construction builds a connection between observables and states.


\paragraph{Dynamics:} Dynamics will be described by a measurable one-parameter group of invertible, measure-preserving transformations $\Phi_t:X\to X$. General dynamics need not be smooth. In physics, we are often interested in smooth flows described by a one-parameter family of volume-preserving diffeomorphisms of $X$. Taking a $t$-derivative of how a point $x\in X$ moves under this flow defines a vector field $V$, which preserves the volume measure if and only if $\mathcal{L}_V\mu=0$; i.e., it is divergence-free with respect to $\mu$. 


% \paragraph{Coarse-graining:} Let $\Sigma_0\subseteq\Sigma$ be the sub-$\sigma$-algebra of macroscopically resolved events. We define a coarse-graining as 
% \[
%   Pf:=\E_\mu[f\mid\Sigma_0].
% \]
% On the Hilbert space $\mH$, $P$ is represented by the orthogonal projection onto $L^2(X,\Sigma_0,\mu)$ \footnote{Note that $P$ acts on functions and not measurable sets.}. Coarse-graining by construction preserves the volume measure $\mu$; however, general states are coarse-grained as $\rho_0=P \rho$. A coarse-graining can be viewed as a conditional expectation $E:L^\infty(X,\Sigma,\mu)\to L^\infty(X,\Sigma_0,\mu\mid \Sigma_0)$ that projects to the algebra of coarse-grained observables.
% It is a defining property of this projection that the fine-grained and coarse-grained states are observationally indistinguishable on the macroscopic algebra, even though they may differ on the full $\sigma$-algebra.


% \paragraph{Net of subregion (subspace) coarse-grainings:} In our measure-theoretic setting above, there is an important class (net) of coarse-grainings that are provided geometrically by restricting to subregions $O\subset X$. For a measurable subregion $O \subset X$, we wish to associate a true unital von Neumann subalgebra $\mathcal M_O \subset \mathcal M=L^\infty(X,\Sigma,\mu)$ that captures the observables strictly localized to $O$. 

% To ensure that $\mathcal M_O$ shares the global identity element $\mathbf{1}_X$, this subalgebra is the space of essentially bounded functions that are arbitrarily measurable within $O$ but strictly constant almost everywhere outside of $O$:
% \[
% \mathcal M_O = \{ f \in L^\infty(X, \Sigma, \mu) \mid f = \chi_O f + c \chi_{O^c} \text{ for some } c \in \mathbb{C} \}
% \]
% where $\chi_O$ and $\chi_{O^c}$ are the indicator functions of the region and its complement, respectively. Similarly, there is a subalgebra associated to the complement $\mathcal M_{O^c}$ \footnote{Note that in this construction, the regions $O$ have the same dimension as the phase space $X$.}

% The projections from the global algebra $\mathcal M$ into the local subalgebra $\mathcal M_O$ are normal conditional expectation $E: \mathcal M \to \mathcal M_O$ \footnote{A conditional expectation is a unit-preserving, positive, ultraweakly continuous linear map satisfying the module property $E(f_1 g f_2) = f_1 E(g) f_2$ for all $f_1, f_2 \in \mathcal M_O$ and $g \in \mathcal M$.}. Given a state $\omega_{O^c}$ on $\mathcal M_{O^c}$, these projections (normal conditional expectations) are defined as
% \[
% E_\omega(f)(x) = \chi_O(x) f(x) + \chi_{O^c}(1) \left( \frac{1}{\omega_{O^c}(1)} \int_{O^c} f(y) \rho_\omega(y) \, d\mu(y) \right)\ .
% \]
% In fact, there is a one-to-one correspondence between the states of $\mathcal M_{O^c}$ and the normal-condition expectations for $\mathcal M_O$. It is straightforward to see that this conditional expectation preserves any state for which the probability density $\rho$ restricted to $O^c$ is proportional to $\rho_\omega$. However, if the volume measure of $O^c$ is infinite, there does not exist any volume-preserving conditional expectation.


% \paragraph{Subspace algebras versus subsystem algebras:} Note that, in this construction, we have a direct sum $\mathcal M_O\oplus \mathcal M_{O^c}$ as opposed to a subsystem where we expect a tensor product structure. For a measurable partition \(X=O\sqcup O^c\),
% \[
% L^\infty(X)
% \simeq L^\infty(O)\oplus L^\infty(O^c).
% \]These summands represent mutually exclusive alternatives: the microstate belongs to one region or the other. Their units are \(\chi_O\) and \(\chi_{O^c}\), and their products vanish.
% However, your unital region algebras are
% \[
% \mathcal A(O)=L^\infty(O)\oplus\mathbb C\chi_{O^c},
% \qquad
% \mathcal A(O^c)=\mathbb C\chi_O\oplus L^\infty(O^c).
% \]They are not themselves the two direct summands. They overlap:
% \[
% \mathcal A(O)\cap\mathcal A(O^c)
% =\operatorname{span}\{\chi_O,\chi_{O^c}\}.
% \]The direct summands are the supported algebras \(\chi_OL^\infty(X)\) and \(\chi_{O^c}L^\infty(X)\).
% For a global product \(X=Y\times Z\), equipped with a product measure,
% \[
% L^\infty(X)
% \simeq L^\infty(Y)\,\bar\otimes\,L^\infty(Z).
% \]The base and hidden-variable algebras are
% \[
% \mathcal N=L^\infty(Y)\otimes1,
% \qquad
% \mathcal R=1\otimes L^\infty(Z).
% \]They represent simultaneously existing degrees of freedom: a microstate specifies both \(y\) and \(z\). Their intersection is \(\mathbb C1\), and together they generate the full algebra.
% For the Hamiltonian interpretation, require a product Poisson structure as well:
% \[
% \{f(y),g(z)\}=0.
% \]
% This factorization does not require the state to factorize; \(\rho(y,z)\) may contain correlations.
% The second qualification is that a general Poisson fibration \(X\to Y\) need not admit such a global product decomposition or a globally defined complementary fiber algebra.
% A precise formulation for the paper would be:
% A partition of phase space gives a direct-sum decomposition into mutually exclusive sectors. A product decomposition of phase space gives a tensor-product decomposition into simultaneous subsystems. General Poisson coarse-grainings extend the subsystem picture without necessarily admitting a global tensor-product factorization.


\paragraph{Coarse-grainings:} Classical coarse-graining retains a subalgebra of observables and discards the information needed to distinguish states that agree on this subalgebra. Two useful constructions arise from a subregion of the space of microstates and from a subsystem of observed variables. Both can be formulated in measure theory, before introducing a Hamiltonian or a Poisson bracket. We first describe these constructions and then ask when the retained observables evolve autonomously.

In the measure-theoretic setting above, a coarse description is specified by a sub-$\sigma$-algebra $\Sigma_0\subseteq\Sigma$, with associated unital von Neumann subalgebra
\begin{align}
    \mathcal N=L^\infty(X,\Sigma_0,\mu)\subseteq\mathcal M\ .
\end{align}
The reference measure $\mu$ determines the coarse-graining map
\begin{align}
    E_N(f)=\mathbb E_\mu[f\mid\Sigma_0]\ .
\end{align}
This is the unique $\mu$-preserving normal conditional expectation onto $\mathcal N$. On $L^2(X,\mu)$ it acts as the orthogonal projection onto $L^2(X,\Sigma_0,\mu)$. A general normal state $\omega_\rho(f)=\int_X\rho f\,d\mu$ is coarse-grained to $\omega_0=\omega_\rho\circ E_N$, whose density is $\rho_0=E_N(\rho)$, using the $L^1$ extension of conditional expectation. The two states agree on $\mathcal N$. Thus the subalgebra specifies the retained information, while the chosen expectation specifies how a state on this subalgebra is extended back to the full algebra.


\paragraph{Entropy of a coarse-graining:} Consider a $\mu$-preserving coarse-graining $E$, a normal state with density $\rho$, and its coarse-grained density $\rho_0=E(\rho)$. Then, the relative entropy (KL-divergence) of these two states associates an entropy to this coarse-graining:
\begin{align}
    \label{eq:KL-classical}
    S(\rho\|\rho_0)=\int_X \rho \log(\rho/\rho_0)d\mu=H(\rho_0)-H(\rho)
\end{align}
where $H(\rho)=-\int_X\rho\log\rho\,d\mu$ is the entropy relative to the chosen reference measure, and the equality assumes that the two entropies are finite. The nonnegativity of relative entropy implies that coarse-graining cannot decrease this entropy. This compares a state with its coarse-grained version; it does not by itself imply monotonic entropy growth under time evolution.

\paragraph{Generating function for entropy:} For faithful densities $\rho,\rho_0>0$ almost everywhere, the multiplicative relative phase 
\[
u_t
=
\left(\frac{\rho}{\rho_0}\right)^{it}
\]
is the generating function for coarse-graining entropy because 
\begin{align}
&-i\left.\frac{\dd}{\dd t}u_t\right|_{t=0}
=
\log\frac{\rho}{\rho_0},\nn\\
&S(\rho\|\rho_0)=-i \int_X \rho (\partial_t u_t\big|_{t=0})d\mu\ .
\end{align}
Here the differentiation under the integral is justified when $\rho\log(\rho/\rho_0)$ is integrable. More generally, for every faithful density $\rho$ we can define the strongly continuous unitary multiplication group
\[
  U_t^\rho:=\Mop_{\rho^{it}}
  =\exp\!\bigl(it\Mop_{\log\rho}\bigr),
  \qquad
  (U_t^\rho\psi)(x)=e^{it\log\rho(x)}\psi(x).
\]
Similarly, the coarse-grained density defines
\[
  U_t^{\rho_0}:=\Mop_{\rho_0^{it}}.
\]
The parameter $t$ is an auxiliary phase parameter. Nothing in this construction identifies it with laboratory time. In particular, $U_t^\rho$ does not transport points of $X$; it only rotates the phase of a Hilbert space vector pointwise\footnote{For faithful normal states on the same abelian algebra, $u_t=(\rho/\rho_0)^{it}$ is the Connes Radon-Nikodym cocycle. Relative modular operators remain nontrivial, even though the modular automorphism group of each state is trivial.}.

\subsection{Subspace (single-particle) and subsystem (multi-particle) coarse-grainings}
\label{subsec:classical-region-subsystem}

Subspace and subsystem coarse-grainings organize incomplete information differently. An example of a subspace coarse-graining resolves microstates within a selected region of the phase space and identifies those outside it. A microstate is either inside the coarse-graining region or outside it. The set of microstates adds up in this case; that is why, in anticipation of the generalization to Hilbert space, we refer to them as {\it subspaces} so that addition of the number of independent microstates becomes a direct sum. As we will see, there is a net of subspace algebras on the phase space of a single particle. As an example of such a net of coarse-grainings, consider the reduced state to the part of the phase space of a pendulum where the angle $\theta$ is in some limited range $\theta\in (-\theta_0,\theta_0)$, and the set of all $\theta_0\in (0,\pi)$. 

Next, consider a microstate described using several (potentially independent) degrees of freedom. In this case, the set of microstates multiplies, and in anticipation of generalization to Hilbert spaces, we refer to this case as {\it subsystems}, so that multiplication is the tensor product. As a motivating example, consider two particles with the phase space $(q_1,p_1,q_2,p_2)$, and the restriction to the first particle as a coarse-graining on the second particle's degrees of freedom. 

Note that since subspaces could have the same dimension as the entire phase space, they can be discussed even in a single-particle setup. A subsystem is instead described by a coarse space $Y$ and a map $\pi: X\to Y$ where $Y$ need not be an embedded submanifold of $X$. There is a net of subsystems on the phase space of many-body systems and field theory.


% Evolution transports the resolved region in the first case and the fibers of indistinguishable microstates in the second. Their invariance has correspondingly different meanings: an invariant region defines a conserved event, while an invariant subsystem defines an autonomous factor dynamics. Ergodicity excludes nontrivial conserved events but permits nontrivial autonomous factors.

\paragraph{Subspace coarse-grainings:}
For a measurable region $O\subset X$, define the algebra generated by observables that are constant outside of $O$:
\begin{align}
    \mathcal M_O=\{f\in\mathcal M:\ f=f \cdot 1_O +c 1_{O^c}\text{ for some }c\in\mathbb C\}
\end{align}
where $1_O$ is the indicator function on $O$, and $(f\cdot 1_O)(x)=f(x)1_O(x)$ is point by point multiplication.
This algebra resolves the microstate completely within $O$ and retains only the event ``outside $O$'' on its complement. Assuming $0<\mu(O)<1$, its conditional expectation is
\begin{align}
    E_O(f)=1_O f+
    1_{O^c}\mathbb{E}_{O^c}(f)
\end{align}
where 
\begin{align}
    \mathbb{E}_{O^c}(f)=\frac{1}{\mu(1_{O^c})}\int_{O^c} f d\mu
\end{align}
is the average of $f$ on $O^c$.
For $O_1\subseteq O_2$, we have $\mathcal M_{O_1}\subseteq\mathcal M_{O_2}$, so this construction gives a net of coarse observable algebras. Since a microstate $x\in X$ is either in $O$ or $O^c$ we have 
\begin{align}
    \mathcal M_X=1_O\mathcal M\oplus 1_{O^c}\mathcal M
    \simeq L^\infty(O,\mu|_O)\oplus L^\infty(O^c,\mu|_{O^c})\ .
\end{align}
Our unital subspace algebras have the form $\mathcal M_O=1_O\mathcal M\oplus\mathbb C1_{O^c}$, so $\mathcal M_O$ and $\mathcal M_{O^c}$ are not themselves the two direct summands. In fact,
\begin{align}
    \mathcal M_O\cap\mathcal M_{O^c}=\operatorname{span}\{1_O,1_{O^c}\}\ .
\end{align}

\paragraph{Subsystem coarse-grainings:}
Suppose instead that the probability space has a product decomposition
\begin{align}
    (X,\Sigma,\mu)=(Y\times Z,\Sigma_Y\otimes\Sigma_Z,\mu_Y\otimes\mu_Z)\ .
\end{align}
A microstate now specifies both an observed variable $y$ and a hidden variable $z$. The observable algebra factorizes as
\begin{align}
    \mathcal M\simeq\mathcal M_Y\,\bar\otimes\,\mathcal M_Z,
    \qquad \mathcal M_Y=L^\infty(Y,\mu_Y),\quad
    \mathcal M_Z=L^\infty(Z,\mu_Z)\ .
\end{align}
Retaining the subsystem $Y$ means retaining $\mathcal N_Y=\mathcal M_Y\otimes 1$. The corresponding conditional expectation averages over the hidden variables:
\begin{align}
    (\mathcal{E}_Y f)(y,z)=\int_Z f(y,z')\,d\mu_Z(z')\ .
\end{align}
This operation forgets the hidden variables at every microstate. The tensor factors describe simultaneously specified degrees of freedom. Although the reference measure is a product, a general state with density $\rho(y,z)$ can have correlations between them. 

More generally, a measurable map $\pi:X\to Y$ between standard probability spaces specifies coarse observables through
\begin{align}
    &\mathcal N_Y=\mathcal M_Y\circ \pi\subset \mathcal M,
    \qquad \mu_Y=\mu\circ \pi^{-1}\nn\\
    &\mathcal M_Y=L^\infty(Y,\mu_Y)
\end{align}

Here $\mathcal N_\pi=\pi^*\mathcal M_Y$ is the pullback of $\mathcal M_Y$, and $\mu_Y=\pi_*\mu$ is the
push-forward measure. Under the standard Borel assumptions, disintegration gives
conditional probability measures $\mu_y$ concentrated on the fibers $\pi^{-1}(y)$.

The conditional expectation is
\begin{align}
    (\mathcal{E}_\pi f)(x)=\int_{\pi^{-1}(\pi(x))}f(x')\,d\mu_{\pi(x)}(x')\ .
\end{align}
The fiber measures may depend on $y$, and no global product decomposition is required. 

Note that both subspace and subsystem conditional expectations are special cases of a common measure-theoretic notion of coarse-graining.

\paragraph{Dynamically closed and invariant coarse-grainings:}
Let $\Phi_t:X\to X$ be a measurable one-parameter group of invertible, measure-preserving transformations. It evolves observables according to
\begin{align}
    \alpha_t(f)=f\circ\Phi_{-t}\ .
\end{align}
For any coarse algebra $\mathcal N\subseteq\mathcal M$, we can follow the family $\mathcal N_t=\alpha_t(\mathcal N)$. We call the coarse description dynamically closed when
\begin{align}\label{eq:classical-coarse-invariance}
    \alpha_t(\mathcal N)=\mathcal N,\qquad t\in\mathbb R\ .
\end{align}
For nonatomic $\mu$ and $0<\mu(O)<1$, invariance of the subspace algebra $\mathcal M_O$ violates ergodicity, as we show below.
For the subsystem $Y$ in a product space $Y\times Z$, dynamical closure means that the $Y$-component of $\Phi_t(y,z)$ depends only on $y$ \footnote{Note that this does not follow from preservation of the product measure.}. For a coarse-graining specified by $\pi$, this means that there is a measure-preserving evolution $\Psi_t$ on the coarse space such that
\begin{align}\label{eq:classical-factor-intertwining}
    \pi\circ\Phi_t=\Psi_t\circ\pi\ ,
\end{align}
understood modulo null sets. Evolving the microstate and then retaining its coarse variables gives the same result as evolving the coarse variables themselves. In ergodic theory such a coarse description is called a {\it factor system}, and we say the coarse-grained variables evolve autonomously. Intuitively, this means that the future of the macroscopic (or observed) variables can be determined from the present macroscopic state, without the need to know the hidden microscopic details. In other words, there is an effective theory that describes the coarse-grained variables.


% For a subregion, evolution transports its associated algebra:
% \begin{align}
%     \alpha_t(\mathcal M_O)=\mathcal M_{\Phi_t(O)}\ .
% \end{align}
% An invariant region therefore defines an invariant coarse algebra. For nonatomic $\mu$ and $0<\mu(O)<1$, the converse holds as well: the single atom representing the unresolved complement identifies $O^c$ from $\mathcal M_O$. Ergodicity excludes nontrivial invariant regions, but it does not exclude nontrivial factor systems. Invariance of a subalgebra permits its observables to evolve into one another; they need not be individually fixed.

%  For example, on the two-torus with uniform product measure, the flow
% \begin{align}
%     \Phi_t(y,z)=\bigl(y+t\sin(2\pi z)\!\!\pmod 1,\ z\bigr)
% \end{align}
% preserves the measure, but the future observed coordinate depends on the hidden coordinate. Hence $\mathcal N_Y$ is not invariant. An autonomous observed subsystem is also weaker than a product dynamics: the hidden variables may evolve in a way that depends on the observed variables.

\paragraph{Effective theory (Compatibility of coarse-graining and evolution):}
For the measure-preserving dynamics above, invariance of $\mathcal N$ is equivalent to
\begin{align}\label{eq:classical-expectation-equivariance}
    \mE_{\pi}\circ\alpha_t=\alpha_t\circ \mE_{\pi},
    \qquad t\in\mathbb R\ .
\end{align}
Indeed, when $\mathcal N_Y$ is invariant, $\alpha_t\circ \mE_{\pi}\circ\alpha_{-t}$ is another $\mu$-preserving conditional expectation onto $\mathcal N_Y$, so uniqueness gives the equality. Conversely, this equality ensures that $\alpha_t$ maps the range of $\mE_{\pi}$ onto itself. The physical interpretation is macroscopic predictability: the future observed variables can be determined from their present values. This compatibility supplies an effective dynamics on the coarse algebra.


\paragraph{Correlation functions of observables and subalgebras:}
For complex functions $f,g\in\mathcal M$, define the correlation function and its connected part as
\begin{align}
    &C_{f,g}(t)=\mu\bigl(f^*\alpha_t(g)\bigr)\nn\\
    &C^{con}_{f,g}(t)
    :=\mu\bigl(f^*\alpha_t(g)\bigr)-\mu(f^*)\mu(g),
    \qquad \mu(f)=\int_X f\,d\mu\ .
\end{align}
Ergodicity is equivalent to the vanishing of the time-averaged connected correlation for every pair of observables:
\begin{align}
    \lim_{T\to\infty}\frac{1}{T}\int_0^T C^{con}_{f,g}(t)\,dt=0\ .
\end{align}
A stronger condition, {\it weak mixing}, requires
\begin{align}
    \lim_{T\to\infty}\frac{1}{T}\int_0^T |C^{con}_{f,g}(t)|^2\,dt=0\ .
\end{align}
The flow is {\it mixing} if $C^{con}_{f,g}(t)\to0$ as $t\to+\infty$. These conditions distinguish cancellation after time averaging from actual loss of correlation at late times. Mixing does not specify a decay rate and does not imply that there are no autonomous subsystems.

Quantitative measures of the correlation between two subspace algebras, or between two subsystems, and their relation to the total variation distance from the product measure, are collected in Appendix~\ref{app:classical-correlations}.

\paragraph{Subsystem coarse-grainings of a free relativistic field theory:} Classical field theory provides a natural setting for a net of subsystem algebras with infinitely many degrees of freedom. For an unconstrained real scalar field on a Cauchy surface $\Sigma$, a phase-space point is a pair of initial data $(\phi,\pi)$. An open spatial region $U\subset\Sigma$ selects the unital polynomial algebra
\begin{align}
    \mathcal N_U
    :=\operatorname{Alg}_{\mathbb R}\{1,\phi(f),\pi(g):f,g\in C_c^\infty(U)\},
    \qquad
    \{\phi(f),\pi(g)\}=\int_\Sigma fg\,d\Sigma\ .
\end{align}
Here $\phi(f)=\int_\Sigma f\phi\,d\Sigma$, and similarly for $\pi(g)$; the other brackets between generators vanish. Thus $\mathcal N_U$ is a Poisson subalgebra, and algebras associated with disjoint spatial regions Poisson commute. Retaining $\mathcal N_U$ resolves the field data in $U$ while discarding the exterior data. Although labeled by a spatial region, this is a subsystem coarse-graining: $U$ selects degrees of freedom, whereas a region of phase space selects entire field configurations.


Choose an ambient Poisson algebra $\mathcal M$ containing these observables and stable under a complete Hamiltonian evolution. With our convention for the pullback, the transported subsystem is
\begin{align}
    \alpha_t^H(F)=F\circ\Phi_{-t}^H,
    \qquad \partial_t\alpha_t^H(F)=\{H,\alpha_t^H(F)\},
    \qquad \mathcal N_{U,t}=\alpha_t^H(\mathcal N_U)\ .
\end{align}
Hamiltonian evolution preserves Poisson brackets, so every $\mathcal N_{U,t}$ remains a Poisson subalgebra. Invariance for all $t$ gives an autonomous reversible subsystem. More generally, propagation can make the transported observables depend on initially discarded data. Whether their time translates generate $\mathcal M$, or instead form a nested family, is a property of the chosen subsystem and dynamics together.


\subsection{Ergodic hierarchy, subspaces and subsystems}
\label{subsec:classical-ergodic-hierarchy}

The ergodic hierarchy describes how much information about a state remains accessible at late times. In the algebraic language, it can be studied both through correlations of observables and through the evolution of coarse observable algebras. For our purposes, the relevant distinction is between conserved events, autonomous subsystems, and subalgebras that are preserved only in one time direction.

\paragraph{Ergodicity:}
Let $(\mathcal M,\mu,\alpha_t)$ be the measure-preserving dynamical system introduced above. The flow is {\it ergodic} if every measurable set invariant under all $\Phi_t$ has measure zero or one. Equivalently, its fixed-point algebra is trivial:
\begin{align}
    \mathcal M^\alpha
    :=\{f\in\mathcal M:\alpha_t(f)=f\text{ for all }t\in\mathbb R\}
    =\mathbb C1\ .
\end{align}
For nonatomic $\mu$ and $0<\mu(O)<1$, ergodicity forbids invariant subspace algebras. If 
\begin{align}
    \alpha_t(\mathcal M_O)=\mathcal M_O,\qquad t\in\mathbb R,
\end{align}
then $\Phi_t(O)=O$ modulo null sets and $1_O\in\mathcal M^\alpha$ (it is a {\it conserved event}), which contradicts ergodicity. However, ergodicity does not exclude an invariant subsystem $\mathcal N\subset\mathcal M$: the condition $\alpha_t(\mathcal N)=\mathcal N$ permits observables in $\mathcal N$ to evolve into one another. Such a subsystem is an autonomous factor of the dynamics (conserved subsystem). Taking two identical independent copies of a classical dynamical system always results in autonomous factors; however, a standard result in ergodic theory implies that if the original system is weakly mixing, the combined dynamical system is ergodic \footnote{For an invertible discrete-time example, take two independent bilateral Bernoulli shifts on $\{0,1\}^{\mathbb Z}$ with their product probability measures. Each shift, and their product, is mixing.}.

\paragraph{Generating property:}
Given a coarse algebra $\mathcal N$, the family $\mathcal N_t=\alpha_t(\mathcal N)$ describes the observables accessible as this coarse description is transported by the dynamics. If $\mathcal N_t=\mathcal N$, evolution remains within the original description. Otherwise, the transported observables can depend on variables that were initially hidden. The algebra
\begin{align}
    \mathcal M_N:=\bigvee_{t\in\mathbb R}\alpha_t(\mathcal N)
\end{align}
contains the information accessible from all time translates of the coarse observables. Here $\bigvee$ denotes the generated von Neumann algebra. The condition $\mathcal M_N=\mathcal M$ says that these translates generate the full observable algebra. This is called the {\it generating property} and is distinct from ergodicity of the dynamics.

%{\paragraph Correlations among subalgebras:} We can also define 


\paragraph{Kolmogorov property:}
Of particular interest are subalgebras satisfying
\begin{align}\label{eq:classical-one-sided-inclusion}
    \alpha_t(\mathcal N_+)\subseteq\mathcal N_+,\qquad t\geq0\ .
\end{align}
The family $\mathcal N_t=\alpha_t(\mathcal N_+)$ is then decreasing: $\mathcal N_{t+s}\subseteq\mathcal N_t$ for $s\geq0$. To interpret this inclusion, one can take $\mathcal N_+$ to be generated by observations on a half-infinite time interval, with the orientation chosen so that positive $t$ removes the earliest part of the record.

The Kolmogorov, or $K$, property can be formulated through the existence of such a one-sided algebra whose translates generate the full algebra and whose tail is trivial:
\begin{align}\label{eq:classical-K-algebra}
    \bigvee_{t\in\mathbb R}\mathcal N_t=\mathcal M,
    \qquad
    \bigcap_{t\geq0}\mathcal N_t=\mathbb C1\ .
\end{align}
The first condition says that the complete record resolves all observables. The second says that no nonconstant observable remains accessible arbitrarily far into the retained half of the record. The analogous conditions with integer time describe a $K$-automorphism. They provide an inclusion-based formulation of loss of memory and motivate the future and past algebras used in quantum ergodic theory~\cite{ouseph2024local}.

\paragraph{Example: a Bernoulli shift:}
For independent uniform bits $x_n$, $n\in\mathbb Z$, with the shift $\alpha(x_n)=x_{n+1}$ and $\mathcal N_+$ the algebra generated by $\{x_n:n\geq0\}$, one has $\alpha^k(\mathcal N_+)=\operatorname{vN}\{x_n:n\geq k\}\subsetneq\mathcal N_+$ for $k>0$, the translates generate $\mathcal M$, and the tail is trivial by the Kolmogorov zero--one law: the retained description forgets one coordinate per step although the shift is invertible and measure preserving. This is a subsystem coarse-graining, the retained and discarded coordinates being tensor factors, and the strict one-sided nesting is possible only for subsystems. In a finite measure space $\Phi_t(O)\subseteq O$ and measure preservation force $\mu(O\setminus\Phi_t(O))=0$, so a region, and with it a subspace algebra $\mathcal M_O$ in the nonatomic setting of Section~\ref{subsec:classical-region-subsystem}, cannot shrink strictly under a measure-preserving flow.

\paragraph{Coarse-graining along a decreasing family:}
Each $\mathcal N_t$ has a $\mu$-preserving conditional expectation $\mathcal{E}_t:\mathcal M\to\mathcal N_t$. For $s\geq t$, the tower property gives
\begin{align}
    \mathcal{E}_s\mathcal{E}_t=\mathcal{E}_s\ .
\end{align}
For a fixed state with density $\rho$, its coarse-grained densities $\rho_t=\mathcal{E}_t(\rho)$ therefore satisfy
\begin{align}
    H(\rho_s)\geq H(\rho_t),
    \qquad
    S(\rho\|\rho_s)\geq S(\rho\|\rho_t),
    \qquad s\geq t,
\end{align}
whenever the entropies are finite. Here the loss of information follows from the nesting of the retained algebras. It does not require irreversible microscopic dynamics. For a trivial tail, the reverse martingale theorem further gives $E_t(f)\to\mu(f)1$ in $L^2$ for every $f\in L^2(X,\mu)$.

% \paragraph{Towards state-dependent subalgebra flows:}
% The structures above depend on a specified dynamics $\alpha_t$. Measure theory alone does not assign a nontrivial flow to a state, and all the coarse algebras considered here admit $\mu$-preserving conditional expectations. In the quantum setting, modular theory supplies a distinguished dynamics, and compatibility with that dynamics controls the existence of state-preserving expectations. The one-sided relation \eqref{eq:classical-one-sided-inclusion} anticipates the inclusion structure of a half-sided modular flow, without by itself supplying its quantum modular properties.

% To make a classical comparison with state-generated modular evolution, we next introduce a Poisson bracket and the Hamiltonian associated with $-\log\rho$. The ergodic properties of this dynamics must be assessed on its appropriate invariant sectors: on a symplectic phase space it conserves $\rho$ pointwise, so a nonconstant smooth density already supplies nontrivial conserved events on the full phase space.

\subsection{Classical avatar of modular flow in Hamiltonian systems}
\label{subsec:classical-state-flow}

So far, the state and the dynamics have been independent inputs. A probability density assigns weights to microstates, but measure theory alone does not turn these weights into a motion of the microstates. A Poisson bracket supplies the missing structure: it converts a function into a Hamiltonian vector field. Applied to the logarithm of a density, it gives a distinguished state-dependent dynamics. This construction will let us compare the dynamics associated with a state and with its coarse description.

\paragraph{Poisson brackets and Hamiltonian motion:}
A Poisson structure on a smooth manifold $X$ is an antisymmetric bivector $\Pi$ whose bracket
\begin{align}
    \{f,g\}=\Pi(df,dg),\qquad f,g\in C^\infty(X),
\end{align}
satisfies the Jacobi identity. The bracket obeys the Leibniz rule in each argument, so fixing a smooth real Hamiltonian $H$ defines a vector field. We use the convention
\begin{align}\label{eq:classical-Hamiltonian-convention}
    X_H(f)&=\{f,H\},\qquad
    \alpha_t^H(f)=f\circ\Phi_{-t}^H,\nn\\
    \delta_H(f)&:=\left.\partial_t\alpha_t^H(f)\right|_{t=0}=\{H,f\},
\end{align}
where $\Phi_t^H$ is the flow of $X_H$. Thus the generator on observables has the opposite sign from the vector field transporting points. The Jacobi identity ensures that Hamiltonian evolution preserves the Poisson bracket. In canonical coordinates, $\{q_i,p_j\}=\delta_{ij}$ and the point trajectories satisfy the usual Hamilton equations.

\paragraph{Casimirs and symplectic leaves:}
A Casimir is a function $c$ with $\{c,f\}=0$ for every smooth $f$. It is conserved by every Hamiltonian, rather than by one chosen dynamics. The points accessible from a given microstate by successive Hamiltonian motions form a symplectic leaf. Every Casimir is constant on each leaf, although global Casimirs need not distinguish all leaves. A leaf is therefore a classical analogue of an irreducible sector for Hamiltonian motion. This analogy concerns the Poisson bracket: the pointwise algebra $L^\infty(X,\mu)$ remains abelian, and restricting it to a leaf does not turn it into a nonabelian factor.

\paragraph{Liouville measure and the state Hamiltonian:}
For the explicit constructions below, take $X$ to be a connected compact symplectic manifold without boundary, with symplectic form $\omega$ and normalized Liouville measure
\begin{align}
    d\mu=\frac{\omega^n/n!}{\int_X\omega^n/n!},\qquad \dim X=2n.
\end{align}
Hamiltonian flows preserve this measure. Equivalently,
\begin{align}\label{statepreserve}
    \int_X\{f,g\}\,d\mu=0
\end{align}
for smooth $f,g$. Compactness makes smooth Hamiltonian flows complete and avoids boundary terms. On a noncompact phase space these properties, together with the required integrability, must be checked separately.

Let $\rho$ be a smooth strictly positive probability density relative to $\mu$. Define its {\it state Hamiltonian} and the associated evolution by
\begin{align}\label{eq:classical-state-Hamiltonian}
    K_\rho=-\log\rho,\qquad
    \Phi_t^\rho:=\Phi_t^{K_\rho},\qquad
    \alpha_t^\rho:=\alpha_t^{K_\rho}.
\end{align}
The infinitesimal evolution of an observable is
\begin{align}\label{eq:classical-state-derivation}
    \delta_\rho(f)=\{K_\rho,f\}=\{f,\log\rho\}.
\end{align}
The flow moves along level sets of the density, since $X_{K_\rho}(\rho)=0$. It preserves both $\mu$ and the state $\omega_\rho(f)=\int_X\rho f\,d\mu$. For a Gibbs density $\rho=Z^{-1}e^{-\beta H}$, we have $K_\rho=\beta H+\log Z$ and $\Phi_t^\rho=\Phi_{\beta t}^H$. Thus the state recovers the Hamiltonian motion, with a rescaled parameter; for a general density there is no prior identification of this parameter with physical time.

\paragraph{The classical KMS identity:}
The logarithm is distinguished by an integration-by-parts identity. Using \eqref{statepreserve},
\begin{align}
    0=\int_X\{g,f\rho\}\,d\mu
    =-\omega_\rho(\{f,g\})+
    \omega_\rho\bigl(f\{g,\log\rho\}\bigr).
\end{align}
Consequently,
\begin{align}\label{KMSclassical}
    \omega_\rho(\{f,g\})=\omega_\rho(f\delta_\rho(g)).
\end{align}
For a Gibbs state this reads $\omega_\beta(\{f,g\})=\beta\omega_\beta(f\delta_H(g))$. This is the classical, or static, KMS condition, written with the time orientation in \eqref{eq:classical-Hamiltonian-convention}~\cite{vanDeVen2022}. It is stronger than stationarity: for a given smooth positive density, testing against all smooth $f$ determines $\delta_\rho(g)$ for every $g$. Conversely, on a connected symplectic manifold, imposing this identity for a specified Hamiltonian dynamics determines a smooth positive density to be Gibbs, up to normalization. The identity also motivates formulations on algebras of local observables in infinite systems, where existence of the dynamics and equilibrium states requires additional hypotheses; a global Gibbs density need not exist.

\paragraph{Poisson modular flow:}
There is a geometric formulation of the same construction. For a smooth positive measure $\nu$, define its Poisson modular vector field by
\begin{align}
    \mathfrak m_\nu(f):=\operatorname{div}_\nu X_f.
\end{align}
This measures the failure of Hamiltonian motions to preserve $\nu$~\cite{Weinstein1997}. For $d\nu=\rho\,d\mu$ and Liouville $\mu$, the divergence formula gives
\begin{align}
    \mathfrak m_\nu(f)=X_f(\log\rho)
    =\{f,-\log\rho\}=X_{K_\rho}(f).
\end{align}
Thus $\Phi_t^\rho$ is precisely this Poisson modular flow; its action on observables uses the inverse pullback, so $\delta_\rho=-\mathfrak m_\nu$. On a general Poisson manifold and for an arbitrary smooth reference measure, the formula instead reads $\mathfrak m_{\rho\mu}=\mathfrak m_\mu+X_{-\log\rho}$. A reference measure invariant under every Hamiltonian need not exist. Our symplectic setting removes this additional obstruction.

The construction uses the Poisson structure in an essential way. It is distinct from the multiplication group $U_t^\rho$ introduced above and from the trivial modular automorphism group of the abelian von Neumann algebra. It also supplies no automatic ergodic or entropy-production property. A nonconstant $\rho$ is itself a conserved observable, so the flow cannot be ergodic on the full probability space. Questions about ergodicity and mixing must be posed on invariant sectors with suitable invariant measures. The comparison of full and coarse state Hamiltonians, however, is meaningful without assuming ergodicity.

\subsection{Subsystem Poisson flows in classical many-body systems and field theory}
\label{subsec:classical-subalgebra-Hamiltonian}

A coarse description determines a reduced state and hence, when the coarse space is Poisson, its own state-generated dynamics. We can then ask whether the reduced dynamics agrees with what the full dynamics does to the observed variables. Comparing the two motions gives a classical model of the subsystem flows studied in this paper. The comparison separates evolution already visible within the retained algebra from evolution involving the discarded variables.

\paragraph{Poisson subsystems:} We say a subalgebra $\mathcal N_Y\subset \mathcal M$ is a Poisson subsystem if for all $b_1,b_2\in \mathcal N_Y$, we have $\{b_1,b_2\}\in \mathcal N_Y$. In the geometric setting below, $\pi:X\to Y$ is a smooth surjective Poisson submersion:
\begin{align}\label{eq:classical-Poisson-subsystem}
    \{\pi^*f,\pi^*g\}_X=\pi^*\{f,g\}_Y,
    \qquad f,g\in C^\infty(Y).
\end{align}
The smooth coarse observables $\mathcal N=\pi^*C^\infty(Y)$ are closed under the bracket. This condition allows us to speak of Hamiltonian dynamics within the coarse description.

The simplest setting makes both the geometric and probabilistic structures explicit. Take a product $X=Y\times Z$ of connected compact symplectic manifolds, with product symplectic structure and normalized Liouville measure $\mu=\mu_Y\otimes\mu_Z$. The projection to $Y$ is Poisson, and the associated von Neumann subalgebra is $\mathcal N_Y=L^\infty(Y,\mu_Y)\otimes1$. For a smooth positive density $\rho$, restriction of the state gives the marginal density
\begin{align}\label{eq:classical-marginal-Hamiltonian}
    \rho_Y(y)=\int_Z\rho(y,z)\,d\mu_Z(z),\qquad
    K_Y=-\log\rho_Y.
\end{align}
The expectation $\mE_Y$ from Section~\ref{subsec:classical-region-subsystem} extends this marginal back to $X$ as $\rho_0(y,z)=\rho_Y(y)$. Here, $\rho_Y=\pi_*\rho$ is the classical analog of the reduced state on $Y$, and $\rho_0=\pi^*\rho_Y=\pi^*\pi_* \rho$ is the pullback of $\rho_Y$ to $X$, which is the coarse-grained version of $\rho$. Its Hamiltonian is $K_0=\pi^*K_Y$, whose flow is
\begin{align}
    \Phi_t^0(y,z)=(\Phi_t^Y(y),z).
\end{align}
This provides a definite lift of the reduced motion to the full phase space. For a general Poisson submersion, a specified coarse Hamiltonian can likewise be pulled back, but the marginal measure need not be Liouville and a complementary subsystem need not exist. We use the product setting below to keep the marginal, its reference measure, and its lifted Hamiltonian unambiguous.


% \NL{Include discussion of singular inclusions, if they exist}

\paragraph{Agreement between reduced and full Poisson dynamics:}
Agreement means that the observed trajectory is determined by its initial coarse state and follows the flow generated by $K_Y$:
\begin{align}\label{eq:classical-state-intertwining}
    \pi\circ\Phi_t^\rho=\Phi_t^Y\circ\pi.
\end{align}
For our complete smooth flows, this is equivalent to the infinitesimal condition
\begin{align}\label{eq:classical-state-defect}
    \{\pi^*g,K_\rho\}_X-\pi^*\{g,K_Y\}_Y
    =\{\pi^*g,K_\rho-K_0\}_X=0,
    \qquad g\in C^\infty(Y).
\end{align}
The left side measures the part of the observed velocity missed by the reduced state Hamiltonian. A Hamiltonian invisible to all coarse observables belongs to the Poisson centralizer
\begin{align}
    \mathcal N^{\mathrm c}_{\mathrm P}
    :=\{h\in C^\infty(X):\{h,b\}=0\text{ for every }b\in\mathcal N\}.
\end{align}
Thus agreement is precisely the condition $K_\rho-K_0\in\mathcal N^{\mathrm c}_{\mathrm P}$. This is an analogue of a relative commutant defined using Poisson brackets. The ordinary operator commutant of a classical subalgebra contains the whole abelian observable algebra and cannot express this condition.

In the product setting, nondegeneracy of the symplectic form and connectedness of $Y$ imply that this centralizer consists of functions of $z$ alone. Hence \eqref{eq:classical-state-intertwining} is equivalent to
\begin{align}\label{eq:classical-product-state-closure}
    K_\rho(y,z)=K_Y(y)+K_Z(z),
    \qquad
    \rho(y,z)=\rho_Y(y)\rho_Z(z).
\end{align}
Normalization of the marginal fixes $\int_Z\rho_Z\,d\mu_Z=1$. The condition \eqref{eq:classical-state-defect} implies that the coarse observables have the same Poisson evolution under $\rho_0$ and $\rho$. In other words:
\begin{align}
    \forall g\in C^\infty(Y):\qquad \alpha^\rho_t\circ\alpha_{-t}^0(\pi^*g)=\pi^*g\ .
\end{align}
This compares full and reduced observable evolution, as in Section~\ref{subsec:quantum-coarse}; the state-preserving expectation criterion itself differs in the classical and quantum theories.

\paragraph{Comparing the two flows:}
Define the relative transformation by first undoing the lifted coarse motion and then applying the full motion:
\begin{align}\label{eq:classical-relative-transformation}
    R_t=\Phi_t^\rho\circ\Phi_{-t}^0,
    \qquad
    \mathcal V_t(f):=f\circ R_t^{-1}
    =\alpha_t^\rho\circ\alpha_{-t}^0(f).
\end{align}
For each $t$, this is a symplectic transformation and a measure-preserving automorphism, respectively. Since the lifted coarse flow preserves $\mathcal N_Y$,
\begin{align}\label{eq:classical-relative-algebra-flow}
    \mathcal V_t(\mathcal N_Y)=\alpha_t^\rho(\mathcal N_Y).
\end{align}
The relative transformation therefore follows the same moving subalgebra as the full dynamics, while canceling the chosen internal coarse evolution. It fixes every coarse observable exactly when \eqref{eq:classical-state-intertwining} holds. For a product state, it can still move hidden variables (the relative commutant): $R_t(y,z)=(y,\Phi_t^Z(z))$.

The infinitesimal generator of $R_t$ is the difference of the state Hamiltonians:
\begin{align}\label{eq:classical-relative-Hamiltonian}
    K_{\mathrm{rel}}&=K_\rho-K_0=-\log\frac{\rho}{\rho_0},\nn\\
    \left.\partial_tR_t\right|_{t=0}&=X_{K_{\mathrm{rel}}},\qquad
    \left.\partial_t\mathcal V_t(f)\right|_{t=0}=\{K_{\mathrm{rel}},f\}.
\end{align}

Every normal state of the hidden factor, with density $\rho_Z$ relative to $\mu_Z$,
gives an expectation $E_{\rho_Z}:\mathcal M\to\mathcal N_Y$:
\begin{equation}
 (E_{\rho_Z}f)(y)=\int_Z\rho_Z(z)f(y,z)\,d\mu_Z(z).
\end{equation}
These states parameterize the expectations with $y$-independent averaging.
General normal expectations allow measurable probability kernels $\nu_y(dz)$
depending on $y$, with $\mu_Y(dy)\nu_y(dz)\ll\mu_Y\otimes\mu_Z$:
$E_\nu f(y)=\int_Z f(y,z)\nu_y(dz)$.

\paragraph{Example: a uniform marginal with nonautonomous observables:}
Take $Y$ and $Z$ to be two-dimensional tori with angular coordinates $(q,p)$ and $(Q,P)$, each of period $2\pi$, and the product symplectic form $dq\wedge dp+dQ\wedge dP$. Consider
\begin{align}
    \rho(q,p,Q,P)=\frac{e^{\epsilon\cos(q-Q)}}{\mathcal Z_\epsilon},
    \qquad
    \mathcal Z_\epsilon=\int_X e^{\epsilon\cos(q-Q)}\,d\mu,
    \qquad \epsilon\ne0.
\end{align}
Translation invariance of the angular integral gives $\rho_Y=1$, so $K_0=0$ and the reduced state predicts no motion. The full state Hamiltonian is $K_\rho=-\epsilon\cos(q-Q)+\log\mathcal Z_\epsilon$, giving
\begin{align}
    \dot q=\dot Q=0,\qquad
    \dot p=-\epsilon\sin(q-Q),\qquad
    \dot P=\epsilon\sin(q-Q).
\end{align}
For example, the bounded coarse observable $e^{ip}$ evolves as
\begin{align}
    \alpha_t^\rho(e^{ip})
    =\exp\bigl(i[p+\epsilon t\sin(q-Q)]\bigr).
\end{align}
Its future value depends on the hidden angle $Q$, so it leaves $\mathcal N_Y$ for $t\ne0$. The full state and its uniform marginal are stationary, but the retained observables do not evolve autonomously. Here $R_t=\Phi_t^\rho$ directly displays the motion that the marginal state Hamiltonian misses. The example also exhibits the failure of monotonicity announced at the beginning of this section: the relative Hamiltonian $K_{\mathrm{rel}}=K_\rho-K_0=-\epsilon\cos(q-Q)+\log\mathcal Z_\epsilon$ is indefinite for every $\epsilon\ne0$, since $\mathcal Z_\epsilon=I_0(\epsilon)<e^{|\epsilon|}$; its $\mu$-average $\log\mathcal Z_\epsilon$ is positive while its $\rho$-average $-S(\rho\|\rho_0)$ is negative. Its quantum counterpart $K_M-K_N$ is a nonnegative quadratic form, by \eqref{eq:modular-inequality} and Proposition~\ref{prop:sufficiency}.

\paragraph{Relative entropy and the information discarded:}
The same relative logarithm appeared in the phase construction at the beginning of this section. Here it generates a vector field, while its expectation gives
\begin{align}\label{eq:classical-relative-entropy-Hamiltonian}
    S(\rho\|\rho_0)=-\omega_\rho(K_{\mathrm{rel}})
    =H(\rho_0)-H(\rho).
\end{align}
The vector field alone does not determine this number: on a connected symplectic manifold it determines its Hamiltonian only up to a constant. The normalized densities fix that constant. On a general Poisson manifold, Casimirs give a larger ambiguity.

Information loss and failure of autonomous coarse dynamics are also distinct. For a product state, the relative transformation acts only on $Z$, yet
\begin{align}
    S(\rho_Y\rho_Z\|\rho_Y\otimes1)
    =\int_Z\rho_Z\log\rho_Z\,d\mu_Z
\end{align}
can be positive. Coarse-graining has discarded the hidden distribution even without correlations. More generally,
\begin{align}
    S(\rho\|\rho_Y\otimes1)
    =S(\rho\|\rho_Y\rho_Z)+S(\rho_Z\|1),
\end{align}
where $\rho_Z$ is the other marginal. The first term is the mutual information and the second measures the hidden marginal relative to the chosen reference. Neither this entropy nor the relative transformation supplies a monotone entropy-production law without an additional coarse-graining or dynamical assumption.

\paragraph{Subregions and the passage to field theory:}
The measurable region algebras $\mathcal M_O$ remain valid coarse descriptions. What can fail is the smoothness needed to assign their coarse densities a Hamiltonian flow. For the region expectation $E_O$,
\begin{align}
    \rho_0=E_O(\rho)=\chi_O\rho+c_O\chi_{O^c},\qquad
    c_O=\frac{1}{\mu(O^c)}\int_{O^c}\rho\,d\mu.
\end{align}
For an open region with smooth boundary this density is generally discontinuous there. In the interior of $O$, $K_0=K_\rho$; in the interior of the complement, $K_0$ is constant. These local statements do not specify a smooth global Hamiltonian flow across the boundary. Additional regularization or boundary conditions are needed. This is a limitation of the Hamiltonian construction, not of the measure-theoretic net.

Finally, every classical coarse algebra still admits a state-preserving conditional expectation for a faithful probability state, even when its Poisson state flow fails to preserve that algebra. In the product example this expectation is explicitly
\begin{align}
    E_{\rho,Y}(f)(y)=\int_Z f(y,z)\frac{\rho(y,z)}{\rho_Y(y)}\,d\mu_Z(z).
\end{align}
Its existence does not imply \eqref{eq:classical-state-intertwining}. In quantum modular theory the two questions merge: by Takesaki's theorem (Section~\ref{subsec:quantum-coarse} and Theorem~\ref{thm:correctable}), a reference-preserving expectation onto a subalgebra exists if and only if the subalgebra is invariant under the modular flow of the reference. This is the additional structure that the quantum theory contributes to the comparison of full and coarse dynamics.

The classical construction thus motivates studying how relative state evolutions transport retained observables. In QFT a spacetime region selects local field observables, whereas a region $O\subset X$ above selects possible microstates. The connection to the subregion algebras of QFT, developed in the companion paper~\cite{companionQFT}, is through inclusions of observable algebras and their state-dependent flows; it does not identify these two notions of region.

\paragraph{A classical half-sided inclusion:}
A simple example of strict nesting is supplied by a classical chiral current $j(x)$ on the line. On rapidly decreasing smooth configurations, choose
\begin{align}
    \{j(x),j(y)\}=-\partial_x\delta(x-y),
    \qquad P=\frac12\int_{\mathbb R}j(x)^2\,dx,
    \qquad \tau_a j(x)=j(x+a)\ .
\end{align}
The nonnegative Hamiltonian $P$ generates $\tau_a$ through $\{P,j(x)\}=\partial_xj(x)$. Let $\mathcal N_b$ be generated by $1$ and $j(f)$ with $f\in C_c^\infty((b,\infty))$. Since $\tau_a j(f)=j(f(\,\cdot-a))$,
\begin{align}
    \tau_a(\mathcal N_b)=\mathcal N_{b+a}\subsetneq\mathcal N_b\quad(a>0),
    \qquad
    \bigcup_{b\in\mathbb R}\mathcal N_b=\mathcal M_{\mathrm{loc}},
    \qquad
    \bigcap_{b\geq0}\mathcal N_b=\mathbb R1\ .
\end{align}
Here $\mathcal M_{\mathrm{loc}}$ is the polynomial algebra of compactly smeared currents; $P$ is a global functional generating its automorphisms. The full translation group is reversible, but its positive half acts on $\mathcal N_0$ by injective, nonsurjective Poisson endomorphisms. The decreasing family progressively discards accessible field degrees of freedom.

The dilation Hamiltonian $D=\frac12\int_{\mathbb R}xj(x)^2\,dx$ gives a further connection with half-sided modular inclusions. Its flow $\delta_s$ satisfies
\begin{align}
    \{D,P\}=P,
    \qquad \delta_s(\mathcal N_b)=\mathcal N_{e^s b},
    \qquad \delta_s\tau_a\delta_{-s}=\tau_{e^s a}\ .
\end{align}
Thus dilations preserve $\mathcal N_0$ and act half-sidedly on $\mathcal N_b$ for $b>0$. These are classical Hamiltonian realizations of the translation--dilation relations. Their construction uses the Poisson structure; identifying the dilation with a state-generated modular evolution requires additional input.

\paragraph{Classical error correction and the classical sufficiency gap.}
The constructions of this section can be read in the language that Section~\ref{sec:QEC} will use for the quantum theory. Restriction of a density to the coarse algebra $\mathcal N=L^\infty(X,\Sigma_0,\mu)$ is an erasure channel. The reference-preserving expectation $E_N$ always exists, and the densities it preserves,
\begin{align}
 \mathcal C_N=\{\rho:\ E_N(\rho)=\rho\},
\end{align}
are exactly the $\Sigma_0$-measurable densities; in the product setting these are the densities $\rho_Y(y)\,1$ whose hidden marginal is the reference $\mu_Z$. Every such density is recovered from its restriction by $E_N$ itself, so $\mathcal C_N$ is an exactly correctable code for the erasure, and its logical algebra is $\mathcal N$. For a general density, the chain rule $S(\rho\|1)=S(\rho\|\rho_0)+S(\rho_0\|1)$ with $\rho_0=E_N(\rho)$, which uses only that $\log\rho_0$ is $\Sigma_0$-measurable, shows that
\begin{align}\label{eq:classical-sufficiency-gap}
 S(\rho\|\rho_0)=S(\rho\|1)-S(\rho_0\|1)=-\omega_\rho(K_{\mathrm{rel}})\ \ge0
\end{align}
is the loss of relative entropy to the reference under restriction to $\Sigma_0$. By the classical equality case of monotonicity, this vanishes if and only if $\Sigma_0$ is a sufficient statistic for the pair $(\rho,\mu)$, that is, if and only if $\rho\in\mathcal C_N$~\cite{OhyaPetz1993}. The expectation value of the relative Hamiltonian, which generates the classical subsystem flow \eqref{eq:classical-relative-transformation}, is therefore the classical sufficiency gap. Two features distinguish this from the quantum case. The correctable code exists for every coarse algebra, independently of the dynamics; whether the retained observables evolve autonomously is a separate question answered by \eqref{eq:classical-product-state-closure}. And there is a single gap, because a commutative algebra has no commutant distinct from itself. In the quantum theory both features change: the existence of a reference-preserving expectation becomes equivalent to autonomy (Takesaki), and the subsystem-flow generator acquires a second gap for the inclusion of commutants (Proposition~\ref{prop:sufficiency}).

\section{Operator algebra background}\label{section:basics}

The observables of a physical system form a von Neumann algebra. Classical observables are in the center of this algebra. A quantum system with no non-trivial classical observables is a von Neumann \emph{factor}. For a finite quantum system, such factors are
finite matrix algebras\footnote{The factor associated to a $d$-level quantum system is the algebra of $d\times d$ complex matrices, $M(d,\mathbb{C})$}. 
In this work, we are primarily interested in the infinite-dimensional factors describing subsystems of many-body quantum systems in the thermodynamic and continuum limit. Infinite-dimensional factors are classified as type $I_\infty$, type $II_1$, type $II_\infty$, and type $III$. As we will see below, this classification is based on the ergodic properties of the modular flow of states. For a review of the definition of a von Neumann algebra, modular flow and basics of modular theory see Appendix \ref{}.

\subsection{Algebra types and ergodicity of modular flow}

Quantum systems come naturally equipped with a notion of modular dynamics that associates to every faithful normal state $\omega$ \footnote{A normal state is the infinite-dimensional generalization of a standard quantum density matrix on the GNS Hilbert space. Infinite-dimensional Von Neumann algebras can have purely abstract, singular states that exist mathematically that cannot be approximated by density matrices. An example in the Hilbert space of a simple harmonic oscillator is the state $\omega(a)=\lim_{n\to \infty}\braket{n|a|n}$. It is singular in that it has vanishing expectation value for all compact operators. Requiring a state to be {\it normal} explicitly rules out these unphysical states.} a unitary flow
$\sigma^{\omega}_t(a)=\Delta_{\omega;M}^{it}a\Delta_{\omega;M}^{-it}$. 
For matrix algebras, every faithful normal state $\omega$ corresponds to a full rank density matrix $\rho_\omega$. it is $\sigma^{\omega}_t(a)=\rho_\omega^{it}a\rho_\omega^{-it}$, where $\rho_\omega$ is the density matrix associated with state $\omega$. This flow preserves the algebra.

The set of operators that are invariant under modular flow (conserved modular charges) forms an algebra of modular charges \footnote{The subalgebra of modular charges is called {the centralizer} in mathematical jargon.}. Using the ergodic properties of modular flows, one classifies infinite-dimensional factors. To state the results easily, we consider connected modular two-point functions:
\begin{align}
    f_{ab,con}^{(\omega)}(t)=\omega(a_1 \sigma_\omega^t(a_2))-\omega(a_1)\omega(a_2)\ .
\end{align}
If there are classical observables, they are inside the intersection of modular charges of all states. Since our focus is modular ergodicity, we focus on von Neumann factors. They can be classified as
\begin{enumerate}
    \item {\bf Type I$_\infty$:} The modular flow of any state is quasi-periodic, never ergodic, and there are Poincaré recurrences. In other words, in any state, all modular two-point functions are quasi-periodic.
    \item {\bf Type II} There exist states where certain modular two-point functions can decay, but the generator of modular flow (modular Hamiltonian) is inside the algebra, which means that their correlators are time-independent and can never decay. Further, if there exists a state (finite trace) with trivial modular flow, we call it type II$_1$, otherwise, it is type II$_\infty$. 
    \item {\bf Type III:} The modular Hamiltonian is never inside the algebra. However, it could still be possible that for select $t$, there exist operators in the algebra such that $\sigma_\omega^t(a)=u_t a u_t^\dagger$. If this never occurs, we say the algebra is type III$_1$. Type III algebras admit no traces or density matrices.
\end{enumerate}
In a theory of ergodic theory specialized to modular flow of states (modular ergodic theory), type III$_1$ factors are the most mixing and chaotic. In physics, the local algebras of causally complete regions in QFT are type III$_1$ factors. 


\subsection{Modular flows versus arbitrary unitary dynamics}\label{sec:modular-vs-dynamics}

% Of course, modular dynamics depends on the state that defines the modular flow.  In matrix algebras, every unitary one-parameter family $U_t=e^{i H t}$ can be viewed as the modular flow for the state $\rho=e^{-H}/\tr(e^{-H})$ \footnote{More formally, any weakly-continuous one-parameter automorphism of the matrix algebra $M_d(\mathbb{C})$ is the modular flow of a faithful state on this algebra.}. Therefore, by focusing on the modular flow of arbitrary states, we are considering general closed dynamics.  In a type I algebra, all automorphisms are inner, meaning there always exists unitaries $u_t$ in the algebra such that 
% \begin{align}
%     \alpha_t(a)=u_t a u_t^\dagger
% \end{align}
% We can also always define the generator $H$ using $u_t=e^{i H t}$, however, $H$ might be unbounded, or more generally it could be that $\tr(e^{-H})=\infty$. Relaxing the normalizability condition $\omega(1)=1$, generalizes our discussion to unnormalized states (weights). In type I algebras, every automorphism can be viewed as the modular flow of a weight.

% In type II algebras, all modular flows are inner; therefore, any dynamics described by outer automorphisms lies beyond modular flows. In type III$_1$ algebras that we are interested in this paper, the modular flow is always outer, and any inner dynamics cannot be viewed as a modular flow.
% Ergodic modular flows are so special that the only quantum systems that admit any (faithful normal) states with such property are exotic infinite-dimensional algebras: type III$_1$ factors~\cite{Connes1973,Longo1982}\footnote{Note that the converse fails, meaning that even in a type III$_1$ algebra, not every faithful normal state has ergodic modular flow; i.e., a trivial centralizer. An example is the Araki-Woods type III$_1$ factors.}.

In a finite quantum system, a closed dynamics and a thermal state are the same piece of data. A Hamiltonian $H\in M_d(\mathbb C)$ generates $\alpha_t=\operatorname{Ad}e^{itH}$, the Gibbs state $\rho_\beta=e^{-\beta H}/\operatorname{tr}e^{-\beta H}$ is the unique KMS state of $\alpha$ at inverse temperature $\beta$, and every faithful state is Gibbs at $\beta=1$ for $H=-\log\rho$, so it determines a dynamics $\sigma^\rho_t=\operatorname{Ad}\rho^{it}$ for which it is thermal. The second half of this correspondence in generalizes to infinite dimensions: every faithful normal state $\omega$ on a von Neumann algebra $\mathcal M$ has a \emph{modular flow} $\sigma^\omega_t=\operatorname{Ad}\Delta_\omega^{it}$, the unique one-parameter group of automorphisms of $\mathcal M$ for which $\omega$ is KMS at $\beta=1$~\cite{Takesaki1970}. 
Then, it is natural to ask: Given a dynamics $\alpha$ on an infinite-dimensional quantum system $\mathcal M$, is it the modular flow of some state? The answer depends on the type of $\mathcal M$ as we explain below.

\paragraph{Dynamics, inner and outer:}
A dynamics is a flow, a pointwise $\sigma$-weakly continuous one-parameter group $t\mapsto\alpha_t\in\operatorname{Aut}(\mathcal M)$. The automorphism $\alpha_t$ is \emph{inner} if it is $\operatorname{Ad}u$ for a unitary $u\in\mathcal M$, and the flow is inner if every $\alpha_t$ is. Inner dynamics is generated by a Hamiltonian inside the algebra, and an outer one is generated by something outside the algebra \footnote{On a factor with separable predual an inner flow is always of the form $\alpha_t=\operatorname{Ad}e^{itH}$ with $H$ a self-adjoint operator affiliated with $\mathcal M$, unique up to an additive constant~\cite{Kallman1971}.}. For example, the observables algebra of infinite spin chain in the infinite-temperature state is the hyperfinite $\mathrm{II}_1$ factor $R=\bigotimes_n(M_2,\operatorname{tr})$, the global rotation $\alpha_t=\bigotimes_n\operatorname{Ad}e^{it\sigma_z}$ is a continuous flow whose would-be generator, the total magnetization $\sum_n\sigma_z^{(n)}$, is not affiliated with $R$ \footnote{This unitary is outer for all $t\notin\pi\mathbb Z$, since $\bigotimes_n\operatorname{Ad}u_n$ is inner iff $\sum_n(1-|\operatorname{tr}u_n|)<\infty$ and here $|\operatorname{tr}e^{it\sigma_z}|=|\cos t|$.}.

% To set the notation, we use $\varphi$, $\psi$ for faithful normal semifinite weights (states included); weights are unavoidable because the natural ``Gibbs'' objects on infinite algebras often have infinite normalization.

The first observation is that the modular flow of any state in type I and type II algebras (semifinite algebras) is always inner. Let $\tau$ be the trace. Every weight is $\varphi=\tau(h\,\cdot\,)$ with $h$ positive, injective and affiliated with $\mathcal M$, $\sigma^\varphi_t=\operatorname{Ad}h^{it}$, and $\varphi$ is a state iff $\tau(h)<\infty$. Then,
\begin{equation}
 \alpha_t=\operatorname{Ad}e^{itH}\ \text{inner}
 \ \Longleftrightarrow\ \alpha=\sigma^{\varphi}\ \text{with}\ \varphi=\tau(e^{-H}\,\cdot\,),
 \qquad
 \alpha\ \text{state-modular}\ \Longleftrightarrow\ \tau(e^{-H})<\infty .
 \label{eq:semifinite-dictionary}
\end{equation}
% (With the physicist's sign $H\to-\beta H$ the last condition is $\tau(e^{-\beta H})<\infty$.) So on a semifinite factor a dynamics admits a KMS weight iff it is inner, and then at every temperature for which the Gibbs weight converges. The type refines this.
% \begin{theorem}[Takesaki, Connes]\label{thm:TC}
% \begin{enumerate}
%  \item[\rm(i)] The modular automorphisms $\sigma^\varphi_t$ are inner for every $t$ if and only if $\mathcal M$ is semifinite, i.e.\ admits a faithful normal semifinite trace~\cite{Takesaki1970}.
%  \item[\rm(ii)] Any two modular flows differ by an inner perturbation: $\sigma^\psi_t=\operatorname{Ad}u_t\circ\sigma^\varphi_t$ with $u_t=[D\psi:D\varphi]_t\in\mathcal M$ a unitary cocycle, $u_{t+s}=u_t\sigma^\varphi_t(u_s)$. Hence the set $T(\mathcal M)=\{t:\sigma^\varphi_t\in\operatorname{Inn}(\mathcal M)\}$ and the class $\delta_{\mathcal M}(t)=[\sigma^\varphi_t]\in\operatorname{Out}(\mathcal M)=\operatorname{Aut}/\operatorname{Inn}$ are independent of $\varphi$. Conversely every such cocycle comes from a unique weight $\psi$~\cite{Connes1973}.
%  \item[\rm(iii)] Consequently a flow $\alpha$ is the modular flow of some weight if and only if $[\alpha_t]=\delta_{\mathcal M}(t)$ for every $t$. It is the modular flow of a \emph{state} if and only if, in addition, the cocycle relating it to the modular flow of a state extends analytically to a strip of width $\tfrac12$~\cite{Connes1973}.
% \end{enumerate}
% \end{theorem}

% Part (ii) says that a factor has, up to inner readjustment, only \emph{one} modular dynamics: $\delta_{\mathcal M}$ is a canonical one-parameter subgroup of $\operatorname{Out}(\mathcal M)$, the ``non-commutative time'' of the algebra~\cite{ConnesTakesaki1977,ConnesRovelli1994}. Part (iii) says that there are no others, and separates states from weights by an analyticity condition, the abstract form of ``the partition function converges''. One corollary is invisible in finite dimensions. If $\alpha$ is modular, then the rescaled flow $t\mapsto\alpha_{\beta t}$ is modular iff $\alpha_{(\beta-1)t}$ is inner for all $t$, i.e.\ iff $(\beta-1)\mathbb R\subseteq T(\mathcal M)$: \emph{a modular flow can be run at a different temperature only on a semifinite algebra.} On a type III algebra the inverse temperature of a modular flow is an invariant of the flow, not a parameter.

%\paragraph{Semifinite algebras: modular means inner with a Gibbs weight.}


% \emph{Type $\mathrm I_d$.} Every automorphism is inner (Skolem--Noether) and $\rho\mapsto\operatorname{Ad}\rho^{it}$ is a bijection between faithful states and flows: dynamics and equilibrium are one datum.

% \emph{Type $\mathrm I_\infty$.} $B(\mathcal H)$ still has no outer automorphisms, and every flow $\operatorname{Ad}e^{itH}$ is the modular flow of the weight $\operatorname{tr}(e^{H}\,\cdot\,)$; but it is the modular flow of a \emph{state} only when $e^{H}$ is trace class.
% The free particle $H=P$ is thermal for a weight and for no normal state, and the trivial flow is modular only for the trace, there being no tracial state.

% \emph{Type $\mathrm{II}_1$.} Now $\tau$ is a state and $\sigma^\tau=\operatorname{id}$: the trace is the infinite-temperature state, and $\operatorname{Ad}e^{itH}$ is state-modular whenever $H$ is bounded above. Modular flows are exactly the inner flows. What is new is that outer automorphisms exist (global symmetry rotations, the Bernoulli shift on $\bigotimes_{n\in\mathbb Z}M_2$) and none of them is modular at any time, although all of them preserve $\tau$: a global symmetry of an infinite system is a dynamics with no thermal state.

In type $\mathrm{II}_\infty$ algebras, a new kind of outer automorphism appears: the trace-scaling automorphisms $\tau\circ\theta_\lambda=\lambda\tau$, which exist for every $\lambda>0$ on $R\otimes B(\ell^2)$. They are outer, since inner automorphisms preserve $\tau$, and thermal for nothing. The continuous trace-scaling flow $\theta_s$, $\tau\circ\theta_s=e^{-s}\tau$, is where type III comes from: $(R\otimes B(\ell^2))\rtimes_\theta\mathbb R$ is the injective $\mathrm{III}_1$ factor~\cite{ConnesTakesaki1977,Haagerup1987}.

\paragraph{Type III: modular flows are a canonical outer class:}
A type III factor has no trace; the picture inverts. No inner flow is modular, because a modular flow would then be inner for all $t$; conversely, every type III factor has outer automorphisms, supplied by its own modular flows at $t\notin T(\mathcal M)$ where the set $T(\mathcal M)$ is called the {\it Connes invariant}. The modular flows are precisely the flows lying in the canonical class $\delta_{\mathcal M}$, which is now a non-trivial one-parameter subgroup of $\operatorname{Out}(\mathcal M)$, and the three subtypes record its shape. For $\mathrm{III}_\lambda$, $0<\lambda<1$, the kernel is $T(\mathcal M)=(2\pi/|\log\lambda|)\mathbb Z$ and the class is a \emph{circle}: Connes' generalized trace is a weight whose modular flow is literally periodic, inner exactly once per period. For $\mathrm{III}_0$ the kernel is a countable subgroup, possibly dense, possibly trivial. For $\mathrm{III}_1$ the kernel is $\{0\}$: $\delta_{\mathcal M}$ is an injective line, every modular automorphism of every weight is outer for every $t\ne0$, and $\mathcal M$ is the crossed product of its $\mathrm{II}_\infty$ core by the trace-scaling flow. $\mathrm{III}_1$ is also the only type on which a state can have trivial centralizer, an ergodic modular flow~\cite{Connes1973,Longo1982}, though not every state of a $\mathrm{III}_1$ factor does.

For a physicist the relevant case is $\mathrm{III}_1$, the type of every local algebra in quantum field theory. The vacuum modular flow of a wedge is the boost at the Unruh temperature~\cite{BisognanoWichmann1975}; the same boost run at any other rate is the modular flow of no state and no weight on $\mathcal A(W)$, and no dynamics generated by an element of $\mathcal A(W)$ is modular at all. Modular dynamics is not, as in finite dimensions, ordinary dynamics with a thermal state attached; it is the one outer dynamics that every state of the algebra defines, up to inner readjustment, with its temperature fixed. 

% In semifinite algebras the obstruction to being modular is analytic, a partition function must converge; in type III it is topological, membership in one curve in $\operatorname{Out}(\mathcal M)$, inside which the ergodic hierarchy of the following sections lives (proofs and a table: Section~\ref{sec:aut-vs-modular}).

% The question we carry from classical dynamics is whether the retained observables evolve autonomously. We first ask this question for finite quantum systems.

% When the two algebras have a common cyclic and separating vector, we compare their modular evolutions through the subsystem flow $D_{N\subset M}(t)=\Delta_{\omega;M}^{-it}\Delta_{\omega;N}^{it}$ of Definition~\ref{def:subsystem-flow} below. We examine the relation between autonomy, conditional expectations, and generating inclusions. Throughout we read these in the language of error correction: the restriction of states to $\mathcal N$ is an erasure channel, a conditional expectation $\mE:\mathcal M\to\mathcal N$ is an encoding of states on $\mathcal N$ that the erasure undoes, and the question whether the retained observables evolve autonomously becomes the question whether the erasure is exactly correctable on the states the encoding produces (Section~\ref{sec:QEC}). These notions do not by themselves imply mixing; the geometric realization for vacuum subregion algebras of QFT is taken up in the companion paper~\cite{companionQFT}.

% Section~\ref{sec:classical} reviews the classical constructions. Two examples guide the comparison: a classical field region selects a subsystem of field observables, and a stationary marginal state need not give autonomous evolution of those observables. The latter distinction motivates comparing full and reduced state-generated flows.

\paragraph{Type III$_1$ modular flows can be ergodic:}

When the algebra $\mathcal M$ has a non-trivial center $Z\subset \mathcal M$, then this center is inside the centralizer of all modular flows:
\begin{align}
    \forall z\in Z,\forall t,\omega:\qquad \Delta_{\omega}^{it}z \Delta_\omega^{-it}=z
\end{align}
This means that these observables are conserved charges, and an obstruction to ergodicity in modular dynamics of arbitrary states. That is why, in quantum ergodic theory, it is natural to restrict to algebras with trivial center, i.e., factors.

For a fixed faithful normal state, modular ergodicity means that the centralizer is trivial:
\begin{align}
    \label{eq:modular-ergodicity}
    \mathcal M^\omega\equiv \{a\in \mathcal M:\Delta_{\omega;M}^{it}a\Delta_{\omega;M}^{-it}=a\qquad \forall t\}=\mathbb{C} 1
\end{align}

For a nontrivial factor, this requires type $III_1$, but it does not follow from the
type alone. For example, consider an infinite tensor product of identical three-level
systems with faithful site density matrix
$\rho=\sum_{i=1}^3p_i\ket{i}\bra{i}$.
If $\log(p_1/p_3)$ and $\log(p_2/p_3)$ are rationally independent, the resulting
Araki--Woods factor is type $III_1$~\cite{ArakiWoods1968}.
Nevertheless, the site spectral projections $\ket{i}\bra{i}_n$ belong to the
centralizer of the product state. 

In fact, states of type III$_1$ algebra with ergodic trivial centralizer satisfy a stronger condition: strong mixing (clustering in modular time) \cite{}.

\begin{lemma}[Strong mixing implies ergodicity]\label{lem:mixing-ergodic}
Let $\omega$ be a faithful normal state on a factor $\mathcal M$ and
\begin{equation}\label{eq:modular-correlation}
 C^\omega_{a,b}(t):=\omega\bigl(a^*\sigma^\omega_t(b)\bigr)-\omega(a^*)\omega(b),\qquad a,b\in\mathcal M,
\end{equation}
its connected modular correlation functions. If $C^\omega_{a,b}(t)\to0$ as $t\to\infty$ for all $a,b\in\mathcal M$ (strong mixing), then $\mathcal M^\omega=\mathbb C1$, and $\mathcal M$ is of type $III_1$.
\end{lemma}
\begin{proof}
For $a\in\mathcal M^\omega$ one has $\sigma^\omega_t(a)=a$, so $C^\omega_{a,a}(t)=\omega(a^*a)-|\omega(a)|^2=\omega\bigl((a-\omega(a))^*(a-\omega(a))\bigr)$ does not depend on $t$. Mixing forces it to vanish, and faithfulness gives $a=\omega(a)1$. The type statement is the result quoted above.
\end{proof}

% s:  Local algebras of a quantum field theory are denoted $\mathcal M(O)$. For a faithful normal state $\omega$ and an algebra $\mathcal C$ for which the reference vector is cyclic and separating we write $\Delta_{\omega;C}$, $J_{\omega;C}$ (or $\Delta_C$, $J_C$ when the state is clear) for the modular data, $\sigma^{\omega;C}_t=\operatorname{Ad}\Delta^{it}_{\omega;C}$ for the modular flow, $K_C=-\log\Delta_{\omega;C}$ for the modular Hamiltonian, $\beta_t=\sigma^{\omega;C}_{-t}$ for the physical modular flow generated by $K_C$, and $S_C(\varphi\|\omega)$ for the relative entropy of the restrictions to $\mathcal C$ (Section~\ref{sec:QEC}); subscripts carry the algebra's letter in roman type. When $\mathcal C$ is the full algebra, the label is dropped: $\sigma^\omega_t=\sigma^{\omega;M}_t$. The subsystem flow is $D_{N\subset M}(t)=\Delta_{\omega;M}^{-it}\Delta_{\omega;N}^{it}$ (Definition~\ref{def:subsystem-flow}), its compression to $\mathcal H_N=\overline{\mathcal N\Omega}$ is $C_{N\subset M}(t)$ (Appendix~\ref{app:compressed-comparison}), and $e_N$ is the projection onto $\mathcal H_N$.
% ---------------------------------------------------------------------------
% Section 3.3, edited. Changes relative to the submitted text are marked
% "% EDIT:" in the margin comments. Labels kept: subsec:quantum-coarse, eq:CE-bimodular,
% condexp, eq:quantum-coarse-entropy-decomposition, eq:quantum-coarse-autonomy,
% eq:takesaki-intertwining, thm:split-no-expectation, lem:expectation-slice,
% prop:conormal-factorization, eq:conormal-factorization.
% New labels: eq:slice-expectation-intro, eq:reduced-dynamics, eq:subsystem-flow-preview.
% ---------------------------------------------------------------------------

\subsection{Infinite-dimensional quantum subsystems (subfactors)}
\label{subsec:quantum-coarse}

In many-body finite quantum systems, the algebra of observables is the tensor product of single-body algebras. For a composite system $A=BR$ with $\mathcal H_{BR}=\mathcal H_B\otimes\mathcal H_R$ and density matrices $\rho_{BR},\rho_B,\rho_R$, the algebras are $\mathcal N=B(\mathcal H_B)\otimes1$, $\mathcal R=1\otimes B(\mathcal H_R)$ and $\mathcal M=\mathcal N\vee\mathcal R=B(\mathcal H_{BR})$.\footnote{Throughout the paper an inclusion of observable algebras is written $\mathcal N\subset\mathcal M$, $\mathcal R=\mathcal N'\cap\mathcal M$ is its relative commutant, and $\mathcal M'$, $\mathcal N'$ are the commutants on the Hilbert space. Roman capitals are reserved for physical systems.} We say that $B$ is a subsystem of $A$; both algebras are factors, so $\mathcal N\subset\mathcal M$ is a \emph{subfactor}. The algebra of $R$ is the \emph{relative commutant} of this inclusion, $\mathcal R=\mathcal N'\cap\mathcal M$.
% EDIT: was "$\mN^c=\mN'\cap\mM$" (old symbol) and "$\mathcal E_{\rho_R}(b\otimes r)=\omega_R(b)(b\otimes1)$" (wrong formula).
The partial trace over $R$ in a fixed state $\chi_R$ of the environment, read in the Heisenberg picture, is the unital completely positive projection
\begin{align}
    \mathcal E_{\chi_R}:\mathcal M\to\mathcal N,\qquad
    \mathcal E_{\chi_R}(b\otimes r)=\chi_R(r)\,(b\otimes 1)\ ,
    \label{eq:slice-expectation-intro}
\end{align}
the simplest example of a \emph{conditional expectation}; the general definition is given below.

Subfactors play a central role in modular ergodic theory. The modular flow of a faithful state $\omega$ of $A$ is a closed, automorphic dynamics of $\mathcal M$, but it can mix $B$ and $R$: from the point of view of the subsystem $B$ it is an open dynamics, in which the relative commutant $\mathcal R$ plays the role of the environment. The modular flow of $\omega$ transports the algebra of $B$ to
% EDIT: was $\sigma_{\omega;\mM}^{it}(\mN)$.
\begin{align}
    \mathcal N_t=\sigma^{\omega;M}_t(\mathcal N)\ .
\end{align}
Since in general $\mathcal N_t$ leaves $\mathcal N$, the effective dynamics seen in $B$ is, whenever a conditional expectation $\mathcal E:\mathcal M\to\mathcal N$ is available, the unital completely positive map
% EDIT: was $\mT_t(B\subset A;\omega)\equiv\mE\circ\sigma^t\circ\mE$; $\mathcal T$ now denotes the erasure channel, and the right-hand $\mE$ is redundant on $\mathcal N$.
\begin{align}
    \Phi^{\omega}_t:=\mathcal E\circ\sigma^{\omega;M}_t\big|_{\mathcal N}:\mathcal N\to\mathcal N\ ,
    \label{eq:reduced-dynamics}
\end{align}
which depends on the state $\omega$ and on the choice of $\mathcal E$, and which is not a semigroup in general. We say that the dynamics can be \emph{reduced} to $B$ if $\mathcal N_t=\mathcal N$ at all times: information in $B$ never leaks into the environment, and $\mathcal N$ is a dynamically \emph{invariant} subalgebra. By Takesaki's theorem, recalled below, this happens if and only if the two modular flows agree on $\mathcal N$,
% EDIT: was "$\mB_0$ ... $\sigma_{\omega;A}^t(b)=\sigma_{\omega;B}^t(b)$" (old notation).
\begin{align}
    \forall b\in\mathcal N:\qquad \sigma^{\omega;M}_t(b)=\sigma^{\omega;N}_t(b)\ ,
\end{align}
where $\sigma^{\omega;N}$ is the modular flow of the restriction $\omega|_N$. In Section~\ref{sec:QEC} we show that this is also the necessary and sufficient condition for $\mathcal N$ to be exactly correctable, by an $\omega$-preserving conditional expectation, against the erasure of its environment. Therefore the comparison
% EDIT: was "$\mathcal D_t=\mE\circ\sigma^{it}_{\omega;A}\circ\sigma^{-t}_{\omega;B}\circ\mE$ ... (see Definition~\ref{def:subsystemflow})": wrong order relative to Definition def:subsystem-flow, spurious $i$, dangling label.
\begin{align}
    b\longmapsto\sigma^{\omega;M}_{-t}\bigl(\sigma^{\omega;N}_t(b)\bigr),\qquad b\in\mathcal N,
    \label{eq:subsystem-flow-preview}
\end{align}
which evolves $b$ with the intrinsic modular flow of $B$ and then undoes the modular flow of $A$, measures the residual dynamics left by the mismatch between the modular flow of the full algebra $\mathcal M$ and that of the subalgebra $\mathcal N$. When $\Omega$ is cyclic for $\mathcal N$ it is implemented by the unitary $D_{N\subset M}(t)=\Delta_{\omega;M}^{-it}\Delta_{\omega;N}^{it}$, which we call the \emph{subsystem flow} (Definition~\ref{def:subsystem-flow}); in general it is implemented by a contraction, the compressed subsystem flow of Section~\ref{sec:QEC}. We argue that the subsystem flow plays a central role in the modular ergodic hierarchy (Section~\ref{subsec:QEC-hierarchy}), which classifies subalgebras from invariant (quantum error-correcting codes) through generating and forgetful to modular $K$-systems, the maximally chaotic case: modular $K$-systems are chaotic at the universal rate $2\pi$ (Section~\ref{subsec:geometry-modular-ergodic})~\cite{OFLLM2024}.
% EDIT: was "(Section~\ref{subsec:qeh})" and "it was shown in \cite{}" (dangling label, empty citation).
Here we review the operator-algebraic framework needed to study subsystem flows; the hierarchy itself is the subject of Section~\ref{sec:QEC} and of Section~\ref{subsec:QEC-hierarchy}.

\paragraph{Dynamics reducible to a subalgebra:} Given a group action $\alpha:G\to\operatorname{Aut}(\mathcal M)$, we say that the dynamics is reducible to $\mathcal N\subset\mathcal M$ if 
\begin{align}
    \forall g\in G,b\in\mathcal N:\qquad \alpha_g(b)\in \mathcal N\ . 
\end{align}
Then, the dynamics can be {\it reduced} to $\mathcal N$. In dynamical systems, this is sometimes called {\it autonomous dynamics}, as well. For the modular flow of $\omega$ this is the condition $\mathcal N_t=\mathcal N$ above.

\paragraph{Irreducible dynamics:} Given a general dynamical system, we say it is irreducible if there is no non-trivial proper subalgebra $\mathcal N\subset \mathcal M$ with reducible dynamics.

\paragraph{Quantum coarse-grainings.}
An inclusion $\mathcal N\subseteq\mathcal M$ specifies the retained observables.
When it exists, a \emph{normal conditional expectation}
$\mE_N:\mathcal M\to\mathcal N$ is a completely positive, unital projection
with the bimodule property
\begin{equation}
 \mE_N(bac)=b\,\mE_N(a)\,c,
 \qquad a\in\mathcal M,\quad b,c\in\mathcal N;
 \label{eq:CE-bimodular}
\end{equation}
% EDIT: "faithful" removed from the definition (it is an additional property, needed in Lemma lem:expectation-slice) and stated separately.
it is \emph{faithful} if $\mE_N(a^*a)=0$ implies $a=0$. A conditional expectation specifies how a state on $\mathcal N$ is extended to $\mathcal M$.
Unlike the classical conditional expectation in Section~\ref{sec:classical}, which exists for every coarse algebra,
a normal expectation onto an infinite-dimensional subfactor need not exist, as we show below. On the GNS Hilbert space, let $e_N$
be the orthogonal projection onto $\overline{\mathcal N\Omega}$.
If $\mE_N$ preserves the reference state $\varphi$, then
$e_Na\Omega=\mE_N(a)\Omega$; this is its Jones projection.
Without state preservation, the Hilbert-space projection does not represent the expectation\footnote{We will see in Section~\ref{sec:QEC} that the QEC interpretation of this is that $\mE_N$ is the encoding of logical states on $\mathcal N$
into states on $\mathcal M$, and restriction to $\mathcal N$ is the erasure it corrects.}.

\paragraph{Tensor-product subsystems:}
% EDIT: $\mathcal M=\mathcal N\otimes\mathcal R$ with $\mathcal N=B(\mathcal H_B)\otimes1$ was inconsistent with the opening paragraph; $E_{\chi_R}$ renamed $\mE_{\chi_R}$ as in \eqref{eq:slice-expectation-intro}.
For a finite composite system, $\mathcal H_{BR}=\mH_B\otimes\mH_R$ and
$\mathcal M=B(\mathcal H_B)\otimes B(\mathcal H_R)$. Retaining $\mathcal N=B(\mathcal H_B)\otimes1_R$ gives the expectations \eqref{eq:slice-expectation-intro},
\begin{equation}\label{condexp}
 \mE_{\chi_R}(b\otimes r)=\chi_R(r)\,(b\otimes1_R),
 \qquad b\in B(\mathcal H_B),\quad r\in B(\mathcal H_R),
\end{equation}
one for each state $\chi_R$ of the environment, and by Lemma~\ref{lem:expectation-slice} below these are all of them.
For a fixed expectation, $\varphi=\varphi\circ\mE_{\chi_R}$ precisely when
$\rho_{BR}=\rho_B\otimes\chi_R$. Thus $I_\rho(B:R)=0$ is necessary, but the hidden
marginal must also equal the chosen reference. More generally,
\begin{equation}\label{eq:quantum-coarse-entropy-decomposition}
 S(\rho_{BR}\|\rho_B\otimes\chi_R)
 =I_\rho(B:R)+S(\rho_R\|\chi_R).
\end{equation}
The two terms measure discarded correlations and discarded information about the
hidden distribution, just as in Section~\ref{sec:classical}.
The subspace and mixed finite-dimensional inclusions are reviewed in
Appendix~\ref{app:finite-coarse}. Our main examples are subfactors, including
local QFT algebras for which this tensor factorization need not hold.

\paragraph{Modular invariant algebras and quantum coarse-graining:}
For a faithful normal state $\omega$, autonomy of $\mathcal N\subset\mathcal M$ means
\begin{equation}\label{eq:quantum-coarse-autonomy}
 \mathcal N_t:=\sigma^{\omega;M}_t(\mathcal N)=\mathcal N,\qquad t\in\mathbb R.
\end{equation}
Takesaki's theorem~\cite{Takesaki1972} makes this equivalent to the existence of
an $\omega$-preserving normal expectation $\mE_\omega:\mathcal M\to\mathcal N$. In that case the two modular flows intertwine,
% EDIT: $V:\overline{\mathcal N\Omega}\hookrightarrow\mathcal H_{BR}$ corrected to $\mathcal H$ (the GNS space of $\omega$, not the system space).
\begin{align}
 \sigma^{\omega;M}_t(b)&=\sigma^{\omega;N}_t(b),\qquad b\in\mathcal N,\nn\\
 \Delta_{\omega;M}^{it}V&=V\Delta_{\omega;N}^{it},
 \qquad V:\overline{\mathcal N\Omega}\hookrightarrow\mathcal H.
 \label{eq:takesaki-intertwining}
\end{align}
This is the form in which Takesaki's theorem is used throughout; its error-correction form is Theorem~\ref{thm:correctable}.
The intrinsic modular operator of $\mathcal N$ acts on $\overline{\mathcal N\Omega}$.
In the tensor-product example, autonomy is equivalent to a product state
$\omega_N\otimes\omega_R$, the Takesaki condition of Section~\ref{sec:intro}; $\mE_{\omega_R}$ preserves it.
For a general inclusion, absence of correlations with $\mathcal N'\cap\mathcal M$ is not sufficient:
that relative commutant need not account for every discarded observable.
% EDIT: dangling "Appendix~\ref{app:compressed-comparison}" (the appendix is commented out) replaced by Section~\ref{sec:QEC}.
Section~\ref{sec:QEC} treats the subsystem flow compressed to
$\overline{\mathcal N\Omega}$, which is the object entering Theorem~\ref{thm:correctable}.

% EDIT: "Modular monotonicity" moved here from Section 5 and rewritten for a reference vector that is cyclic and
% separating for M but not necessarily cyclic for N. Labels eq:Tomita-extension, eq:modular-inequality,
% eq:characteristic-bounded are kept, so all existing cross-references continue to resolve.
\paragraph{Modular monotonicity:}
The modular operators of $\mathcal M$ and of $\mathcal N$ act on different spaces, $\mathcal H$ and $\mathcal H_N=\overline{\mathcal N\Omega}$, but they are tied together by the Tomita operators, which require no cyclicity of $\Omega$ for $\mathcal N$: since $\Omega$ is separating for $\mathcal M$ it is separating for $\mathcal N$, and it is cyclic for $\mathcal N$ on $\mathcal H_N$. Let $S_{\omega;M}=J_{\omega;M}\Delta_{\omega;M}^{1/2}$ on $\mathcal H$ and $S_{\omega;N}=J_{\omega;N}\Delta_{\omega;N}^{1/2}$ on $\mathcal H_N$ be the closures of $a\Omega\mapsto a^*\Omega$, $a\in\mathcal M$, and of $n\Omega\mapsto n^*\Omega$, $n\in\mathcal N$. For $n\in\mathcal N$,
\begin{equation}
  S_{\omega;M}\,Vn\Omega=n^*\Omega=V\,S_{\omega;N}\,n\Omega ,
  \label{eq:Tomita-extension}
\end{equation}
so the Tomita operator of $\mathcal M$ extends that of $\mathcal N$ along $V$: $S_{\omega;M}V\supseteq VS_{\omega;N}$, because $S_{\omega;M}$ is closed. Inserting the polar decompositions and using that the modular conjugations are antiunitary involutions,
\begin{equation}
  \Delta_{\omega;M}^{1/2}\,V\xi=J_{\omega;M}VJ_{\omega;N}\,\Delta_{\omega;N}^{1/2}\xi ,
  \qquad \xi\in\operatorname{Dom}\Delta_{\omega;N}^{1/2},
  \label{eq:modular-extension}
\end{equation}
where $J_{\omega;M}VJ_{\omega;N}$ is a linear isometry of $\mathcal H_N$ into $\mathcal H$. This one identity is the source of every monotonicity statement in the paper.

\begin{lemma}[Modular monotonicity]\label{lem:modular-monotonicity}
Let $\Omega$ be cyclic and separating for $\mathcal M$, let $\mathcal N\subset\mathcal M$, and write $\Delta_M=\Delta_{\omega;M}$, $\Delta_N=\Delta_{\omega;N}$.
\begin{enumerate}
 \item[\rm(i)] $V\operatorname{Dom}\Delta_N^{1/2}\subseteq\operatorname{Dom}\Delta_M^{1/2}$ and $\|\Delta_M^{1/2}V\xi\|=\|\Delta_N^{1/2}\xi\|$ for $\xi\in\operatorname{Dom}\Delta_N^{1/2}$. In the order of closed quadratic forms on $\mathcal H_N$,
 \begin{equation}
  \Delta_N\;\geq\;V^*\Delta_MV ,
  \label{eq:modular-inequality}
 \end{equation}
 where $V^*\Delta_MV$ denotes the closed form $\xi\mapsto\|\Delta_M^{1/2}V\xi\|^2$ on $V^{-1}\operatorname{Dom}\Delta_M^{1/2}$; for $e_N=1$ this is $\Delta_N\geq\Delta_M$.
 \item[\rm(ii)] For $0\le\alpha\le\tfrac12$, $V\operatorname{Dom}\Delta_N^{\alpha}\subseteq\operatorname{Dom}\Delta_M^{\alpha}$ and $\|\Delta_M^{\alpha}V\xi\|\le\|\Delta_N^{\alpha}\xi\|$ for $\xi\in\operatorname{Dom}\Delta_N^{\alpha}$. Hence $\Delta_M^{\alpha}V\Delta_N^{-\alpha}$, defined on the dense subspace $\operatorname{Ran}\Delta_N^{\alpha}$, extends to a contraction $W_\alpha:\mathcal H_N\to\mathcal H$,
 \begin{equation}
  \|\Delta_M^{\alpha}V\Delta_N^{-\alpha}\|\le1\quad(0\le\alpha\le\tfrac12),
  \qquad W_0=V,\qquad W_{1/2}=J_{\omega;M}VJ_{\omega;N},
  \label{eq:characteristic-bounded}
 \end{equation}
 and the two endpoints are isometries.
\end{enumerate}
\end{lemma}
\begin{proof}
(i) is \eqref{eq:modular-extension}, since $J_{\omega;M}VJ_{\omega;N}$ preserves norms; the form statement is the definition of the order of closed forms, with equality on the smaller domain. (ii) is the Heinz--Kato inequality~\cite{Kato1952}: if $A$ on $\mathcal H$ and $B$ on $\mathcal H_N$ are positive self-adjoint and $T:\mathcal H_N\to\mathcal H$ is a contraction with $T\operatorname{Dom}B\subseteq\operatorname{Dom}A$ and $\|AT\xi\|\le\|B\xi\|$ on $\operatorname{Dom}B$, then $T\operatorname{Dom}B^{s}\subseteq\operatorname{Dom}A^{s}$ and $\|A^{s}T\xi\|\le\|B^{s}\xi\|$ for $0\le s\le1$. Apply it with $A=\Delta_M^{1/2}$, $B=\Delta_N^{1/2}$, $T=V$ and $s=2\alpha$, using (i). For $V=1$ this is the L\"owner--Heinz theorem, the operator monotonicity of $z\mapsto z^{s}$ for $0\le s\le1$. The endpoints are $W_0=V$ and \eqref{eq:modular-extension}.
\end{proof}

The lemma is a quantum statement with no classical counterpart: in Section~\ref{subsec:classical-state-flow} the Poisson Hamiltonian of a state has no analogue of the Tomita operator and satisfies no order relation under coarse-graining. Its role in what follows is this. 
% For an inclusion of subfactors $\mN\subset \mM$ we define the {\it compressed subsystem flow} as
% \begin{equation}
%  C_{N\subset M}(t)=V^*\Delta_M^{-it}V\,\Delta_N^{it}
%  \quad\hbox{on }\mathcal H_N=\overline{\mathcal N\Omega},
%  \qquad V:\mathcal H_N\hookrightarrow\mathcal H_\mM ,
%  \label{eq:compressed-recall}
% \end{equation}
% % which coincides with the subsystem flow $D_{N\subset M}(t)$ when $\Omega$ is also cyclic for $\mathcal N$. 
% Here $\Delta_N$ is the modular operator of $\omega|_N$ acting on $\mathcal H_N$. Continued to imaginary modular time, $z=t+i\alpha$ with $0\le\alpha\le\tfrac12$, it reads
% \begin{equation}
%  C_{N\subset M}(t+i\alpha)=V^*\,\Delta_M^{-it}\,W_\alpha\,\Delta_N^{it}
%  =V^*X_\alpha(t)V,
%  \qquad X_\alpha(t):=\Delta_M^{-it}\,W_\alpha\,\Delta_N^{it}\,V^* ,
%  \label{eq:compressed-strip}
% \end{equation}
% so the imaginary-time values of the compressed subsystem flow are compressions to $\mathcal H_N$ of contractions $X_\alpha(t)$ on $\mathcal H$, which are partial isometries with initial space $\mathcal H_N$ on both edges of the strip, $\alpha=0$ and $\alpha=\tfrac12$. 
% The order relation \eqref{eq:modular-inequality} is the edge $\alpha=\tfrac12$ of this statement and the Heinz--Kato inequality fills in the strip. 
% For $e_N=1$ the edge operators are unitary, $D_{N\subset M}(t)=\Delta_M^{-it}\Delta_N^{it}$ and $D_{N\subset M}(t+\tfrac{i}{2})=\Delta_M^{-it}J_{\omega;M}J_{\omega;N}\Delta_N^{it}$, and \eqref{eq:characteristic-bounded} is Borchers' bound $\|\Delta_M^{\alpha}\Delta_N^{-\alpha}\|\le1$~\cite{revolutionizing}; analyticity in the interior of the strip and strong continuity on its boundary are proved there for $e_N=1$, and the argument, which uses only the uniform bound \eqref{eq:characteristic-bounded} and the two unitary groups, applies to \eqref{eq:compressed-strip} with $V$ inserted (Definition~\ref{def:subsystem-flow} and the lemma following it).

\paragraph{Common cyclic and separating vector:} A vector $\ket\Omega$ is cyclic and separating for a von Neumann algebra $\mathcal C$
if $\mathcal C\Omega$ is dense and $c\Omega=0$ implies $c=0$. For an inclusion we
call $\Omega$ a \emph{common} cyclic and separating vector if it has this property
for both algebras. Such vectors exist in abundance in the setting of interest:
every type $III$ von Neumann algebra on a separable Hilbert space is in standard
form~\cite{Blackadar2006}, and for two type $III$ algebras on a common separable
Hilbert space the set of common cyclic and separating vectors is a dense $G_\delta$
subset~\cite{dixmiermarechal}. The vacuum of a quantum field theory is a common
cyclic and separating vector for every pair of nested local algebras by the
Reeh--Schlieder theorem.

% EDIT: this paragraph was placed before the Takesaki paragraph; moved here so that the three mechanisms are announced immediately before the definitions and results that make them precise. "The two examples" (tensor amplification, adjacent local algebras) now live in Section 4, and the sentence says so.
In infinite dimensions the inclusion need not admit a tensor-product description, and three mechanisms obstruct conditional expectations: a trivial relative commutant $\mathcal N'\cap\mathcal M=\mathbb C1$ (irreducible subfactors, which cannot occur for a proper inclusion of type $I$ factors, and which null-translated wedge inclusions can provide in interacting QFT, subject to the theory's lightlike localization properties); the absence of a normal product state across $\mathcal N$ and $\mathcal R=\mathcal N'\cap\mathcal M$ for a conormal inclusion $\mathcal N\vee\mathcal R=\mathcal M$, which excludes a faithful normal expectation by Proposition~\ref{prop:conormal-factorization}; and the split property, the existence of a type $I$ factor between $\mathcal N$ and $\mathcal M$, which is statistical independence of $\mathcal N$ and $\mathcal M'$ rather than a dynamical condition and excludes normal expectations onto a type $III$ subalgebra by Theorem~\ref{thm:split-no-expectation}. The definitions and results that follow make these precise; two examples, tensor amplification and adjacent local algebras of a conformal net, are discussed in Section~\ref{sec:QEC}, and the error-correction reading of all three mechanisms is collected in Section~\ref{subsec:QEC-hierarchy}. For finite algebras a faithful normal trace gives a trace-preserving expectation onto every von Neumann subalgebra; a type $I$ ambient factor does not, in general, admit a normal expectation onto an arbitrary subfactor.

\paragraph{Conormal inclusions:} An inclusion of von Neumann algebras $\mathcal N\subset \mathcal M$ is \emph{conormal} if
\begin{equation}
 \mathcal N\vee\mathcal R=\mathcal M.
\end{equation}
\paragraph{Standard and split inclusions:}
An inclusion $\mathcal N\subset\mathcal M$ is
\emph{standard} with respect to $\Omega\in\mathcal H$ if $\Omega$ is cyclic
and separating for $\mathcal M$, $\mathcal N$ and the relative commutant $\mathcal R=\mathcal N'\cap \mathcal M$. It is \emph{split} if a type I factor $\mathcal F$ exists with
$\mathcal N\subset\mathcal F\subset\mathcal M$.

\begin{theorem}[Split inclusions and normal expectations]\label{thm:split-no-expectation}
If $\mathcal N$ is a type III factor and
$\mathcal N\subset\mathcal M$ is split, then no normal conditional expectation
$\mE:\mathcal M\to\mathcal N$ exists.  
\end{theorem}

\begin{proof}
Let $\mathcal N\subset\mathcal F\subset\mathcal M$ be an interpolating type I
factor and suppose that $\mE:\mathcal M\to \mathcal N$ is normal.  Its restriction
\begin{equation}
 \mE_{\mathcal F}:=\mE|_{\mathcal F}:\mathcal F\longrightarrow\mathcal N
\end{equation}
is a normal norm-one projection with range $\mathcal N$.  Tomiyama's type
theorem states that a normally expected von Neumann subalgebra of a type I von
Neumann algebra is type I \cite{Tomiyama1959}\cite{Takesaki1972}. This contradicts $\mathcal N$ being a type III factor.
\end{proof}

Faithfulness of $\mE$ and preservation of a distinguished state play no role in this theorem, and neither does standardness.

\begin{lemma}\label{lem:expectation-slice}
If $\mathcal N$ is a factor and $\mE:\mathcal M\to\mathcal N$ is a faithful
normal conditional expectation, there is a faithful normal state $\chi_R$ on
the relative commutant $\mathcal R$ such that
\begin{equation}
 \mE(c)=\chi_R(c)\mathbbm 1,
 \qquad \mE(b c)=\chi_R(c)b,
 \qquad b\in\mathcal N,\quad c\in\mathcal R.
\end{equation}
\end{lemma}

\begin{proof}
For $b\in\mathcal N$ and $c\in\mathcal R$,
\begin{equation}
 b \mE(c)=\mE(b c)=\mE(c b)=\mE(c)b.
\end{equation}
Thus $\mE(c)\in Z(\mathcal N)=\mathbb C\mathbbm 1$, so
$\mE(c)=\chi_R(c)\mathbbm 1$ for a normal state $\chi_R$.  If $c\geq0$ and
$\chi_R(c)=0$, faithfulness of $\mE$ implies $c=0$; hence $\chi_R$ is faithful.  The
second identity is the bimodule property.
\end{proof}

For a faithful normal state $\varphi$ on $\mathcal N$, the state
$\varphi_0:=\varphi\circ \mE$ is faithful and normal and satisfies
\begin{equation}
 \varphi_0(b c)=\varphi_0(b)\chi_R(c).
\end{equation}

\begin{proposition}[conormality and tensor product]\label{prop:conormal-factorization}
Under the hypotheses of Lemma~\ref{lem:expectation-slice}, multiplication
extends to a normal $*$-isomorphism
\begin{equation}
 m:\mathcal N\overline\otimes\mathcal R
 \overset{\cong}{\longrightarrow}\mathcal N\vee\mathcal R,
 \qquad m(b\otimes c)=b c.
\end{equation}
If the inclusion is conormal, then
\begin{equation}\label{eq:conormal-factorization}
\mathcal M\cong
 \mathcal N\overline\otimes\mathcal R .
\end{equation}
\end{proposition}

\begin{proof}
Multiplication is a $*$-homomorphism on $\mathcal N\odot\mathcal R$ because
the two algebras commute.  Moreover,
\begin{equation}
 \varphi_0\bigl(m(b\otimes c)\bigr)
 =(\varphi\otimes\chi_R)(b\otimes c).
\end{equation}
The GNS representation of the faithful normal product state
$\varphi\otimes\chi_R$ is the spatial tensor-product representation.  Hence
multiplication extends normally to $\mathcal N\overline\otimes\mathcal R$.
Faithfulness of the product state gives injectivity, while surjectivity follows
from the definition of $\mathcal N\vee\mathcal R$.  Conormality then yields
\eqref{eq:conormal-factorization}; compare
\cite{Faulkner2020,FaulknerSperanza2024}.
\end{proof}

\paragraph{Index of an inclusion:} The Pimsner--Popa index $\operatorname{Ind}(\mE)$ of a conditional expectation $\mE:\mathcal M\to \mathcal N$ is the infimum of $\lambda\geq1$ such that, for all positive operators $x_+\in\mathcal M$,
\begin{align}
    \mE(x_+)\geq \lambda^{-1}x_+\ .
\end{align}
If no finite bound exists, $\operatorname{Ind}(\mE)=\infty$. For an inclusion of factors, the minimal index is
\begin{align}
    [\mathcal M:\mathcal N]_0=\inf_\mE \operatorname{Ind}(\mE)\ .
\end{align}
The index is a measure of the relative size of $\mathcal M$ to $\mathcal N$. When the relative commutant is infinite dimensional, or there are no normal conditional expectations, we say the index is infinite.

Finite index does not by itself specify a matrix relative commutant or conormality, and infinite index does not force the relative commutant to be infinite-dimensional. The index measures the relative size of $\mathcal M$ and $\mathcal N$ through the best constant in $\mE(x)\ge\lambda^{-1}x$; when no normal expectation exists, the index is infinite.

\paragraph{Local algebras of QFT (no normal conditional expectations):}
In a local conformal net $\{\mathcal{A}(I)\}$ on the circle, let $I_1,I_2$ be disjoint
open intervals with closures sharing exactly one endpoint, and set $I=\operatorname{int}\overline{I_1\cup I_2}$.
Haag duality and strong additivity give $\mathcal{A}(I_1)'\cap\mathcal{A}(I)=\mathcal{A}(I_2)$
and $\mathcal{A}(I_1)\vee\mathcal{A}(I_2)=\mathcal{A}(I)$, so the inclusion
$\mathcal N:=\mathcal{A}(I_1)\subset\mathcal{M}:=\mathcal{A}(I)$ is conormal with
$\mathcal R=\mathcal A(I_2)$. The adjacent pair has no spatial tensor-product
factorization in a nontrivial conformal net~\cite{Wiesbrock1993}.
Proposition~\ref{prop:conormal-factorization} therefore excludes faithful normal
expectations. The same conclusion holds for a nonfaithful normal expectation:
composing it with a normal state on $\mathcal N$ would give a normal product
state on $\mathcal N\vee\mathcal R$, whereas existence of such a state for commuting
factors entails a spatial tensor-product representation~\cite{DoplicherLongo1984}.
Thus no normal expectation $\mE:\mathcal A(I)\to\mathcal A(I_1)$ exists.
The obstruction concerns normal product states and expectations for this inclusion;
it does not exclude error correction for other channels or chosen families of states.


\section{Invariant subsystems as QEC codes}\label{sec:QEC}

Sections~\ref{sec:classical} and~\ref{section:basics} identified two objects attached to an inclusion of von Neumann algebras $\mN\subset \mM$ and a faithful normal state $\omega$: the autonomy condition $\sigma^{\omega}_t(\mathcal N)=\mathcal N$ and the existence of an $\omega$-preserving normal conditional expectation $\mE:\mathcal M\to\mathcal N$. Here, we start by reviewing Takesaki's theorem, which ties the two together, and provides an operational interpretation for invariant subsystems in the language of operator-algebraic quantum error correction~\cite{Faulkner2020,FaulknerSperanza2024}. We define {\it subsystem flow} and argue it is the natural quantitative measure of the failure of error correction. 
% Throughout, $\mathcal N\subset\mathcal M$ is an inclusion of factors on a Hilbert space $\mathcal H$, $\Omega$ is cyclic and separating for $\mathcal M$, $\omega=\braket{\Omega|\cdot\,\Omega}$, and $\mathcal R=\mathcal N'\cap\mathcal M$.




\subsection{QEC and subsystem flow in matrix algebras}
A finite composite system has observable algebra $\mathcal M=\mathcal N\otimes\mathcal R$. A general unitary flow on $\mathcal M$ sends $\mathcal N_0\equiv \mathcal N\otimes 1_R$ and $\mathcal R_0\equiv 1_B\otimes \mathcal R$ to a new embedding
\begin{align}
&\mathcal N_t=\iota_t(\mathcal N)\equiv U_t(\mathcal N\otimes 1_R)U_t^\dagger\subset\mathcal M  \nn\\
&\mathcal R_t=\iota_t^c(\mathcal R)\equiv U_t(1\otimes \mathcal R)U_t^\dagger
\end{align}
where the map $\iota_t^c$ is conjugate to $\iota_t$. It results in new decompositions $\mathcal M=\mathcal N_t\mathcal R_t$ with $[\mathcal N_t,\mathcal R_t]=0$.
Every such flow can be realized as the modular flow of a state on $\mathcal M$ defined by
\begin{align}
    \mathcal N_t=\rho_{BR}^{it}\mathcal N_0\rho_{BR}^{-it},\qquad \rho_{BR}=e^{-i\log U_1}.
\end{align}
For this dynamics to preserve $\mathcal N$ it must be that $\mathcal N_t=\mathcal N_0$ for all time, which occurs if and only if the state $\rho_{BR}$ is uncorrelated between $\mathcal N$ and $\mathcal R$:
\begin{align}
 &   \rho_{BR}=\rho_B\otimes \rho_R:\qquad     \rho_{BR}^{it}(b\otimes 1)\rho_{BR}^{-it}=\rho_B^{it}b\rho_B^{-it}\otimes 1_R.
\end{align}
and there is a one-to-one correspondence between states of $R$ and projections:
\begin{align}
    \mE_{\rho_R}:\mathcal M\to \mathcal N_0, \qquad \mE_{\rho_R}(b\otimes r)=(b\otimes 1_R)\: \tr(\rho_R r)
\end{align}
which leaves the uncorrelated state invariant:
\begin{align}
    \rho_{BR}=\rho_B\otimes \rho_R:\qquad \tr(\rho_{BR}\mE_{\rho_R}(a))=\tr(\rho_R a)
\end{align}
We also have
\begin{align}\label{KnillLaff}
    \mE_{\rho_R}(1_B\otimes r)=\mE_{\rho_R}(1_B\otimes 1_R)\tr(\rho_R r)=(1_B\otimes 1_R)\tr(\rho_R r)
\end{align}
The condition for autonomous evolution of subalgebra $\mathcal N_0$ identifies an uncorrelated state $\rho_{BR}=\rho_B\otimes \rho_R$. In the language of QEC, this is the condition that the erasure of $\mathcal R_0$ is unitarily correctable on the code subalgebra $\mE_{\rho_R}(\mathcal M)=\mathcal N_0$, and the equation \eqref{KnillLaff} is the algebraic version of the Knill-Laflamme condition for subsystem QEC codes. It also implies 
\begin{align}
    \mE_{\rho_R}(\rho_{BR}^{it}(1\otimes r)\rho_{BR}^{-it})=\mE_{\rho_R}(1\otimes r)=(1_B\otimes 1_R)\tr_R(\rho_R r)
\end{align}
which is a totally forgetful map, erasing all the information in $R$. 

The leakage of information from $\mathcal N_0$ into $\mathcal R_0=1_B\otimes \mathcal R$ can be attributed to correlations in the state $\rho_{BR}$ between $B$ and $R$. The information-disturbance theorem in QEC says that an encoding map $\mE_{\rho_R}(b\otimes 1)$ is a good QEC code to protect $\mathcal N$ if it is a good forgetful map on the relative commutant $\mathcal R$. 



\subsection{Modular invariant subsystem as exact correctability}\label{subsec:QEC-takesaki}

Here, we generalize the discussion of matrix algebras above to arbitrary von Neumann algebras. 

\paragraph{The erasure channel.}
In the Schr\"odinger picture, restriction of normal states
\begin{equation}
 \mathcal T:\mathcal M_*\to\mathcal N_*,\qquad \varphi\mapsto\varphi|_N,
 \label{eq:erasure-channel}
\end{equation}
is the \emph{erasure channel} associated with the inclusion: an observer confined to $\mathcal N$ has lost access to every observable not in $\mathcal N$. In the Heisenberg picture it is the inclusion map $\mathcal N\hookrightarrow\mathcal M$. When $\mathcal M=\mathcal N\vee\mathcal R$ and multiplication gives a spatial tensor factorization $\mathcal M\cong\mathcal N\overline\otimes\mathcal R$, this is the ordinary erasure of the factor $\mathcal R$ and $\varphi\mapsto\varphi|_{\mathcal R}$ is its complementary channel. Without such a factorization, the relative commutant need not account for every discarded observable, and the erasure is simply the restriction \eqref{eq:erasure-channel}. Section~\ref{subsec:QEC-KL} identifies the complementary channel for a general inclusion and states correctability in Knill--Laflamme form.

\paragraph{Encoding by a conditional expectation.}
A normal conditional expectation $\mE:\mathcal M\to\mathcal N$ specifies how a state on $\mathcal N$ is extended to $\mathcal M$:
\begin{equation}
 \mE_*:\mathcal N_*\to\mathcal M_*,\qquad \varphi_N\mapsto\varphi_N\circ\mE .
 \label{eq:encoding}
\end{equation}
Since $\mE$ is a projection onto $\mathcal N$, $(\varphi_N\circ\mE)|_N=\varphi_N$, so $\mathcal T\circ\mE_*=\operatorname{id}_{\mathcal N_*}$. Thus $\mE_*$ is an \emph{encoding} of the logical states $\mathcal N_*$ into $\mathcal M_*$ that is exactly correctable against the erasure \eqref{eq:erasure-channel}, with recovery map $\mE_*$ itself:
\begin{equation}
 \mathcal C_{\mE}:=\mE_*(\mathcal N_*)=\{\varphi\in\mathcal M_*:\ \varphi\circ\mE=\varphi\}
 \label{eq:code}
\end{equation}
is the \emph{code} of $\mE$. It is a norm-closed convex set of normal states, and the logical algebra is $\mathcal N$ itself: $\mE(n)=n$ for $n\in\mathcal N$. In the Heisenberg picture $\mE(n)=n$ says that the logical observables themselves are untouched; this much is automatic for a restriction, and the content of correctability is that the state on all of $\mathcal M$, correlations included, is reconstructed from its restriction (Section~\ref{subsec:QEC-KL}). Which subalgebras of $\mathcal M$ admit such an expectation is not a kinematic question; by Takesaki's theorem it is the dynamical question of modular autonomy, and it is this that distinguishes the quantum from the classical theory, where every subalgebra is expected with respect to the reference measure (Section~\ref{sec:classical}).

By Lemma~\ref{lem:expectation-slice}, every faithful normal expectation onto a factor is determined on $\mathcal N\vee\mathcal R$ by a faithful normal state $\chi$ on $\mathcal R$ through $\mE(nr)=\chi(r)\,n$. In the tensor-product case the code \eqref{eq:code} is the set of product states $\varphi_N\otimes\chi$ with the hidden factor fixed; different choices of $\chi$ give different codes with the same logical algebra. The entropy decomposition \eqref{eq:quantum-coarse-entropy-decomposition},
\[
 S(\rho_{BR}\|\rho_B\otimes\chi_R)=I_\rho(B:R)+S(\rho_R\|\chi_R),
\]
separates the two ways a state can fail to lie in the code: correlations between the logical and hidden factors, and a hidden marginal different from the reference $\chi$.

Note that in passive QEC, the goal is to choose an encoding $\mE_*$ such that all errors commute with the encoded subalgebra. In a sense, the encoding $\mE_*$ is the recovery map in passive QEC.

\paragraph{State preservation and zero coarse-graining entropy.}
A state $\omega$ lies in the code of $\mE$ precisely when
\begin{equation}
 S(\omega\|\omega\circ\mE)=0\quad\Longleftrightarrow\quad\omega=\omega\circ\mE .
 \label{eq:QEC-entropy}
\end{equation}
This is equality of a state with its specified extension. It does not by itself establish a code: exact error correction requires a channel, a recovery map, and a family of states on which their composition is the identity. The family here is \eqref{eq:code}, and \eqref{eq:QEC-entropy} says that $\omega$ belongs to it.

The following theorem collects the equivalent forms of the correctable case. 
\begin{theorem}[Correctable inclusions]\label{thm:correctable}
Let $\mathcal N\subset\mathcal M$ be an inclusion of factors, $\Omega$ cyclic and separating for $\mathcal M$ and separating for $\mathcal N$, and $\omega=\braket{\Omega|\cdot\,\Omega}$. The following are equivalent.
\begin{enumerate}
 \item[\rm(i)] \emph{Modular autonomy:} $\sigma^{\omega}_t(\mathcal N)=\mathcal N$ for all $t\in\mathbb R$.
 \item[\rm(ii)] \emph{Expected inclusion:} there is a normal conditional expectation $\mE:\mathcal M\to\mathcal N$ with $\omega\circ\mE=\omega$.
 \item[\rm(iii)] \emph{Trivial subsystem flow:} $C_{N\subset M}(t)=1_{\mathcal H_N}$ for all $t\in\mathbb R$.
\end{enumerate}
When these hold, $\mE$ is unique, its GNS projection is $e_N=VV^*$, and $\mathcal T\circ\mE_*=\operatorname{id}$ on $\mathcal N_*$ while $\mE_*\circ\mathcal T=\operatorname{id}$ on the code $\mathcal C_{\mE}\ni\omega$.
\end{theorem}

\begin{proof}
(i)$\Leftrightarrow$(ii) is Takesaki's theorem~\cite{Takesaki1972}; uniqueness and the identification of the GNS projection with $e_N$ are part of the same theorem, since $\mE(m)\Omega=e_Nm\Omega$ for an $\omega$-preserving expectation.

(ii)$\Rightarrow$(iii). For an $\omega$-preserving expectation the modular flows intertwine on $\mathcal N$,
$\sigma^{\omega}_t(n)=\sigma^{\omega;N}_t(n)$, and on the GNS space
$\Delta_M^{it}V=V\Delta_N^{it}$, \eqref{eq:takesaki-intertwining}. Hence
$C_{N\subset M}(t)=V^*V\Delta_N^{-it}\Delta_N^{it}=1$.

(iii)$\Rightarrow$(i). If $C(t)=1$ then $V^*\Delta_M^{-it}V=\Delta_N^{-it}$ is unitary on $\mathcal H_N$, so for $\psi\in\mathcal H_N$,
$\|e_N\Delta_M^{-it}V\psi\|=\|V\psi\|=\|\Delta_M^{-it}V\psi\|$,
which forces $\Delta_M^{-it}V\psi\in\operatorname{Ran}V$. Thus $\Delta_M^{-it}$ leaves $\mathcal H_N$ invariant and restricts there to $V\Delta_N^{-it}V^*$. For $n\in\mathcal N$,
\[
 \sigma^{\omega}_{-t}(n)\Omega=\Delta_M^{-it}n\Omega
 =V\Delta_N^{-it}n\Omega=\sigma^{\omega;N}_{-t}(n)\Omega .
\]
Since $\Omega$ is separating for $\mathcal M$, $\sigma^{\omega}_{-t}(n)=\sigma^{\omega;N}_{-t}(n)\in\mathcal N$ for all $t$, which is (i).

The final statements follow from $\mE_*$ being a right inverse of $\mathcal T$ on $\mathcal N_*$ and the identity $\varphi\circ\mE=\varphi$ on the code.
\end{proof}

Two remarks limit what Theorem~\ref{thm:correctable} asserts. First, it concerns the specific channel \eqref{eq:erasure-channel} and the specific family \eqref{eq:code}; failure of (ii) for the state $\omega$ does not exclude expectations preserving another faithful normal state, nor error correction for a different channel or family of states. Second, neither the Jones projection $e_N$ nor a product of modular conjugations is, by itself, a recovery channel: recovery is implemented by $\mE_*$, which exists only under the hypotheses of the theorem.



\subsection{Knill--Laflamme conditions and the forgetful complement}\label{subsec:QEC-KL}

Theorem~\ref{thm:correctable} identifies exact correctability with modular autonomy. This subsection records how the same statement reads in the standard language of quantum error correction: as a Knill--Laflamme condition, and as its equivalent formulation that the complementary channel be \emph{forgetful} on the code. The two are the same set of equations, and in the modular setting they are the Jones-projection form of Takesaki's condition.

\paragraph{The finite-dimensional dictionary.}
Let $\Phi(\rho)=\sum_kE_k\rho E_k^\dagger$ be a channel on $B(\mathcal H_S)$ with Stinespring isometry $W\xi=\sum_kE_k\xi\otimes\ket{k}_E$, so that $\Phi=\tr_E\circ\Ad W$. The complementary channel is $\Phi^c=\tr_S\circ\Ad W$, explicitly
\begin{equation}
 \Phi^c(\rho)=\sum_{k,l}\tr\!\bigl(E_l^\dagger E_k\,\rho\bigr)\,\ket{k}\!\bra{l}.
 \label{eq:complementary-channel}
\end{equation}
A code subspace $P\mathcal H_S$ is exactly correctable for $\Phi$ if and only if~\cite{KnillLaflamme1997}
\begin{equation}
 P\,E_l^\dagger E_k\,P=\lambda_{lk}\,P\qquad\hbox{for all }k,l,
 \label{eq:KL}
\end{equation}
for some matrix $\lambda$. For a state $\rho$ supported on $P$ the matrix elements of \eqref{eq:complementary-channel} are $\Phi^c(\rho)_{kl}=\tr(PE_l^\dagger E_kP\,\rho)$. Hence \eqref{eq:KL} holds if and only if
\begin{equation}
 \Phi^c(\rho)=\Lambda\quad\hbox{for every state $\rho$ supported on }P,\qquad\Lambda_{kl}=\lambda_{lk},
 \label{eq:forgetful}
\end{equation}
because the code states span all operators on $P\mathcal H_S$. Condition \eqref{eq:forgetful} says that the complementary channel is constant, or forgetful, on the code: the environment ends in the same state whatever logical state was encoded. Conditions \eqref{eq:KL} and \eqref{eq:forgetful} are therefore one statement read from two sides. From the system's side, no pair of errors can be told apart on the code except through the code-independent Gram matrix $\lambda$; from the environment's side, the environment, whose state records which error occurred, carries no information about the logical input. The content of the Knill--Laflamme theorem is the further implication that \eqref{eq:forgetful} suffices for recovery by a channel acting on the system alone. Its approximate version is the information--disturbance trade-off~\cite{SchumacherWestmoreland2002,KretschmannSchlingemannWerner2008}, and for a logical \emph{algebra} rather than a subspace the scalar $\lambda_{lk}$ is replaced by an element of the commutant of the logical algebra on the code~\cite{BenyKempfKribs2007}.

\paragraph{The modular dictionary.}
For the erasure channel \eqref{eq:erasure-channel} the pieces are as follows. The system is $\mathcal M$ and the logical algebra is $\mathcal N$. The errors are the elements of the relative commutant $\mathcal R=\mathcal N'\cap\mathcal M$: in the tensor-product case $\mathcal M=\mathcal N\overline\otimes\mathcal R$, erasure of the factor $\mathcal R$ has Kraus operators $\bra{k}$ on the erased factor, and the operators $E_l^\dagger E_k=\ket{k}\!\bra{l}$ of \eqref{eq:KL} are the matrix units of $\mathcal R$. The complementary channel of the erasure on a vector state $\psi$ is restriction to the commutant $\mathcal N'$, the Stinespring environment of the inclusion on the GNS space; it contains $\mathcal R$, the part of the environment inside $\mathcal M$, and $\mathcal M'$, the reference that purifies $\mathcal M$. The code subspace is $\mathcal H_N=\overline{\mathcal N\Omega}$, with the Jones projection $e_N=VV^*$ as code projection. In this language Theorem~\ref{thm:correctable} is a Knill--Laflamme condition with operator-valued coefficients.

\begin{proposition}[Knill--Laflamme form of Theorem~\ref{thm:correctable}]\label{prop:KL}
Under the hypotheses of Theorem~\ref{thm:correctable}, let $e_N$ be the projection onto $\mathcal H_N=\overline{\mathcal N\Omega}$. The following are equivalent.
\begin{enumerate}
 \item[\rm(i)] There is a normal conditional expectation $\mE:\mathcal M\to\mathcal N$ with $\omega\circ\mE=\omega$.
 \item[\rm(ii)] $e_N\,\mathcal M\,e_N\subseteq\mathcal N e_N$: the compression of every observable of $\mathcal M$ to the code subspace is a logical observable.
 \item[\rm(iii)] $\sigma^\omega_t(\mathcal N)=\mathcal N$ for all $t\in\mathbb R$.
\end{enumerate}
When they hold, $e_N\,x\,e_N=\mE(x)\,e_N$ for all $x\in\mathcal M$.
\end{proposition}

\begin{proof}
(i)$\Leftrightarrow$(iii) is Takesaki's theorem~\cite{Takesaki1972}.

(i)$\Rightarrow$(ii). For $n,n'\in\mathcal N$ and $x\in\mathcal M$,
\[
 \langle n\Omega|\,x\,n'\Omega\rangle=\omega(n^*xn')=\omega\bigl(\mE(n^*xn')\bigr)=\omega\bigl(n^*\mE(x)n'\bigr)=\langle n\Omega|\,\mE(x)\,n'\Omega\rangle,
\]
by invariance of $\omega$ and the bimodule property. Since $\mE(x)\in\mathcal N$ commutes with $e_N\in\mathcal N'$, this is $e_Nxe_N=\mE(x)e_N$.

(ii)$\Rightarrow$(i). Define $\mE(x)\in\mathcal N$ by $e_Nxe_N=\mE(x)e_N$. It is well defined: the central support of $e_N$ in $\mathcal N'$ is the projection onto $\overline{\mathcal N'\mathcal N\Omega}\supseteq\overline{\mathcal N'\Omega}=\mathcal H$, since $\Omega$ is separating for $\mathcal N$ and hence cyclic for $\mathcal N'$, so $n\mapsto ne_N$ is a normal $*$-isomorphism of $\mathcal N$ onto $\mathcal Ne_N$. Thus $\mE$ is the composition of the normal completely positive map $x\mapsto e_Nxe_N$ with the inverse of this isomorphism: it is linear, normal, completely positive and unital. It is $\mathcal N$-bimodular because $e_N\in\mathcal N'$ gives $e_N\,nxn'\,e_N=n\,(e_Nxe_N)\,n'$, and it preserves $\omega$ because $\omega(\mE(x))=\langle\Omega|\mE(x)e_N\Omega\rangle=\langle\Omega|e_Nxe_N\Omega\rangle=\omega(x)$.
\end{proof}

The identity $e_Nxe_N=\mE(x)e_N$ is \eqref{eq:KL} with an operator-valued coefficient. For the errors themselves, $c\in\mathcal R$, the coefficient $\mE(c)=\chi(c)$ is a scalar by Lemma~\ref{lem:expectation-slice}, and the condition is the Knill--Laflamme condition in its original form. For a general $x\in\mathcal M$ the coefficient is the logical operator $\mE(x)$, which expresses the Schr\"odinger-picture requirement of Section~\ref{subsec:QEC-erasure} that the whole state on $\mathcal M$, and not only the expectation values of the logical algebra, be reconstructed from the restriction to $\mathcal N$. The Heisenberg-picture requirement that the expectation values of logical observables survive is automatic for the restriction to $\mathcal N$; the content of exact correctability here lies in the correlations between $\mathcal N$ and the discarded observables, and this is what the operator-valued coefficient records.

\begin{corollary}[Conormal inclusions: Knill--Laflamme is a forgetful complement]\label{cor:KL-conormal}
Assume in addition that the inclusion is conormal, $\mathcal N\vee\mathcal R=\mathcal M$. Then {\rm(i)--(iii)} are further equivalent to each of the following, for a faithful normal state $\chi$ on $\mathcal R$.
\begin{enumerate}
 \item[\rm(iv)] \emph{Knill--Laflamme for the gauge algebra:} $e_N\,c\,e_N=\chi(c)\,e_N$ for all $c\in\mathcal R$.
 \item[\rm(v)] \emph{Forgetful complement:} $\langle\psi|c\,\psi\rangle=\chi(c)$ for every unit vector $\psi\in\mathcal H_N$ and every $c\in\mathcal R$; the complementary channel $\psi\mapsto\psi|_{\mathcal R}$ is constant on the code subspace.
 \item[\rm(vi)] \emph{No logical--gauge correlations in the reference:} $\omega(nc)=\omega(n)\chi(c)$ for all $n\in\mathcal N$, $c\in\mathcal R$.
\end{enumerate}
\end{corollary}

\begin{proof}
(iv)$\Leftrightarrow$(v) is polarization. (v)$\Leftrightarrow$(vi): for $n,n'\in\mathcal N$ and $c\in\mathcal R$, $\langle n\Omega|c\,n'\Omega\rangle=\omega(n^*cn')=\omega(c\,n^*n')$ because $c$ commutes with $\mathcal N$, so (v) reads $\omega(cn'')=\chi(c)\omega(n'')$ for all $n''\in\mathcal N$, with $\chi=\omega|_{\mathcal R}$, faithful since $\omega$ is. (i)$\Rightarrow$(vi) is Lemma~\ref{lem:expectation-slice}, applied to the expectation of Theorem~\ref{thm:correctable}, which is faithful because $\omega$ is. (vi)$\Rightarrow$(i): $\mathcal R$ is a factor, since an element of its centre commutes with $\mathcal N$ and with $\mathcal R$, hence with $\mathcal M$. Thus $\omega$ is a faithful normal product state on the commuting factors $\mathcal N$ and $\mathcal R$, which generate $\mathcal M$; as in the proof of Proposition~\ref{prop:conormal-factorization}, multiplication extends to a spatial isomorphism $\mathcal N\overline\otimes\mathcal R\cong\mathcal M$ carrying $\omega|_N\otimes\chi$ to $\omega$~\cite{DoplicherLongo1984}, and $\mE=\operatorname{id}\otimes\chi$ is a normal conditional expectation preserving $\omega$.
\end{proof}

Read together with Theorem~\ref{thm:correctable}, the corollary says that for conormal inclusions Takesaki's condition \emph{is} the Knill--Laflamme condition for erasure of the gauge algebra, and that the forgetful complement is its environment-side reading. The intuition is the one behind the Takesaki condition of Section~\ref{sec:intro}: the modular Hamiltonian $K_M=-\log\Delta_M$ generates the dynamics for which $\omega$ is an equilibrium state; correlations between $\mathcal N$ and $\mathcal R$ in $\omega$ are an interaction term in $K_M$ coupling logical to gauge degrees of freedom; and such a term is precisely what transports logical observables out of $\mathcal N$. In finite dimensions, where $\mathcal N=\bigoplus_aM_{n_a}\otimes1_{m_a}$, all of (i)--(vi) reduce to $\rho=\bigoplus_ap_a\,\rho_a\otimes\tau_a$, the standard form of a code state of an operator-algebra code (Appendix~\ref{app:finite-coarse}).

\paragraph{Non-conormal inclusions.}
Without conormality the gauge algebra $\mathcal R$ does not account for the discarded observables, and (iv)--(v) are not sufficient for (i). For an irreducible subfactor, $\mathcal R=\mathbb C1$ and (iv), (v) hold trivially, while (i) fails for every proper inclusion with $\Omega$ cyclic for $\mathcal N$ (Corollary~\ref{cor:common-cyclic-noncorrectable}). The environment of the erasure is then the full commutant $\mathcal N'$, which contains $\mathcal R\vee\mathcal M'=(\mathcal R'\cap\mathcal M)'$ and equals it only when $\mathcal N=\mathcal R'\cap\mathcal M$; and ``constant on the code'' cannot hold on $\mathcal N'$, because the reference $\mathcal M'$ varies over the code by design. Its correct replacement is: for every $\psi\in\mathcal H_N$ the restriction $\mathcal N'\to\mathcal M'$ is sufficient for the pair $(\psi,\omega)$, i.e.\ the environment holds no information distinguishing a code vector from the reference beyond what the reference system $\mathcal M'$ already holds. This follows from (i). The support of $\omega$ on $\mathcal N'$ is $e_N$, so for $\psi\in\mathcal H_N$ the relative entropies on $\mathcal N'$ reduce to the compression $e_N\mathcal N'e_N$, which is the commutant of $\mathcal N$ on $\mathcal H_N$; the decomposition of Proposition~\ref{prop:sufficiency} below therefore holds for the compressed generator $V^*K_MV-K_N$ of Appendix~\ref{app:compressed-comparison} on $\mathcal H_N$, and (i) gives $C_{N\subset M}(t)\equiv1$, hence $V^*K_MV=K_N$ on $\mathcal H_N$ and the vanishing of both nonnegative gaps.


\section{Subsystem flow and ergodic subsystems}





% \subsection{Subsystem flow}


Let $\Omega$ be cyclic and separating for $\mathcal M$, and set
$\mathcal H_N=\overline{\mathcal N\Omega}$ for a unital inclusion $\mathcal N\subset\mathcal M$.
The restricted state is faithful, and its modular operator $\Delta_N$ acts on
$\mathcal H_N$. With $V:\mathcal H_N\hookrightarrow\mathcal H$ and $e_N=VV^*$, 
define the {\it compressed subsystem flow}
\begin{equation}
 C_{N\subset M}(t)=V^*\Delta_M^{-it}V\Delta_N^{it},
 \qquad t\in\mathbb R.
\end{equation}
This is a strongly continuous contraction on $\mathcal H_N$ fixing $\Omega$. As we saw in the Taksaki condition, if an $\Omega$-preserving expectation exists, the subalgebra $\mN$ is invariant under its modular flow, and the compressed subsystem flow is trivial
$C_{N\subset M}(t)=1_{\mathcal H_N}$.

To display the compression algebraically, extend the intrinsic unitary to
$U_N(t)=V\Delta_N^{it}V^*+(1-e_N)$ on $\mathcal H$ and set
$u_t=\Delta_M^{-it}U_N(t)$, $\widehat C(t)=e_Nu_te_N=VC(t)V^*$.
Writing $\widetilde\sigma^{\omega;N}_{-t}=\operatorname{Ad}U_N(-t)$ gives
\begin{align}
 \widehat C(t+s)&=\widehat C(t)\widetilde\sigma^{\omega;N}_{-t}(\widehat C(s))
                  +\mathcal L(t,s),\nn\\
 \mathcal L(t,s)&=e_Nu_t(1-e_N)\widetilde\sigma^{\omega;N}_{-t}(u_s)e_N.
\end{align}
The second term records passage through the orthogonal complement. This algebraic
identity does not itself specify an open-system master equation.
On vectors $\psi\in\operatorname{Dom}K_N$ with $V\psi\in\operatorname{Dom}K_M$,
\begin{equation}
 -i\left.\partial_t C_{N\subset M}(t)\psi\right|_{t=0}
 =(V^*K_MV-K_N)\psi,\qquad K_C=-\log\Delta_C.
\end{equation}
This compares modular generators on the stated domain, rather than asserting
a strong derivative on the full Hilbert space or a dissipative generator.
The contraction is a Hilbert-space operator. Conjugation
$a\mapsto C(t)^*aC(t)$ is CP on $B(\mathcal H_N)$ but is unital only if
$C(t)^*C(t)=1$; preservation of $\mathcal N$ would require a further condition.
An interpretation as a quantum channel therefore needs additional structure.

\subsection{Subsystem flow as a sufficiency witness}\label{subsec:QEC-sufficiency}

The subsystem flow does more than detect whether the inclusion is correctable: its generator quantifies, state by state, how far the erasure is from being correctable. The relevant notion is Petz's \emph{sufficiency}~\cite{Petz1986,Petz1988,OhyaPetz1993}. Let $\alpha:\mathcal N\to\mathcal M$ be a unital normal completely positive map (here the inclusion) and $\varphi,\omega$ faithful normal states on $\mathcal M$. The restriction $\alpha$ is \emph{sufficient} for the pair $(\varphi,\omega)$ if there is a unital normal completely positive map $\beta:\mathcal M\to\mathcal N$ with $\varphi\circ\alpha\circ\beta=\varphi$ and $\omega\circ\alpha\circ\beta=\omega$; equivalently, both states are recovered from their restrictions to $\mathcal N$ by one and the same channel.

\begin{theorem}[Petz]\label{thm:Petz}
For faithful normal states $\varphi,\omega$ on $\mathcal M$ with $S(\varphi\|\omega)<\infty$, the restriction to $\mathcal N$ is sufficient for $(\varphi,\omega)$ if and only if
\[
 S_M(\varphi\|\omega)=S_N(\varphi|_N\|\omega|_N) .
\]
In that case a recovery channel is given by the Petz map of the inclusion with respect to $\omega$. Equivalently, the Connes cocycle $[D\varphi:D\omega]_t=\Delta_{\varphi|\omega}^{it}\Delta_{\omega}^{-it}$ belongs to $\mathcal N$ for every $t\in\mathbb R$~\cite{Petz1988,JencovaPetz2006}.
\end{theorem}

The cocycle condition is the form of Petz's theorem closest to the subsystem flow. The Connes cocycle compares the modular groups of two states on one algebra, $\sigma^{\varphi}_t=\operatorname{Ad}[D\varphi:D\omega]_t\circ\sigma^{\omega}_t$; the subsystem flow $D_{N\subset M}(t)=\Delta_M^{-it}\Delta_N^{it}$ compares the modular groups of one state on two algebras. Both are cocycles rather than one-parameter groups, and the generator of the second produces exactly the entropy differences that appear in Theorem~\ref{thm:Petz}. The link is the following identity, which expresses the modular Hamiltonian of a state as the imbalance between the information held by an algebra and by its commutant.

\begin{lemma}[Modular energy as an entropy imbalance]\label{lem:modular-energy}
Let $\mathcal C$ be a von Neumann algebra on $\mathcal H$, let $\Omega$ and $\Psi$ be unit vectors that are cyclic and separating for $\mathcal C$, with induced states $\omega$ and $\varphi$, and suppose that $S_C(\varphi\|\omega)$ and $S_{C'}(\varphi\|\omega)$ are finite. Then $t\mapsto\langle\Psi|\Delta_{\omega;C}^{it}|\Psi\rangle$ is differentiable at $t=0$ and
\begin{equation}
 -i\,\partial_t\langle\Psi|\Delta_{\omega;C}^{it}|\Psi\rangle\big|_{t=0}
 =S_{C'}(\varphi\|\omega)-S_C(\varphi\|\omega).
 \label{eq:modular-energy}
\end{equation}
If moreover $\Psi$ lies in the form domain of $\log\Delta_{\omega;C}$, the left side is the quadratic form $\langle\Psi|\log\Delta_{\omega;C}|\Psi\rangle$.
\end{lemma}

\begin{proof}
Write $\omega',\varphi'$ for the states induced on $\mathcal C'$. For cyclic and separating vectors $\Phi,\Psi$ the relative modular operators $\Delta_{\Phi|\Psi;C}$ of \eqref{eq:relative-modular} satisfy~\cite{Araki1976,Araki1977}
\begin{equation}
 \Delta_{\Phi|\Psi;C}^{it}=[D\Phi:D\Psi]^{C}_t\,\Delta_{\Psi;C}^{it},
 \qquad
 \Delta_{\Phi|\Psi;C'}=\Delta_{\Psi|\Phi;C}^{-1},
 \label{eq:relative-modular-cocycles}
\end{equation}
where $[D\Phi:D\Psi]^{C}_t\in\mathcal C$ is the Connes cocycle of the two states on $\mathcal C$. The first identity changes the state in the first slot; applied on $\mathcal C'$ and transported by the second identity, it changes the state in the second slot: $\Delta_{\Phi|\Psi;C}^{it}=\Delta_{\Phi;C}^{it}\,[D\Phi':D\Psi']^{C'}_t$. Applying both to $\Delta_{\omega;C}=\Delta_{\omega|\omega;C}$ gives
\begin{equation}
 \Delta_{\omega;C}^{it}=[D\omega:D\varphi]^{C}_t\;\Delta_{\varphi;C}^{it}\;[D\varphi':D\omega']^{C'}_t .
 \label{eq:three-factor}
\end{equation}
Set $a_t=[D\omega:D\varphi]^{C}_t\in\mathcal C$ and $b_t=\Delta_{\varphi;C}^{it}[D\varphi':D\omega']^{C'}_t\Delta_{\varphi;C}^{-it}\in\mathcal C'$, both unitary with $a_0=b_0=1$. Since $\Delta_{\varphi;C}^{it}\Psi=\Psi$, \eqref{eq:three-factor} gives
\begin{align}
 \langle\Psi|\Delta_{\omega;C}^{it}|\Psi\rangle&=\langle\Psi|a_tb_t|\Psi\rangle,\nn\\
 \langle\Psi|a_t|\Psi\rangle&=\varphi\bigl([D\omega:D\varphi]_t\bigr),
 \qquad
 \langle\Psi|b_t|\Psi\rangle=\varphi'\bigl([D\varphi':D\omega']_t\bigr).
\end{align}
By the first identity in \eqref{eq:relative-modular-cocycles}, $\varphi([D\omega:D\varphi]_t)=\langle\Psi|\Delta_{\omega|\varphi;C}^{it}|\Psi\rangle=\int e^{it\lambda}\,d\mu(\lambda)$, with $\mu$ the spectral measure of $\log\Delta_{\omega|\varphi;C}$ in $\Psi$. Here $\int e^{\lambda}d\mu=\|\Delta_{\omega|\varphi;C}^{1/2}\Psi\|^2=\|\Omega\|^2=1$ bounds the positive part of $\lambda$, and $S_C(\varphi\|\omega)=-\int\lambda\,d\mu<\infty$ then bounds the negative part, so $\int|\lambda|\,d\mu<\infty$ and $\varphi([D\omega:D\varphi]_t)$ is differentiable at $0$ with derivative $-iS_C(\varphi\|\omega)$~\cite{OhyaPetz1993}. The same argument on $\mathcal C'$ gives $\partial_t\varphi'([D\omega':D\varphi']_t)|_{t=0}=-iS_{C'}(\varphi\|\omega)$, and since $[D\varphi':D\omega']_t=[D\omega':D\varphi']_t^*$, the derivative of $\langle\Psi|b_t|\Psi\rangle$ at $0$ is $+iS_{C'}(\varphi\|\omega)$. Finally
\begin{equation}
 \langle\Psi|a_tb_t|\Psi\rangle-1
 =\bigl(\langle\Psi|a_t|\Psi\rangle-1\bigr)+\bigl(\langle\Psi|b_t|\Psi\rangle-1\bigr)
 +\langle(a_t^*-1)\Psi|(b_t-1)\Psi\rangle,
\end{equation}
and $\|(a_t^*-1)\Psi\|^2=\|(a_t-1)\Psi\|^2=2\operatorname{Re}\bigl(1-\langle\Psi|a_t|\Psi\rangle\bigr)=o(t)$, because the derivative of $\langle\Psi|a_t|\Psi\rangle$ at $0$ is purely imaginary; likewise for $b_t$. The cross term is therefore $o(t)$, and
$-i\,\partial_t\langle\Psi|\Delta_{\omega;C}^{it}|\Psi\rangle|_{t=0}=-S_C(\varphi\|\omega)+S_{C'}(\varphi\|\omega)$. The last statement is the spectral theorem for $\log\Delta_{\omega;C}$ on its form domain. In finite dimensions $\log\Delta_{\omega;C}=\log\rho^\omega_{C}\otimes1-1\otimes\log\rho^\omega_{C'}$ on $\mathcal H_C\otimes\mathcal H_{C'}$, and \eqref{eq:modular-energy} reduces to the equality of the von Neumann entropies of the two marginals of a pure state.
\end{proof}

\begin{proposition}[Sufficiency decomposition of the subsystem-flow generator]\label{prop:sufficiency}
Let $\Omega$ be a common cyclic and separating vector for $\mathcal N\subset\mathcal M$, and let $\Psi$ be a unit vector, cyclic and separating for $\mathcal N$ and for $\mathcal M$, for which the four relative entropies $S_M(\varphi\|\omega)$, $S_N(\varphi\|\omega)$, $S_{N'}(\varphi\|\omega)$, $S_{M'}(\varphi\|\omega)$ are finite and $\Psi\in\operatorname{Dom}K_M\cap\operatorname{Dom}K_N$, where $\varphi=\braket{\Psi|\cdot\,\Psi}$. Then
\begin{equation}
 -i\,\partial_t\langle\Psi|D_{N\subset M}(t)|\Psi\rangle\big|_{t=0}
 =\langle\Psi|K_M-K_N|\Psi\rangle
 =\underbrace{\bigl[S_M(\varphi\|\omega)-S_N(\varphi\|\omega)\bigr]}_{\text{gap for }\mathcal N\subset\mathcal M}
 +\underbrace{\bigl[S_{N'}(\varphi\|\omega)-S_{M'}(\varphi\|\omega)\bigr]}_{\text{gap for }\mathcal M'\subset\mathcal N'} ,
 \label{eq:sufficiency-decomposition}
\end{equation}
with both gaps nonnegative. The generator vanishes on $\Psi$ if and only if the restriction $\mathcal M\to\mathcal N$ is sufficient for $(\varphi,\omega)$ and the restriction $\mathcal N'\to\mathcal M'$ is sufficient for $(\varphi,\omega)$.
\end{proposition}

\begin{proof}
On $\operatorname{Dom}K_M\cap\operatorname{Dom}K_N$ the generator of $D_{N\subset M}(t)=\Delta_M^{-it}\Delta_N^{it}$ at $t=0$ is $K_M-K_N$, which is the first equality. Lemma~\ref{lem:modular-energy}, applied to $\mathcal C=\mathcal M$ and to $\mathcal C=\mathcal N$, gives $\langle\Psi|K_M|\Psi\rangle=S_M(\varphi\|\omega)-S_{M'}(\varphi\|\omega)$ and $\langle\Psi|K_N|\Psi\rangle=S_N(\varphi\|\omega)-S_{N'}(\varphi\|\omega)$; subtracting and regrouping gives the second equality. Nonnegativity of each bracket is monotonicity of relative entropy under $\mathcal N\subset\mathcal M$ and $\mathcal M'\subset\mathcal N'$. Vanishing of the sum of two nonnegative terms is vanishing of each, and by Theorem~\ref{thm:Petz} each vanishes exactly when the corresponding restriction is sufficient for $(\varphi,\omega)$.
\end{proof}

\begin{corollary}\label{cor:positive-generator}
Under the hypotheses of Proposition~\ref{prop:sufficiency}, if $\mathcal N\subset\mathcal M$ is proper then the quadratic form $\Psi\mapsto\langle\Psi|K_M-K_N|\Psi\rangle$ is nonnegative and does not vanish identically on any dense subspace of $\operatorname{Dom}K_M\cap\operatorname{Dom}K_N$ that is a core for $K_M$: there exist states $\varphi$ for which the erasure \eqref{eq:erasure-channel} is not sufficient, and the subsystem flow registers them.
\end{corollary}
\begin{proof}
Nonnegativity is \eqref{eq:sufficiency-decomposition}. If the form vanished on a dense subspace $D\subseteq\operatorname{Dom}K_M\cap\operatorname{Dom}K_N$, polarization would give $K_M\Psi=K_N\Psi$ for $\Psi\in D$; if $D$ is a core for $K_M$, the closure of $K_M|_D$ is $K_M$ and is contained in the self-adjoint operator $K_N$, so $K_M=K_N$, hence $\Delta_{\omega;M}=\Delta_{\omega;N}$ and $\sigma^{\omega}_t(\mathcal N)=\mathcal N$ for all $t$, contradicting Corollary~\ref{cor:common-cyclic-noncorrectable}.
\end{proof}

Proposition~\ref{prop:sufficiency} is the precise sense in which the subsystem flow is a generating function for error correction. Its generator is a sum of two sufficiency gaps, one for the coarse-graining and one for the dual coarse-graining of the commutants; the gaps are independent in general, and the vector $\Psi$ selects the state being tested; on the logical orbit $\{u\Omega:u\in\mathcal N\ \hbox{unitary}\}$ the second gap vanishes and the form reduces to the first (Section~\ref{subsec:QEC-KL}). In the classical setting of Section~\ref{sec:classical} only the first gap is present, the dual inclusion being trivial, and sufficiency reduces to the product-state condition \eqref{eq:classical-product-state-closure}.

The analytic continuation of $D_{N\subset M}$ to the strip $0\le\operatorname{Im}z\le1/2$ (Section~\ref{subsec:quantum-coarse}) refines this statement from the generator to the full flow: the contraction $D_{N\subset M}(i\alpha)=\Delta_M^{\alpha}\Delta_N^{-\alpha}$, $0\le\alpha\le1/2$, is the operator whose norm bound \eqref{eq:characteristic-bounded} expresses monotonicity of the modular operators, and its matrix elements interpolate between the sufficiency gaps at $\alpha=0$ and the modular conjugations at $\alpha=1/2$. The relative version of this operator, built from the relative modular operators of a pair $(\varphi,\omega)$, is the subject of the equality conditions in Theorem~\ref{thm:Petz}; the subsystem flow is its specialization to the comparison of one state on two algebras.


\paragraph{Failure of Knill--Laflamme is measured by the subsystem flow.}
When (i)--(iii) fail, the two gaps of Proposition~\ref{prop:sufficiency} below are the two sides of \eqref{eq:KL}--\eqref{eq:forgetful} made quantitative: $S_M(\psi\|\omega)-S_N(\psi\|\omega)$ measures the failure of recovery from the retained algebra, and $S_{N'}(\psi\|\omega)-S_{M'}(\psi\|\omega)$ measures what the environment has learned. They are independent as functionals of $\psi$, and no inequality relates them pair by pair. For $\psi=u\Omega$ with $u$ a unitary of $\mathcal N$ the second gap vanishes identically, since $\psi$ agrees with $\omega$ on $\mathcal N'$, while the Connes cocycle of $(\psi,\omega)$ on $\mathcal M$ is $u\,\sigma^\omega_t(u)^*$, so by Theorem~\ref{thm:Petz} the first gap at $u\Omega$ vanishes exactly when $\sigma^\omega_t(u)\in\mathcal N$ for all $t$: Petz sufficiency, quantified over the logical unitaries, is Takesaki's condition. On the logical orbit $\{u\Omega\}$ the quadratic form of Corollary~\ref{cor:positive-generator} therefore reduces to a single gap, and it is this gap that detects non-correctability. This is no conflict with the information--disturbance trade-off of~\cite{KretschmannSchlingemannWerner2008}, which quantifies over a code subspace rather than over a pair of states: the span of the logical orbit is all of $\mathcal H_N$, on which the environment does learn. The trade-off survives in the modular setting only at the level of the algebra, as the equivalence of Proposition~\ref{prop:KL}, and the subsystem flow is its generating function state by state.

\subsection{History, memory, generating and forgetful flows}\label{subsec:QEC-history-memory}

More generally, for a state that has correlations between $R$ and $B$ we can define a {\it history algebra} $\mathcal N_{\rm hist}:=\vee_{t\in\mathbb{R}}\mathcal N_t$ and a {\it memory algebra} $\mathcal N_{\rm mem}:=\cap_{t\in\mathbb{R}}\mathcal N_t$. By construction, both algebras evolve autonomously in $\mathcal M$. If the relative commutant $\mathcal N_{\rm hist}^c=\mathcal N'_{\rm hist}\cap \mathcal M$ is non-trivial, then the degrees of freedom in it are untouched by the dynamics, and we will see in the next section that in the language of QEC, the erasure of $\mathcal N_{\rm hist}^c$ degrees of freedom can be unitarily corrected on $\mathcal N_{\rm hist}$. The information-disturbance relation can be related to the fact that $\mathcal N_{\rm hist}^c=\mathcal R_{\rm mem}$ and $\mathcal R_{\rm hist}\subset \mathcal N_{\rm mem}^c$. 



The ergodic content of the subsystem flow enters through the orbit of $\mathcal N$ under the modular flow. Let $\alpha_t=\sigma^{\omega}_t$, which is implemented on $\mathcal H$ by $\Delta_M^{it}$ and preserves $\mathcal M$ and $\mathcal M'$.

\begin{definition}[History and memory algebras]\label{def:history-memory}
For a subalgebra $\mathcal N\subset\mathcal M$,
\begin{equation}
 \mathcal N_{\rm hist}:=\bigvee_{t\in\mathbb R}\alpha_t(\mathcal N),\qquad
 \mathcal N_{\rm mem}:=\bigcap_{t\in\mathbb R}\alpha_t(\mathcal N).
\end{equation}
The inclusion is \emph{generating} if $\mathcal N_{\rm hist}=\mathcal M$ (Definition~\ref{def:generating}) and \emph{forgetful} if $\mathcal N_{\rm mem}=\mathbb C1$.
\end{definition}

Both algebras are invariant under $\alpha_t$ by construction. By Takesaki's theorem each therefore admits an $\omega$-preserving normal expectation: $\mathcal N_{\rm mem}$ is the largest subalgebra of $\mathcal N$ that is exactly correctable in the sense of Theorem~\ref{thm:correctable}, and $\mathcal N_{\rm hist}$ is the smallest exactly correctable algebra containing $\mathcal N$. In the language of noiseless subsystems, $\mathcal N_{\rm mem}$ is the noiseless subalgebra of $\mathcal N$ for the composite channel ``evolve by the modular flow, then restrict to $\mathcal N$'': the observables of $\mathcal N$ that remain accessible from $\mathcal N$ at every modular time. These statements are about algebras and expectations; they do not assert that $\mE\circ\alpha_t|_N$ is a Markovian semigroup, which would require additional structure (Appendix~\ref{app:compressed-comparison}).

\begin{lemma}[History--memory duality]\label{lem:history-memory}
Let $\alpha_t=\operatorname{Ad}U_t$ be a one-parameter group of automorphisms of $B(\mathcal H)$ with $U_t\mathcal MU_t^*=\mathcal M$, hence also $U_t\mathcal M'U_t^*=\mathcal M'$, and let history and memory be defined as in Definition~\ref{def:history-memory} for any von Neumann algebra on $\mathcal H$. For every subalgebra $\mathcal N\subset\mathcal M$:
\begin{enumerate}
 \item[\rm(a)] On $B(\mathcal H)$ the duality is exact,
 \begin{equation}
  (\mathcal N_{\rm hist})'=(\mathcal N')_{\rm mem},\qquad (\mathcal N_{\rm mem})'=(\mathcal N')_{\rm hist}.
  \label{eq:history-memory-BH}
 \end{equation}
 \item[\rm(b)] Inside $\mathcal M$, with $\mathcal R=\mathcal N'\cap\mathcal M$,
 \begin{equation}
  (\mathcal N_{\rm hist})'\cap\mathcal M=\mathcal R_{\rm mem},\qquad
  \mathcal R_{\rm hist}\subseteq(\mathcal N_{\rm mem})'\cap\mathcal M ,
  \label{eq:history-memory}
 \end{equation}
 with equality in the second relation when $\mathcal N$ and $\mathcal R_{\rm hist}$ each coincide with their double relative commutant in $\mathcal M$; this holds in particular when $\mathcal M=\mathcal N\overline\otimes\mathcal R$ and $\mathcal R$ is autonomous.
\end{enumerate}
\end{lemma}
\begin{proof}
(a) For von Neumann algebras the commutant of a join is the intersection of the commutants and the commutant of an intersection is the join of the commutants, and $\operatorname{Ad}U_t$ commutes with taking commutants; hence $(\bigvee_t\alpha_t(\mathcal N))'=\bigcap_t\alpha_t(\mathcal N')$ and $(\bigcap_t\alpha_t(\mathcal N))'=\bigvee_t\alpha_t(\mathcal N')$.
(b) Intersecting the first identity of (a) with $\mathcal M=\alpha_t(\mathcal M)$ gives $\bigcap_t\alpha_t(\mathcal N')\cap\mathcal M=\bigcap_t\alpha_t(\mathcal N'\cap\mathcal M)=\mathcal R_{\rm mem}$. For the inclusion, $\alpha_t(\mathcal N'\cap\mathcal M)\subseteq\alpha_t(\mathcal N')\cap\mathcal M\subseteq(\bigcap_s\alpha_s(\mathcal N))'\cap\mathcal M$ for every $t$; take the join over $t$. For the equality statement write $\mathcal P^c=\mathcal P'\cap\mathcal M$. The first relation of (b) applied to $\mathcal R$ in place of $\mathcal N$ reads $(\mathcal R_{\rm hist})^c=(\mathcal R^c)_{\rm mem}$, and $\mathcal R^c=\mathcal N$ when $\mathcal N=\mathcal N^{cc}$; taking relative commutants once more, $(\mathcal R_{\rm hist})^{cc}=(\mathcal N_{\rm mem})^c$, which equals $\mathcal R_{\rm hist}$ when $\mathcal R_{\rm hist}=(\mathcal R_{\rm hist})^{cc}$. In the tensor-product case $\mathcal N=\mathcal N^{cc}$, and if $\mathcal R$ is autonomous then $\mathcal R_{\rm hist}=\mathcal R=\mathcal R^{cc}$.
\end{proof}

\begin{corollary}[Generation forces a forgetful environment]\label{cor:generating-forgetful}
If $\mathcal M$ is a factor and $\mathcal N\subset\mathcal M$ is generating, then $\mathcal R_{\rm mem}=\mathbb C1$: no nontrivial observable of the relative commutant is permanently retained by the modular flow. In particular this holds whenever $\Omega$ is a common cyclic and separating vector for $\mathcal N$ and $\mathcal M$ {\rm(Lemma~\ref{lem:common-cyclicity-generating})}.
\end{corollary}
\begin{proof}
$\mathcal R_{\rm mem}=(\mathcal N_{\rm hist})'\cap\mathcal M=\mathcal M'\cap\mathcal M=\mathbb C1$.
\end{proof}

Corollary~\ref{cor:generating-forgetful} is the modular form of the decoupling principle of error correction: if the logical algebra $\mathcal N$ dynamically generates the full algebra, then its environment $\mathcal R$ retains no permanent memory. Since the Stinespring environment of the erasure is the full commutant $\mathcal N'$ (Section~\ref{subsec:QEC-KL}), part (a) of the lemma completes the picture: $(\mathcal N')_{\rm mem}=(\mathcal N_{\rm hist})'\supseteq\mathcal M'$, with equality exactly when $\mathcal N$ is generating. The reference $\mathcal M'$ is the frozen part of the environment, fixed setwise by the modular flow of $\omega$ and therefore always in its memory; $\mathcal R$ is the part the flow can move, and generation says that it forgets everything. The operators in $(\mathcal N_{\rm mem})'\cap\mathcal M$ that are not in $\mathcal R_{\rm hist}$ commute with all of the permanent memory of $\mathcal N$ but cannot be generated from within $\mathcal R$; they are the modular analogue of boundary degrees of freedom. Note that the corollary does not say $\mathcal N$ itself is forgetful. Common cyclicity gives generation and excludes an $\omega$-preserving expectation onto $\mathcal N$; it neither supplies a decay rate for connected correlations nor a trivial one-sided tail for $\mathcal N$.

To make the discussion less abstract, below, we give a few examples of the concepts introduced in this subsection:

\paragraph{Finite-dimensional example (generating inclusion)}
Everything in this section already happens for two qubits. Take $\mathcal M=M_2\otimes M_2$, $\mathcal N=M_2\otimes1$, $\mathcal R=1\otimes M_2$, and the reference density matrix $\rho=(1-\epsilon)\ket{\Phi^+}\bra{\Phi^+}+\epsilon\,\mathbbm 1/4$ with $\ket{\Phi^+}$ maximally entangled and $0<\epsilon<1$, purified by a vector $\Omega$ that is cyclic and separating for $\mathcal M$. Its spectrum $\{1-3\epsilon/4,\epsilon/4,\epsilon/4,\epsilon/4\}$ is not a product of two two-point spectra, so $\rho\ne U(\rho_B\otimes\rho_R)U^\dagger$ for every unitary $U$: no encoding of one qubit into the two has $\rho$ as a code state, and Theorem~\ref{thm:correctable} fails for $\mathcal N$ and for every $U\mathcal NU^\dagger$. The modular Hamiltonian $K_M=-\log\rho$ differs from a multiple of the identity by a multiple of $\ket{\Phi^+}\bra{\Phi^+}=\tfrac14(\mathbbm 1+X\otimes X-Y\otimes Y+Z\otimes Z)$, a pure interaction term, and $i[K_M,Z\otimes1]\propto X\otimes Y+Y\otimes X$ leaves $\mathcal N$: this is the leakage the subsystem flow measures. Since $\rho_B=\rho_R=\mathbbm 1/2$, a logical unitary $u\otimes1$ changes neither marginal, so for $\Psi=(u\otimes1)\Omega$ the whole sufficiency gap of Proposition~\ref{prop:sufficiency} sits in the correlations between $\mathcal N$ and $\mathcal R$. Finally the inclusion is generating: if $1\otimes r\in\mathcal R_{\rm mem}$ then $i[K_M,1\otimes r]\in\mathcal R$, and averaging this commutator over the unitaries of $\mathcal N$, which fix $\mathcal R$ pointwise and average $\ket{\Phi^+}\bra{\Phi^+}$ to $\mathbbm 1/4$, shows that $1\otimes r$ commutes with $\ket{\Phi^+}\bra{\Phi^+}$, hence $r\propto1$. Thus $\mathcal R_{\rm mem}=\mathbb C1$, so by Lemma~\ref{lem:history-memory} the algebra $\mathcal N_{\rm hist}$, which is of the form $M_2\otimes\mathcal P$, has trivial relative commutant and equals $\mathcal M$. 
% The environment is forgetful and the inclusion generating although $\Omega$ is not cyclic for $\mathcal N$ and nothing is of type $III$; what finite dimensions cannot supply is the half-sided nesting of Section~\ref{subsec:geometry-modular-ergodic}.


\paragraph{Tensor amplification (an expected, nongenerating inclusion):}
Let $\mathcal N$ be a type $III_1$ factor and set $\mathcal M=\mathcal N\overline\otimes M_n(\mathbb C)$,
with $n>1$. Then $\mathcal R=1\otimes M_n(\mathbb C)$ and the inclusion is conormal.
The slice expectations \eqref{condexp} exist, and the minimal Jones--Kosaki index
is $n^2$. Finite index by itself neither specifies a matrix relative commutant
nor implies conormality.
For a faithful product state $\omega_M=\omega_N\otimes\omega_R$, the outer modular
flow preserves $\mathcal N$, so its modular orbit generates only $\mathcal N$.
The site spectral projections of $\omega_R$ lie in the centralizer of $\omega_M$.
In the product GNS representation, the compressed subsystem flow acts only on
the hidden factor. The reference vector is not cyclic for $\mathcal N$ on the full Hilbert
space, consistently with Lemma~\ref{lem:common-cyclicity-generating}.





\subsection{Common cyclic and separating vectors}

When $\mN\subset \mM$ have a common cyclic and separating vector, in the GNS Hilbert space of this vector, we have $e_N=1$. We can set $V=1$, and we denote the subsystem flow by $D_{N\subset M}(t)$, which is unitary for all real $t$.  In the algebraic quantum field theory literature, it is Borchers' \emph{characteristic function} of the inclusion~\cite{revolutionizing}; we use that name only when quoting results from~\cite{revolutionizing}.

% \begin{definition}[Characteristic flow]\label{def:subsystem-flow}
%     Given a pair of von Neumann algebras $\mathcal N\subset\mathcal M$ and a common cyclic and separating vector $\ket{\Omega}\in\mathcal{H}$ inducing the state $\omega$, the subsystem flow of the inclusion is
%     \begin{equation}
%         D_{N\subset M}(t):=\Delta_{\omega;M}^{-it}\Delta_{\omega;N}^{it},\qquad t\in\mathbb R.
%     \end{equation}
% \end{definition}
% For $b\in\mathcal N$, its adjoint action and the transported algebra are
% \begin{equation}\label{eq:characteristic-algebra-transport}
%  \operatorname{Ad}D_{N\subset M}(t)(b)=\sigma^{\omega;M}_{-t}\bigl(\sigma^{\omega;N}_t(b)\bigr),
%  \qquad \operatorname{Ad}D_{N\subset M}(t)(\mathcal N)=\sigma^{\omega;M}_{-t}(\mathcal N).
% \end{equation}
% This compares the intrinsic subsystem motion with the outer modular motion.
% It requires no conditional expectation. In general $D(t)$ is a cocycle rather than
% a one-parameter group; it need not belong to $\mathcal M$.
Characteristic flow compares the information accessible in the two algebras: the expectation value of its generator $K_M-K_N$ in a vector state is the sum of two nonnegative relative-entropy gaps, one for the inclusion $\mathcal N\subset\mathcal M$ and one for the dual inclusion $\mathcal M'\subset\mathcal N'$ (Proposition~\ref{prop:sufficiency}). This is a comparison at fixed state, not a dissipative entropy-production law.


Modular monotonicity gives bounded imaginary-time values $D_{N\subset M}(i\alpha)$ for $0\leq\alpha\leq1/2$. The subsystem flow extends to the strip $S(0,1/2):=\{z\in\mathbb C:0\leq\operatorname{Im}z\leq1/2\}$ \cite{revolutionizing}:
\begin{lemma}
    Given a pair of von Neumann algebras $\mathcal N\subset\mathcal M$ with a common cyclic and separating vector $\ket{\Omega}$, the subsystem flow $D_{N\subset M}(t)$ extends to the strip $\mathcal{S}(0,1/2)$. The extension is analytic in the open strip and satisfies $\|D(z)\|\leq1$ on its closure.

    Moreover, on the boundary $\partial\mathcal{S}(0,1/2) = \{z\in\mathbb{C}:\Im{z} = 0\}\sqcup\{z\in\mathbb{C}:\Im{z} = 1/2\}$, $D_{N\subset M}$ is strongly continuous.
\end{lemma}
For a full discussion we refer the reader to \cite{revolutionizing}, whose results we now summarize:
\begin{enumerate}
    \item $D(0)=1$, and $D_{N\subset M}(t)$ and $D_{N\subset M}(t +\frac{i}{2})$ are unitary and strongly continuous in $t$;
    \item $D_{N\subset M}(z)$ is analytic on $0<\operatorname{Im}z<1/2$ and contractive on its closure;
    \item $D_{N\subset M}(t)\ket{\Omega} = \ket{\Omega}$;
    \item The cocycle condition holds: $D(t+s)=\Delta_M^{-it}D(s)\Delta_M^{it}D(t)$. Here conjugation by $\Delta_M^{it}$ acts on $B(\mH)$, since $D(s)$ need not lie in $\mathcal M$;
    \item The intrinsic boundary symmetry is
    \[
     D(t+i/2)^*J_M D(t)=D(t)^*J_M D(t+i/2),
    \]
    independent of $t$. For the given inclusion its value is $J_N$, equivalently
    $D(t+i/2)=J_M D(t)J_N$. In particular, $D_{N\subset M}(\frac{i}{2}) = J_MJ_N$. This unitary operator is crucial in subfactor theory creating a chain of algebras:
    \begin{equation}
        \mathcal M\supset\mathcal N\supset D_{N\subset M}(\frac{i}{2})^*\mathcal MD_{N\subset M}(\frac{i}{2})\supset D_{N\subset M}(\frac{i}{2})^*\mathcal ND_{N\subset M}(\frac{i}{2})\supset...
    \end{equation}
    \item For all $t\in\mathbb{R}$, $D_{N\subset M}(t)D_{N\subset M}(\frac{i}{2})^*\mathcal MD_{N\subset M}(\frac{i}{2})D_{N\subset M}(t)^*\subset\mathcal M$ holds.
\end{enumerate}
Of properties 5 and 6 only the value $D_{N\subset M}(i/2)=J_MJ_N$ at the top of the strip is used below: the matrix elements of $D$ interpolate between the sufficiency gaps at $z=0$ and the modular conjugations at $z=i/2$ (Section~\ref{subsec:QEC-sufficiency}). The tower in property 5 is Borchers' route to the subfactor structure of the inclusion and is not needed here. The following theorem verifies the intuition that the subsystem flow fully captures the structure of an inclusion of von Neumann algebras with a common cyclic and separating vector~\cite{revolutionizing}:
\begin{theorem}
    Given a von Neumann algebra $\mathcal M$ and a cyclic and separating vector $\ket{\Omega}$, each subalgebra $\mathcal N\subset\mathcal M$ for which $\ket{\Omega}$ is cyclic is uniquely determined by its subsystem flow $D_{N\subset M}$, and every operator-valued function with the properties listed above is the subsystem flow of such a subalgebra.
\end{theorem}

\begin{lemma}[Common cyclicity implies generation]\label{lem:common-cyclicity-generating}
If $\mathcal N\subset\mathcal M$ has a common cyclic and separating vector $\Omega$, then
$\bigvee_{t\in\mathbb R}\sigma^{\omega;M}_t(\mathcal N)=\mathcal M$.
\end{lemma}
\begin{proof}
Let $\mP=\bigvee_t\sigma^{\omega;M}_t(\mathcal N)$. It is modular invariant, so Takesaki's theorem
gives an $\Omega$-preserving expectation $E_{\mP}:\mathcal M\to\mP$.
Since $\mathcal N\subset\mP$ and $\Omega$ is cyclic for $\mathcal N$, the projection onto
$\overline{\mP\Omega}$ is $1$. Hence $E_{\mP}(a)\Omega=a\Omega$ for every $a\in\mathcal M$.
Separatingness for $\mathcal M$ gives $E_{\mP}(a)=a$.
\end{proof}
The same argument shows that a proper common-cyclic inclusion has no
$\Omega$-preserving normal expectation onto $\mathcal N$. Formulated in the language of QEC, in an inclusion in a common cyclic and separating vector, the code subspace is the whole Hilbert space; i.e., $e_N=1$. Since such inclusions are generating, we have the following result:
% These statements do not assume ergodicity or mixing and do not establish the
% classical $K$-property, which additionally requires one-sided nesting and a trivial tail.
% The type $III_1$ local algebras in QFT provide the main examples; they are studied geometrically in the companion paper~\cite{companionQFT}.

\begin{corollary}\label{cor:common-cyclic-noncorrectable}
For the inclusion $\mN\subset \mM$ in a common cyclic and separating vector $\Omega$, no subalgebra of $\mM$ exists whose erasure \eqref{eq:erasure-channel} is exactly correctable.    
\end{corollary}

\begin{proof}
With $e_N=1$, (ii) would give $\mE(m)\Omega=m\Omega$ and, by separation, $\mE(m)=m$ for all $m$, forcing $\mathcal N=\mathcal M$ (Lemma~\ref{lem:common-cyclicity-generating} and the remark following it).
\end{proof}


In QFT, due to the Reeh-Schlieder theorem, inclusions of type III$_1$ algebras with common cyclic and separating vectors are commonplace. We saw that such inclusions do not admit normal expectations. In other words, there are no product states between the subalgebra and the relative commutant. This can be interpreted as the statement that in every state of QFT, there are correlations between $\mN$ and its relative commutant in $\mM$. For adjacent local algebras, by strong additivity, we expect that such inclusions are also conormal. QFT also provides with natural examples of split inclusions. The split property is a condition of statistical independence of $\mathcal N$ and $\mathcal M'$: the intermediate type $I$ factor gives a spatial tensor representation of $\mathcal N\vee\mathcal M'$~\cite{standardinclusion}, so states on $\mathcal N$ and $\mathcal M'$ can be prepared independently. It neither asserts that $\omega$ is a product state nor supplies a decay rate, and it is not a dynamical mixing condition.

% \paragraph{Modular monotonicity.}
% The common cyclic and separating vector controls the relative size of the two
% modular operators. For every $b\in\mathcal N$,
% \begin{equation}
%   S_{\omega;M}\,b\Omega
%   = J_{\omega;M}\Delta_{\omega;M}^{1/2}b\Omega
%   = b^*\Omega
%   = J_{\omega;N}\Delta_{\omega;N}^{1/2}b\Omega
%   = S_{\omega;N}\,b\Omega ,
%   \label{eq:Tomita-extension}
% \end{equation}
% so the Tomita operator of $\mathcal M$ extends that of $\mathcal N$. The associated closed
% quadratic forms are therefore ordered,
% \begin{equation}
%   \Delta_N\;\geq\;\Delta_M,
%   \label{eq:modular-inequality}
% \end{equation}
% and operator monotonicity of $z\mapsto z^\alpha$ for $0\le\alpha\le1$ gives
% $\Delta_N^{\alpha}\geq\Delta_M^{\alpha}$. For $0\le\alpha\le1/2$ this
% form order yields $\operatorname{Dom}\Delta_N^{\alpha}\subseteq\operatorname{Dom}\Delta_M^{\alpha}$,
% so $\Delta_M^{\alpha}\Delta_N^{-\alpha}$ is densely defined and closable, with
% \begin{equation}
%   \|\Delta_M^{\alpha}\Delta_N^{-\alpha}\|^2
%   =\|\Delta_N^{-\alpha}\Delta_M^{2\alpha}\Delta_N^{-\alpha}\|\le1 ,
%   \qquad 0\le\alpha\le\tfrac12 .
%   \label{eq:characteristic-bounded}
% \end{equation}
% This is the quantum statement with no classical counterpart: in
% Section~\ref{subsec:classical-state-flow} the Poisson Hamiltonian of a state
% has no analogue of the Tomita operator and satisfies no such order relation
% under coarse-graining. The bound \eqref{eq:characteristic-bounded} is what makes
% the following extension possible.


% The operator-algebraic error-correction reading of this dichotomy, in which a
% state-preserving expectation is the exactly correctable case and the subsystem
% flow measures the failure of correctability, is developed in
% Section~\ref{sec:QEC}.


% \begin{definition}[Standard inclusion]
% \label{def:standard-inclusion}
% An inclusion of von Neumann algebras $\mathcal N \subset \mathcal M$ is called \emph{standard} if there exists a vector $\Omega$ that is cyclic and separating for $\mathcal N$, $\mathcal M$, and the relative commutant $R = \mathcal N' \cap \mathcal M$.
% \end{definition}

\begin{lemma}[Standard inclusions of factors are generating and forgetful]
\label{lem:standard-forgetful}
Let $\mathcal N \subset \mathcal M$ be a standard inclusion of factors. Then the inclusion is generating and forgetful, meaning $\mathcal N_{\rm mem} = \mathbb C1$.
\end{lemma}
\begin{proof}
Since in a standard inclusion $\mN$, $\mM$ and the relative commutant $R = \mathcal N' \cap \mathcal M$ share a common cyclic and separating vector, we know that the inclusion is generating. Since $R \subset \mathcal M$, the modular orbit of $R$ generates the full ambient algebra, giving $R_{\rm hist} = \mathcal M$. 

By Lemma~\ref{lem:history-memory-duality}(b), the history of the relative commutant is contained within the commutant of the memory relative to $\mathcal M$:
\begin{equation}
    R_{\rm hist} \subseteq \mathcal N_{\rm mem}' \cap \mathcal M.
\end{equation}
Substituting $R_{\rm hist} = \mathcal M$ yields $\mathcal M \subseteq \mathcal N_{\rm mem}' \cap \mathcal M$, which forces the equality $\mathcal N_{\rm mem}' \cap \mathcal M = \mathcal M$. This equality dictates that every element of $\mathcal N_{\rm mem}$ commutes with every element of $\mathcal M$. 

Because $\mathcal N_{\rm mem}$ is a subalgebra of $\mathcal M$ (via $\mathcal N$), it must reside entirely within the center of $\mathcal M$:
\begin{equation}
    \mathcal N_{\rm mem} \subseteq \mathcal M' \cap \mathcal M = \mathcal Z(\mathcal M).
\end{equation}
Since $\mathcal M$ is a factor, its center is trivial ($\mathcal Z(\mathcal M) = \mathbb C1$). This restricts the memory algebra to $\mathcal N_{\rm mem} = \mathbb C1$, proving the inclusion is forgetful.
\end{proof}

\subsection{Modular ergodic hierarchy, chaos and emergent geometry}\label{subsec:QEC-hierarchy}

The properties relevant to this work can now be summarized in the language of modular ergodic theory. 
% Two channels appear in this paper and should be kept apart: the erasure of Section~\ref{subsec:QEC-erasure}, restriction of states to $\mathcal N$, is the \emph{noise}; the modular flow $\sigma^\omega_t$ is a reversible \emph{dynamics}. Information is lost only by their composition, ``evolve by $\sigma^\omega_t$, then restrict to $\mathcal N$'', and the right-hand column below reads each property for that composite channel.
\begin{center}
\begin{tabular}{lll}
 \textit{Property} & \textit{Modular formulation} & \textit{Error-correction reading}\\[4pt]
 Ergodicity & $\mathcal M^\omega=\mathbb C1$ & no observable is conserved by the dynamics\\
 Mixing & $C_{a,b}^\omega(t)\to0$ for all $a,b\in\mathcal M$ & connected correlations decay\\
 Autonomy of $\mathcal N$ & $\sigma^\omega_t(\mathcal N)=\mathcal N$ for all $t$ & $\mathcal N$ exactly correctable (Thm.~\ref{thm:correctable})\\
 Generation & $\mathcal N_{\rm hist}=\mathcal M$ & environment is forgetful (Cor.~\ref{cor:generating-forgetful})\\
 Forgetfulness & $\mathcal N_{\rm mem}=\mathbb C1$ & no noiseless subalgebra of $\mathcal N$ under the dynamics\\
 Half-sided inclusion & $\sigma^\omega_t(\mathcal N)\subseteq\mathcal N$ for one sign of $t$ & one-sided loss of access
\end{tabular}
\end{center}
Here $C^\omega_{a,b}(t)$ is the connected modular correlation function \eqref{eq:modular-correlation}; mixing requires its decay for every pair, is a property of the state on $\mathcal M$ rather than of the inclusion, and implies ergodicity (Lemma~\ref{lem:mixing-ergodic}). The classical $K$-property of Section~\ref{subsec:classical-ergodic-hierarchy} is the conjunction of half-sided inclusion, generation, and forgetfulness. Its modular counterpart is the \emph{modular $K$-system} built on Wiesbrock's half-sided modular inclusions~\cite{Wiesbrock1993}.

We define a quantum $K$-system as
\begin{definition}[Quantum $K$-system]
\label{def:quantum-K-system}
A quantum $K$-system is a proper inclusion of von Neumann algebras $\mathcal N \subset \mathcal M$ with a common cyclic and separating vector $\Omega$ that satisfies three conditions:
\begin{enumerate}
    \item \textbf{Half-sided nesting:} $\sigma^\omega_t(\mathcal N) \subseteq \mathcal N$ for $t \le 0$, where $\omega = \braket{\Omega | \cdot \, \Omega}$.
    \item \textbf{Generation:} $\mathcal N_{\rm hist} \equiv \bigvee_{t \in \mathbb R} \sigma^\omega_t(\mathcal N) = \mathcal M$.
    \item \textbf{Forgetfulness (Trivial Tail):} $\mathcal N_{\rm mem} \equiv \bigcap_{t \in \mathbb R} \sigma^\omega_t(\mathcal N) = \mathbb C1$.
\end{enumerate}
\end{definition}



We already saw that common cyclicity gives generation (Lemma~\ref{lem:common-cyclicity-generating}), and generation forces the environment to be forgetful (Corollary~\ref{cor:generating-forgetful}), and standard inclusions are generating and forgetful (Corollary~\ref{lem:standard-forgetful}). The remaining ingredient is half-sided inclusion. Classically, the one-sided inclusion \eqref{eq:classical-one-sided-inclusion} was imposed on an externally specified measure-preserving flow and had no consequences beyond those listed in Section~\ref{subsec:classical-ergodic-hierarchy}. In the quantum theory, the flow is not external; it is the modular flow of the state, and inclusion under it is a rigid condition. We show in this subsection that it produces a positive-energy translation group whose generator is the generator of the subsystem flow; that the modular flow transports this generator exponentially, at the universal rate $2\pi$ that saturates the modular chaos bound; and that one further hypothesis, standardness, produces an entire conformal group. None of these structures is put in by hand: a state and a nested pair of subalgebras suffice. This is the sense in which geometry emerges from modular chaos, and it is the operator-algebraic mechanism behind the emergence of bulk translations from boundary modular flows in holography~\cite{deBoerLamprou2020,LeutheusserLiu2023}.

\paragraph{Half-sided modular inclusions.}
\begin{definition}[Half-sided modular inclusion, modular $K$-system]\label{def:modular-K}\label{def:HSMI}
Let $\mathcal N\subset\mathcal M$ be a proper inclusion of von Neumann algebras with a common cyclic and separating vector $\Omega$, and $\omega=\braket{\Omega|\cdot\,\Omega}$. The pair $(\mathcal N\subset\mathcal M,\Omega)$ is a \emph{half-sided modular inclusion} if
\begin{equation}
 \sigma^{\omega}_t(\mathcal N)\subseteq\mathcal N\qquad\hbox{for all } t\le0 .
 \label{eq:HSMI}
\end{equation}
It is a \emph{modular $K$-system} if in addition its tail is trivial, $\mathcal N_{\rm mem}=\bigcap_{t\in\mathbb R}\sigma^\omega_t(\mathcal N)=\mathbb C1$.\footnote{In the physical orientation $\beta_t=\sigma^\omega_{-t}$ of Section~\ref{section:basics}, \eqref{eq:HSMI} reads $\beta_t(\mathcal N)\subseteq\mathcal N$ for $t\ge0$: it is the classical one-sided condition \eqref{eq:classical-one-sided-inclusion} with the external flow replaced by the modular flow of the state. This orientation agrees with the companion paper and with~\cite{Wiesbrock1993,revolutionizing}, where positive physical translations nest $\mathcal N$ into itself. Inclusions of the opposite orientation, $\sigma^\omega_t(\mathcal N)\subseteq\mathcal N$ for $t\ge0$, are obtained by passing to commutants, and every statement below has a mirror image for them.}
\end{definition}

Half-sided nesting with a common cyclic and separating vector is stronger than it looks. Generation, $\mathcal N_{\rm hist}=\mathcal M$, is automatic by Lemma~\ref{lem:common-cyclicity-generating}, and so, by Corollary~\ref{cor:generating-forgetful}, is forgetfulness of the environment $\mathcal R=\mathcal N'\cap\mathcal M$. Forgetfulness of $\mathcal N$ itself, the trivial tail, is not automatic, and we return to it below. What nesting adds is the following theorem of Wiesbrock and Borchers, which we state in the notation of this paper.

\begin{theorem}[Wiesbrock, Borchers~{\cite{Wiesbrock1993,Borchers1992,ArakiZsido2005}}]\label{thm:HSMI}
Let $(\mathcal N\subset\mathcal M,\Omega)$ be a half-sided modular inclusion. There is a unique strongly continuous unitary group $T(x)=e^{ixP}$ on $\mathcal H$ with the following properties.
\begin{enumerate}
 \item[\rm(i)] \emph{Positive energy, invariant vacuum:} $P\ge0$ and $T(x)\Omega=\Omega$.
 \item[\rm(ii)] \emph{Half-sided translations:} $T(x)\mathcal MT(-x)\subseteq\mathcal M$ for $x\ge0$, and $\mathcal N=T(1)\mathcal MT(-1)$.
 \item[\rm(iii)] \emph{Modular covariance and detailed balance:}
 \begin{equation}
  \Delta_M^{it}T(x)\Delta_M^{-it}=\Delta_N^{it}T(x)\Delta_N^{-it}=T(e^{-2\pi t}x),
  \qquad
  J_MT(x)J_M=J_NT(x)J_N=T(-x).
  \label{eq:Borchers-covariance}
 \end{equation}
 \item[\rm(iv)] \emph{The generator is the subsystem-flow generator:} $\Delta_N^{it}=T(1)\Delta_M^{it}T(-1)$, and, with $K_C=-\log\Delta_C$,
 \begin{equation}
  2\pi P=K_M-K_N
  \label{eq:HSMI-P}
 \end{equation}
 on $\operatorname{Dom}K_M\cap\operatorname{Dom}K_N$, the right side being understood as the closure of the difference on this domain~\cite{ArakiZsido2005}.
\end{enumerate}
Conversely, if $T(x)=e^{ixP}$ is a unitary group with $P\ge0$, $T(x)\Omega=\Omega$ and $T(x)\mathcal MT(-x)\subseteq\mathcal M$ for $x\ge0$, then \eqref{eq:Borchers-covariance} holds~\cite{Borchers1992} and $(T(1)\mathcal MT(-1)\subset\mathcal M,\Omega)$ is a half-sided modular inclusion.
\end{theorem}

The converse is the easy direction, and it shows how the orientation in \eqref{eq:HSMI} is fixed: $\Omega$ is cyclic and separating for $\mathcal N=T(1)\mathcal MT(-1)$ because $T(1)$ is unitary and fixes $\Omega$, and \eqref{eq:Borchers-covariance} gives
\begin{equation}
 \sigma^\omega_t(\mathcal N)=T(e^{-2\pi t})\,\mathcal M\,T(-e^{-2\pi t})
 =T(1)\,T(e^{-2\pi t}-1)\,\mathcal M\,T(1-e^{-2\pi t})\,T(-1)\subseteq\mathcal N
 \qquad(t\le0),
 \label{eq:HSMI-orbit}
\end{equation}
since $e^{-2\pi t}-1\ge0$ for $t\le0$. The hard direction is the existence of $T$ with the stated generator; the domain subtleties in (iv) were settled in~\cite{ArakiZsido2005}.

The theorem should be read as follows. Its input is a state and two nested algebras, with no geometry anywhere. Its output is a \emph{lightray}: a one-parameter group of translations with positive energy, with $\mathcal N$ the translate of $\mathcal M$ by one unit, and with a modular flow that acts on these translations as a boost, rescaling the translation parameter by $e^{-2\pi t}$. In the vacuum of a quantum field theory this is the Bisognano--Wichmann configuration: $\mathcal M$ is the algebra of a wedge $W$, $\mathcal N$ that of the wedge translated by one unit along a lightlike direction in its edge plane, $T$ the lightlike translations, $P$ the lightlike momentum, and \eqref{eq:Borchers-covariance} the statement that the vacuum modular flow of $W$ is the boost, under which lightlike momentum is rescaled. The content of Theorem~\ref{thm:HSMI} is that this configuration is not a feature of field theory: every half-sided modular inclusion carries an intrinsic lightray, extracted from the modular data alone.

\paragraph{Error-correction reading.}
Relation \eqref{eq:HSMI-P} identifies the emergent translation energy with the generator of the subsystem flow. On $\operatorname{Dom}K_M\cap\operatorname{Dom}K_N$,
\begin{equation}
 -i\,\partial_t D_{N\subset M}(t)\big|_{t=0}=K_M-K_N=2\pi P\ \ge\ 0 ,
 \label{eq:HSMI-generator}
\end{equation}
so by Proposition~\ref{prop:sufficiency} the lightlike energy of a state is a sum of two sufficiency gaps,
\begin{equation}
 2\pi\langle\Psi|P|\Psi\rangle
 =\bigl[S_M(\varphi\|\omega)-S_N(\varphi\|\omega)\bigr]
 +\bigl[S_{N'}(\varphi\|\omega)-S_{M'}(\varphi\|\omega)\bigr],
 \qquad \varphi=\braket{\Psi|\cdot\,\Psi},
 \label{eq:HSMI-energy-gaps}
\end{equation}
for $\Psi$ as in that proposition: the energy that the emergent translations assign to a state is the entropic cost of erasing the complement of $\mathcal N$, together with the dual cost for the commutants. A state is simultaneously recoverable with $\Omega$ from $\mathcal N$ and from $\mathcal M'$ only if it carries no translation energy, $P\Psi=0$. In the dictionary of Section~\ref{sec:QEC}, a half-sided modular inclusion is therefore the maximally non-correctable case with respect to the reference state: no expectation preserves $\omega$, the orbit of $\mathcal N$ generates $\mathcal M$, and the environment is forgetful. The first of these deserves its own proof, because it uses nothing but \eqref{eq:Borchers-covariance}.

\begin{proposition}\label{prop:no-vacuum-expectation}
Let $(\mathcal N\subset\mathcal M,\Omega)$ be a half-sided modular inclusion. There is no normal conditional expectation $\mE:\mathcal M\to\mathcal N$ with $\omega\circ\mE=\omega$.
\end{proposition}

\begin{proof}
By Takesaki's theorem (Theorem~\ref{thm:correctable}) such an expectation would give $\sigma^{\omega}_t(\mathcal N)=\mathcal N$ for all $t\in\mathbb R$. By \eqref{eq:HSMI-orbit}, $\sigma^\omega_t(\mathcal N)=T(e^{-2\pi t})\mathcal MT(-e^{-2\pi t})$, and $e^{-2\pi t}$ ranges over $(0,\infty)$. Invariance for all $t$ forces $T(x)\mathcal MT(-x)=\mathcal N$ for every $x>0$. For $m\in\mathcal M$ the elements $T(x)mT(-x)\in\mathcal N$ converge strongly to $m$ as $x\to0$, so $\mathcal M\subseteq\mathcal N$, contradicting properness.
\end{proof}

The proposition excludes expectations preserving the reference state only. Expectations preserving another faithful normal state or weight are not excluded by nesting; for a conormal inclusion they are excluded by Proposition~\ref{prop:conormal-factorization}, and for a standard inclusion that is neither split nor conormal no general no-expectation theorem is known. Lechner and Scotford~\cite[Section~2]{LechnerScotford2022} state that half-sided modular inclusions with unique vacuum admit no normal conditional expectation and that they are never split; the second assertion is the incompatibility of standard half-sided modular inclusions with the split property~\cite{Wiesbrock1993}, while the first, read as an unrestricted statement, requires an additional mechanism of the conormal type beyond Proposition~\ref{prop:no-vacuum-expectation}. The general criterion remains Takesaki's: a faithful normal expectation exists exactly when some faithful normal semifinite weight on $\mathcal M$, semifinite on $\mathcal N$, has a modular flow that preserves $\mathcal N$~\cite{Takesaki1972}.

\paragraph{Maximal modular chaos.}
Borchers' covariance relation is a statement about the subsystem flow. Differentiating \eqref{eq:Borchers-covariance} at $x=0$ gives, on the appropriate domains,
\begin{align}
 \Delta_M^{it}\,P\,\Delta_M^{-it}&=e^{-2\pi t}P,\qquad [K_M,P]=-2\pi i\,P,\nn\\
 \operatorname{Ad}\Delta_M^{-it}\bigl(K_M-K_N\bigr)&=e^{2\pi t}\bigl(K_M-K_N\bigr),
 \label{eq:modular-Lyapunov}
\end{align}
where the second form uses \eqref{eq:HSMI-P} and the physical modular flow $\beta_t=\operatorname{Ad}\Delta_M^{-it}$ of Section~\ref{section:basics}. The generator of the subsystem flow is thus an eigen-operator of the modular dynamics of $\mathcal M$ with exponent $2\pi$: under physical modular time the sum of the two sufficiency gaps \eqref{eq:HSMI-energy-gaps} grows as $e^{2\pi t}$. This is the maximal rate compatible with the analyticity of $D_{N\subset M}(z)$ on the strip $0\le\operatorname{Im}z\le1/2$ (Section~\ref{sec:QEC}), in the same way that the Maldacena--Shenker--Stanford bound follows from analyticity of thermal correlators on a strip~\cite{MSS2016}. The modular form of the bound is the modular chaos bound of de Boer and Lamprou~\cite{deBoerLamprou2020}; its derivation from the matrix elements of the subsystem flow is given in the companion paper~\cite{companionQFT}. Half-sided modular inclusions saturate it, and the saturation is characteristic: Borchers' theorem runs in the other direction as well, so that for a unitary group $T(x)=e^{ixP}$ fixing $\Omega$ and nesting $\mathcal M$ for $x\ge0$, positivity of $P$ is equivalent to the pair of relations \eqref{eq:Borchers-covariance}~\cite{revolutionizing}. In the terminology of~\cite{deBoerLamprou2020}, $P$ is a modular scrambling mode of $K_M$, and the subsystem flow $D_{N\subset M}(t)=e^{itK_M}e^{-itK_N}$ of a half-sided modular inclusion is its generating function. What is new relative to a merely generating inclusion is not that the modular orbit of $\mathcal N$ fills $\mathcal M$, which common cyclicity already guarantees, but that it does so at the maximal Lyapunov exponent, and that the quantity which grows is the failure of correctability, state by state.

\paragraph{Modular $K$-systems and the tail.}
Nesting and common cyclicity give generation and maximal chaos, but not a trivial tail. The tail of a half-sided modular inclusion has a transparent description. By \eqref{eq:HSMI-orbit},
\begin{equation}
 \mathcal N_{\rm mem}=\bigcap_{x>0}T(x)\mathcal MT(-x),
 \qquad
 (\mathcal N_{\rm mem})'=\bigvee_{x>0}T(x)\mathcal M'T(-x)=(\mathcal N')_{\rm hist},
 \label{eq:HSMI-tail}
\end{equation}
the second identity being Lemma~\ref{lem:history-memory}(a). Thus the tail of $\mathcal N$ is trivial if and only if the translates of the commutant $\mathcal M'$, which is nested the opposite way, $T(-x)\mathcal M'T(x)\subseteq\mathcal M'$ for $x\ge0$, generate all of $B(\mathcal H)$: forgetfulness of $\mathcal N$ is generation for the dual inclusion $\mathcal M'\subset\mathcal N'$ on the full Hilbert space. This can fail. If $(\mathcal N_0\subset\mathcal M_0,\Omega_0)$ is a half-sided modular inclusion and $\mathcal B$ is any factor with cyclic and separating vector $\Omega_{\mathcal B}$, then $(\mathcal N_0\overline\otimes\mathcal B\subset\mathcal M_0\overline\otimes\mathcal B,\ \Omega_0\otimes\Omega_{\mathcal B})$ is again a half-sided modular inclusion, with translations $T_0(x)\otimes1$, but its tail contains the spectator $1\otimes\mathcal B$, and $\Omega_0\otimes\eta$ is a translation-invariant vector for every $\eta$. The two deficiencies are the same: a nontrivial tail is a part of the system that the emergent translations never move. For the local algebras of a quantum field theory with unique vacuum the tail is trivial by locality. In a chiral net with $\mathcal M=\mathcal A((0,\infty))$ and $\mathcal N=\mathcal A((1,\infty))$, the tail $\bigcap_{x}\mathcal A((x,\infty))$ commutes with every local algebra $\mathcal A(I)$, since $I$ lies to the left of $(x,\infty)$ for $x$ large, and hence lies in the centre of the global algebra, which is trivial when the vacuum is unique. A modular $K$-system in the sense of Definition~\ref{def:modular-K} is therefore a half-sided modular inclusion whose emergent translations leave nothing behind, and uniqueness of the translation-invariant vector is its natural sufficient condition in field theory. With this understood, the classical statement that $K$-systems are the most chaotic members of the ergodic hierarchy acquires a quantitative modular counterpart:

\begin{corollary}[Modular $K$-systems are maximally chaotic~{\cite{OFLLM2024}}]\label{cor:maximal-chaos-K-system}
Let $(\mathcal N\subset\mathcal M,\Omega)$ be a half-sided modular inclusion. Then $\mathcal N_{\rm hist}=\mathcal M$, $\mathcal R_{\rm mem}=\mathbb C1$, and the subsystem-flow generator $K_M-K_N=2\pi P$ is transported by the modular flow of $\mathcal M$ at the exponent $2\pi$. If moreover $\mathcal N_{\rm mem}=\mathbb C1$, the pair is a modular $K$-system: it satisfies the three defining conditions of the classical $K$-property \eqref{eq:classical-K-algebra} for the physical modular flow $\beta_t=\sigma^\omega_{-t}$, and it saturates the modular chaos bound.
\end{corollary}
\begin{proof}
Generation is Lemma~\ref{lem:common-cyclicity-generating}, forgetfulness of the environment is Corollary~\ref{cor:generating-forgetful}, the transport relation is \eqref{eq:modular-Lyapunov}, and the nesting condition is \eqref{eq:HSMI} in the physical orientation.
\end{proof}

\paragraph{Standard inclusions and the emergent conformal group.}
Theorem~\ref{thm:HSMI} extracts one lightlike direction from a nested pair. Wiesbrock showed that one more hypothesis, standardness in the sense of Section~\ref{subsec:quantum-coarse}, extracts a whole conformal line~\cite{Wiesbrock1993b}, and Guido, Longo and Wiesbrock showed that nothing is lost in the process~\cite{GuidoLongoWiesbrock1998}. Recall that $(\mathcal N\subset\mathcal M,\Omega)$ is standard if $\Omega$ is cyclic and separating for $\mathcal N$, $\mathcal M$ and for the relative commutant $\mathcal R=\mathcal N'\cap\mathcal M$.

\begin{theorem}[Wiesbrock; Guido--Longo--Wiesbrock~{\cite{Wiesbrock1993b,GuidoLongoWiesbrock1998}}]\label{thm:wiesbrock-PSL}
Let $(\mathcal N\subset\mathcal M,\Omega)$ be a standard half-sided modular inclusion with translation group $T(x)=e^{ixP}$.
\begin{enumerate}
 \item[\rm(i)] $(\mathcal R\subset\mathcal M,\Omega)$ is a half-sided modular inclusion of the opposite orientation. Its nesting group is $\widetilde T(y)=e^{-iyQ}$ with $Q\ge0$, $\widetilde T(y)\Omega=\Omega$, $\mathcal R=\widetilde T(1)\mathcal M\widetilde T(-1)$,
 \begin{equation}
  2\pi Q=K_M-K_R,
  \qquad
  \Delta_M^{it}\,\widetilde T(y)\,\Delta_M^{-it}=\widetilde T(e^{2\pi t}y).
  \label{eq:HSMI-Q}
 \end{equation}
 \item[\rm(ii)] With $D:=K_M/2\pi$, the three self-adjoint generators satisfy, on a common core,
 \begin{equation}
  [D,P]=-iP,\qquad [D,Q]=iQ,\qquad [P,Q]=2iD,
  \label{eq:sl2-relations}
 \end{equation}
 and generate a positive-energy unitary representation of the universal covering group of $PSL(2,\mathbb R)$ fixing $\Omega$, in which $\Delta_M^{it}$ are the dilations, $T(x)$ the translations and $\widetilde T(y)$ the special conformal transformations of the line. The inversion is implemented by $J_MJ_{\widehat{\mathcal N}}$ with $\widehat{\mathcal N}:=\mathcal N\vee J_M\mathcal NJ_M$, and $Q=J_MJ_{\widehat{\mathcal N}}\,P\,J_{\widehat{\mathcal N}}J_M$.
 \item[\rm(iii)] The assignment $(0,\infty)\mapsto\mathcal M$, $(1,\infty)\mapsto\mathcal N$, $(0,1)\mapsto\mathcal R$ extends uniquely to a strongly additive, M\"obius-covariant local net of von Neumann algebras on the circle with vacuum vector $\Omega$, and every such net arises from a standard half-sided modular inclusion in this way.
\end{enumerate}
\end{theorem}

Two remarks make the theorem concrete. First, the sign in \eqref{eq:HSMI-Q} is forced: $K_M-K_R\ge0$ for every inclusion with a common cyclic and separating vector, by the modular monotonicity \eqref{eq:modular-inequality}, so the nesting group of an inclusion of the opposite orientation necessarily has a \emph{negative} generator $-Q$; the positive operator $Q$ is unitarily equivalent to $P$ by the inversion, as (ii) states. Under the physical modular flow the two translation energies are transported oppositely, $\beta_t(P)=e^{2\pi t}P$ and $\beta_t(Q)=e^{-2\pi t}Q$: the subsystem flow of $\mathcal N\subset\mathcal M$ is maximally growing while that of $\mathcal R\subset\mathcal M$ is maximally decaying, and the two together with $K_M$ close into $\mathfrak{sl}(2,\mathbb R)$. Second, in the net of (iii) the three algebras are the algebras of a half-line, of a translated half-line, and of the bounded interval between them; the emergent translations, dilations and special conformal transformations are literally the conformal transformations of the line, and the modular flows of all three algebras are geometric. The modular data of one standard nested pair thus determine a conformal field theory on the circle, together with its space-time symmetry.

\paragraph{Emergent geometry.}
The pattern of this subsection is the one by which spacetime is reconstructed from modular data. A single half-sided modular inclusion yields a lightray, its translations and its boost; a standard one yields a conformal line and the M\"obius group; and a finite family of half-sided modular inclusions in a suitable relative position, together with the modular intersections of the algebras involved, yields the Poincar\'e group of a higher-dimensional spacetime, with the modular flows of wedge algebras as its boosts~\cite{KahlerWiesbrock2001}. In each case the geometric quantity is a difference of modular Hamiltonians, that is, the generator of a subsystem flow, and the geometric transformation group is obtained by exponentiating these generators. The quantum field theories of the companion paper realize this pattern in the vacuum: the subsystem flows of wedges and balls are computed there explicitly, their generators are the lightlike momenta and conformal generators of Minkowski space, and the exponential dependence on $e^{2\pi z}$ that pervades those computations is the universal signature \eqref{eq:modular-Lyapunov} of half-sided nesting rather than an accident of free fields~\cite{companionQFT}.

The holographic version is the same mechanism with the roles of bulk and boundary added. In the large-$N$ limit the algebras of single-trace operators in a boundary time band are type $\mathrm{III}_1$ factors, nested time bands form half-sided modular inclusions, and the translation group of Theorem~\ref{thm:HSMI} is an emergent bulk translation that the boundary dynamics does not contain among its symmetries~\cite{LeutheusserLiu2023}. The modular chaos bound of~\cite{deBoerLamprou2020} is the strip analyticity of Section~\ref{sec:QEC} applied to a family of boundary subalgebras, and its saturation identifies the modular scrambling modes with bulk null translations along a Ryu--Takayanagi surface. In the present language these modes are the generators $K_M-K_N$ of the subsystem flows of nested boundary subalgebras, their growth at the exponent $2\pi$ is \eqref{eq:modular-Lyapunov}, and the sufficiency decomposition \eqref{eq:HSMI-energy-gaps} says what grows: the entropic cost of restricting a state to the smaller subalgebra, together with its dual. Emergent bulk geometry, in this reading, is the exponentiation of the failure of quantum error correction along a modular $K$-system.


\section{Discussion}\label{sec:discussion}

We have developed the subsystem flow as an object of ergodic theory and shown that it is the natural quantitative expression of quantum error correction for the dynamics generated by a state. Classically, every state on a Poisson manifold generates a flow, the relative transformation between the flows of a state and of its coarse-graining is a measure-preserving symplectic cocycle with generator $-\log(\rho/\rho_0)$, and its expectation is the Kullback--Leibler divergence, which is the classical sufficiency gap for the pair consisting of the state and the reference. Every classical coarse-graining is exactly correctable against erasure with respect to the reference, independently of whether the retained observables evolve autonomously. Quantum mechanically the two questions merge: by Takesaki's theorem, the subalgebras that admit a reference-preserving encoding correctable against erasure are exactly the modularly autonomous ones, and we proved that both are equivalent to triviality of the subsystem flow (Theorem~\ref{thm:correctable}). In the standard language of error correction this is a Knill--Laflamme condition with the Jones projection as code projection; for conormal inclusions it is literally the statement that the complementary channel into the relative commutant is forgetful on the code, and the information--disturbance trade-off survives in the modular setting at the level of the algebra but not pair by pair (Section~\ref{subsec:QEC-KL}). When the subsystem flow is not trivial, its generator, evaluated on a state, is the sum of two Petz sufficiency gaps (Proposition~\ref{prop:sufficiency}): the subsystem flow is a generating function for the failure of recoverability, with the second gap, absent classically, coming from the inclusion of commutants.

The modular ergodic hierarchy organizes the ways correctability fails. The history and memory algebras of a subalgebra are the smallest correctable algebra containing it and the largest correctable algebra inside it, and their exact duality (Lemma~\ref{lem:history-memory}) shows that a generating inclusion forces its environment to be forgetful, the modular form of decoupling. Half-sided modular inclusions are the maximally non-correctable case: no expectation preserves the reference, the orbit generates the full algebra, the environment is forgetful, the subsystem-flow generator is Borchers' positive translation generator, and it is transported by the modular flow at the exponent $2\pi$ that saturates the modular chaos bound. At the other end, $L^2$-nuclearity is the trace-class property of the subsystem flow at imaginary modular time $i/4$; it attaches to the inclusion a modular Hamiltonian $H_{N\subset M}$ with discrete spectrum, vacuum ground state, a partition function equal at $\beta=1/2$ to the nuclearity index, and a spectral form factor invariant under the modular flow of the subalgebra. Modular $K$-systems have infinite modular partition function, so the chaotic and the nuclear regimes are the two ends of one hierarchy, with real modular time governing Lyapunov growth and imaginary modular time governing the spectrum. This is the sense in which the ergodic hierarchy of a many-body system, formulated subsystem by subsystem, is a hierarchy of correctability.

Several distinctions have been kept separate because they are frequently identified. Common cyclicity of the reference vector gives generation and excludes a reference-preserving expectation; it does not give mixing, a decay rate for connected correlations, or a trivial one-sided tail. The split property is statistical independence of $\mathcal N$ and $\mathcal M'$ and excludes every normal expectation onto a type $III$ subalgebra, but it is not a dynamical condition. Conormality concerns $\mathcal N$ and $\mathcal R=\mathcal N'\cap\mathcal M$, and together with a faithful normal expectation it forces a spatial tensor factorization, which is what fails for adjacent local algebras. Finite index does not specify a matrix relative commutant, and infinite index does not force an infinite-dimensional one. The absence of a reference-preserving expectation does not exclude expectations preserving other faithful states, nor error correction for other channels or families of states; the general criterion remains Takesaki's.

Several directions remain open. The sufficiency decomposition refers to the generator of the subsystem flow; the full flow $D_{N\subset M}(z)$ on the strip should admit an interpretation in terms of a one-parameter family of recovery channels interpolating between the two modular conjugations, and its matrix elements should control approximate recoverability in the sense of the Fawzi--Renner and Junge--Renner--Sutter--Wilde--Winter inequalities. The modular Hamiltonian of an inclusion is one half of the subsystem-flow generator when the two modular operators commute; the relation between its spectrum and the sufficiency gaps in the general case, and the level statistics of that spectrum for interacting theories (Conjecture~\ref{conj:modular-SFF}), are not known. The compressed subsystem flow of Appendix~\ref{app:compressed-comparison} is a contraction whose conjugation action is completely positive but not in general unital; identifying the structure under which it is a channel on $\mathcal N$ would extend the dictionary beyond the common-cyclic case. For modular $K$-systems the relation between half-sided nesting and a trivial tail is not settled by common cyclicity, and the mechanism that excludes expectations preserving states other than the reference, established here for conormal adjacent regions, is not known for standard inclusions that are neither split nor conormal. In the companion paper~\cite{companionQFT} the subsystem flows of wedges and balls in vacuum quantum field theory are computed explicitly, the modular chaos bound is derived from their analytic structure, and both ends of the hierarchy are realized geometrically.

\section*{Acknowledgements}
We thank \ldots

\appendix

\section{Correlation coefficients for classical subspaces and subsystems}\label{app:classical-correlations}

This appendix records the measures of correlation between two coarse algebras of Section~\ref{subsec:classical-region-subsystem}, for a subspace pair and for a subsystem pair, and their relation to the total variation distance from the product measure. They quantify the distinction between the two kinds of coarse-graining but are not used in the main text.

\paragraph{Correlations between subspaces:} Write $C(f,g)=C_{f,g}(0)$ and similarly for its connected part. Of course, the definition of correlation above changes if we multiply the function by a constant. To account for this and define a measure of correlations between two subalgebras, we define
\begin{align}
    C(O_1,O_2)\equiv \sup_{\text{supp}(f_i)\in O_i,|f_i|_\infty\leq 1}|C(f_1,f_2)|
\end{align}
where $|f|_\infty=\sup_{x\in X}|f(x)|$. It follows from the definition above that 
\begin{align}
    C(O_1,O_2)\leq \int_{O_1\cap O_2} 1=\mu(O_1\cap O_2)=C(1_{O_1},1_{O_2})
\end{align}
Therefore, we have
\begin{align}
    C(O_1,O_2)=C(1_{O_1},1_{O_2})\ .
\end{align}

For connected correlations, define the event coefficient
\begin{equation}
 \alpha(O_1,O_2)=
 \sup_{\substack{A\subset O_1\\B\subset O_2}}
 |\mu(A\cap B)-\mu(A)\mu(B)|.
\end{equation}
In particular,
\begin{equation}
 C^{con}_{1_{O_1},1_{O_2}}(t)
 =\mu\bigl(O_1\cap\Phi_t(O_2)\bigr)-\mu(O_1)\mu(O_2).
\end{equation}
This measures the correlation of an event with a transported event.

\paragraph{Correlations between subsystems:}
Let $X=Y\times Z$ carry a joint probability measure $\mu_{YZ}$ with marginals
$\mu_Y,\mu_Z$. The unconnected correlation has supremum one over functions bounded
by one, because constant functions attain it. For connected correlations, distinguish
\begin{align}
 \alpha(Y,Z)&=\sup_{\substack{A\subset Y\\B\subset Z}}
 |\mu_{YZ}(A\times B)-\mu_Y(A)\mu_Z(B)|,\nn\\
 \Gamma(Y,Z)&=\sup_{\substack{f,g\ {\rm real}\\\|f\|_\infty,\|g\|_\infty\leq1}}
 |\mu_{YZ}(fg)-\mu_Y(f)\mu_Z(g)|=4\alpha(Y,Z).
\end{align}
The factor four follows by rescaling real functions from $[-1,1]$ to $[0,1]$
and writing the latter as integrals of indicator functions. With
$\nu=\mu_{YZ}-\mu_Y\otimes\mu_Z$, total variation satisfies
\begin{equation}
 \alpha(Y,Z)\leq\operatorname{TV}(\mu_{YZ},\mu_Y\otimes\mu_Z)
 :=\sup_E|\nu(E)|=\tfrac12\|\nu\|_1,
 \qquad \Gamma(Y,Z)\leq2\operatorname{TV}.
\end{equation}
Here $\|\nu\|_1$ is the total variation norm of the signed measure. All these
coefficients vanish precisely when the joint measure is a product.


\section{Infinite-dimensional quantum systems (von Neumann algebras)}

In this subsection, we review and fix our notation for general quantum systems (von Neumann algebras). Readers familiar with standard definitions can skip this subsection.

\paragraph{Algebra of quantum observables:}
Quantum observables are described by a 
\emph{von Neumann algebra} $\mathcal M$ a unital $*$-subalgebra of the bounded operators
$B(\mathcal{H})$ on some separable Hilbert space $\mathcal{H}$, closed in the
weak operator topology.
% The commutativity of the classical algebra is replaced by the requirement that $\mathcal M$
% be closed under the adjoint; it need not be abelian.  Indeed the classical theory is recovered
% by taking $\mathcal M=L^\infty(X,\Sigma,\mu)$, acting by multiplication on
% $\mathcal{H}=L^2(X,\Sigma,\mu)$.
A von Neumann algebra $\mathcal M$ is a \emph{factor} if its center
$\mathcal{Z}(\mathcal M) = \mathcal M\cap\mathcal M'$ consists only of scalar multiples
of the identity, where $\mathcal M' = \{b\in B(\mathcal{H}) : [b,a]=0\text{ for all }a\in\mathcal M\}$
is the commutant.  Physically, a factor algebra describes a system with no classical
superselection sectors \footnote{As we review in Section~\ref{sec:classical}, in classical physics, the commutator is replaced by the Poisson bracket and the commutant $\mathcal M'$ is the quantum analog of the Poisson
centralizer introduced in Section~\ref{sec:classical}. It consists of all observables invisible to $\mathcal M$. The classical analog of a factor is a non-degenerate Poisson bracket that restricts to a symplectic leaf.}.

% -------------------------------------------------------------------
\paragraph{Quantum states:}

A \emph{normal state} on $\mathcal M$ is a positive, normalized, $\sigma$-weakly continuous
linear functional $\varphi:\mathcal M\to\mathbb{C}$, with $\varphi(\id)=1$.
For a finite matrix algebra it has the form $\varphi(a)=\Tr(\rho_\varphi a)$,
where $\rho_\varphi$ is positive with unit trace. Type $III$ algebras have no
intrinsic tracial density matrix; the normal functional is the general definition.

A state $\varphi$ is \emph{faithful} if $\varphi(a^*a)=0$ implies $a=0$.
In finite dimensions this means that its density matrix has no zero eigenvalues.
Faithfulness is the counterpart of a strictly positive classical probability density
and gives a separating vector in the GNS construction.

A state $\varphi$ is \emph{tracial} if $\varphi(ab)=\varphi(ba)$ for all $a,b\in\mathcal M$.
On a finite algebra, a faithful normal tracial state plays the role of the reference
Liouville measure. Its modular flow is trivial, despite faithfulness. In the type $III$
algebras relevant to QFT, we use faithful normal states without a trace.

\paragraph{GNS Hilbert space and modular operators.}
Given a von Neumann algebra $\mathcal M$ and a faithful normal state $\varphi$, the
Gelfand--Naimark--Segal (GNS) construction produces a Hilbert space $\mathcal{H}_\varphi$,
a cyclic unit vector $|\Omega\rangle\in\mathcal{H}_\varphi$ with
$\varphi(a)=\langle\Omega|a|\Omega\rangle$, and a $*$-representation
$\pi:\mathcal M\to B(\mathcal{H}_\varphi)$ \footnote{The classical analog is $\mathcal{H}=L^2(X,\mu)$
with the constant function $|\id\rangle$ as the reference vector.}.

The Tomita--Takesaki Modular theory assigns to the pair $(\mathcal M,|\Omega\rangle)$ an
antilinear operator $S_\varphi$ defined on the dense domain
$\mathcal M|\Omega\rangle$ by
\begin{equation}
  S_\varphi\, a|\Omega\rangle = a^*|\Omega\rangle, \qquad a\in\mathcal M.
  \label{eq:Tomita-S}
\end{equation}
Its polar decomposition $S_\varphi = J_\varphi\,\Delta_\varphi^{1/2}$ defines the
\emph{modular conjugation} $J_\varphi$, an antiunitary involution, and the
\emph{modular operator} $\Delta_\varphi$, a positive self-adjoint operator. 

% -------------------------------------------------------------------
\paragraph{Modular dynamics:}
The Tomita--Takesaki theorem~\cite{Takesaki1970} guarantees that the one-parameter group
of unitaries $\Delta_\varphi^{it}$, $t\in\mathbb{R}$, implements an automorphism group of
$\mathcal M$:
\begin{equation}
  \sigma_t^\varphi(a) = \Delta_\varphi^{it}\, a\, \Delta_\varphi^{-it},
  \qquad a\in\mathcal M,\quad t\in\mathbb{R}.
  \label{eq:modular-flow}
\end{equation}
This is the \emph{modular flow} (or \emph{modular automorphism group}) of $\varphi$. 

The modular operator $\Delta_\varphi$ and $J_\varphi$ encode the non-commutativity of
$\mathcal M$. Section~\ref{subsec:classical-state-flow} describes a related
state-generated Poisson flow: for a smooth positive density $\rho$, the Hamiltonian
$K_\rho=-\log\rho$ defines a flow preserving that state. This construction uses
the Poisson bracket and differs from the trivial modular automorphism group of
an abelian von Neumann algebra. In both constructions the reference state is stationary:
$\varphi\circ\sigma_t^\varphi=\varphi$.

With the convention \eqref{eq:modular-flow}, the \emph{KMS condition} reads
\begin{equation}
 F_{a,b}(t)=\varphi\bigl(a\,\sigma_t^\varphi(b)\bigr),
 \qquad F_{a,b}(t-i)=\varphi\bigl(\sigma_t^\varphi(b)a\bigr).
 \label{eq:KMS-function}
\end{equation}
The function is bounded and analytic on $-1<\operatorname{Im}z<0$, with continuous
boundary values. Equivalently, $\beta_t^\varphi:=\sigma_{-t}^\varphi
=\operatorname{Ad}\Delta_\varphi^{-it}$ gives the upper-strip condition at inverse
temperature one. Its generator is $K_\varphi=-\log\Delta_\varphi$, while
$\partial_t\sigma_t^\varphi(b)|_{t=0}=i[\log\Delta_\varphi,b]$ on the generator domain.
The finite imaginary-time KMS shift cannot be replaced by this derivative.
The Poisson integration-by-parts identity \eqref{KMSclassical} is the classical
KMS condition; its relation to quantum KMS requires a semiclassical construction
and a specified time convention~\cite{vanDeVen2022}.


\section{Finite-dimensional coarse-grainings}\label{app:finite-coarse}

The tensor-product example \eqref{condexp} is one of three useful finite-dimensional
inclusions. Here the ambient algebra is $B(\mathcal H)$.
\begin{enumerate}
 \item {\bf Subspace inclusions:} Let $\mathcal H=P\mathcal H\oplus(1-P)\mathcal H$.
 The retained unital algebra is
 \[
  \mathcal N_P=P B(\mathcal H)P\oplus\mathbb C(1-P).
 \]
 Its center is $\mathbb CP\oplus\mathbb C(1-P)$. For a fixed state $\chi_{1-P}$
 on the complementary corner, the expectation is
 \[
  E_{\chi_{1-P}}(a)=PaP+
  \chi_{1-P}\bigl((1-P)a(1-P)\bigr)(1-P).
 \]
 It preserves precisely the block-diagonal states whose complementary block
 is proportional to the density matrix of $\chi_{1-P}$.
 Block diagonality alone does not suffice for a fixed expectation.
 For $P_i\leq P_{i+1}$, $\mathcal N_{P_i}\subseteq\mathcal N_{P_{i+1}}$.
 \item {\bf Subsystem inclusions:} If $\mathcal H=\mathcal H_B\otimes\mathcal H_R$,
 the retained algebra is $B(\mathcal H_B)\otimes1_R$.
 Its expectations and preserved states are those of \eqref{condexp}.
 \item {\bf Mixed inclusions:} A general unital finite-dimensional subalgebra has
 \[
  \mathcal H=\bigoplus_a\mathcal H_B^{(a)}\otimes\mathcal H_R^{(a)},\qquad
  \mathcal N=\bigoplus_a B(\mathcal H_B^{(a)})\otimes1_R^{(a)}.
 \]
 The center records the block label $a$. Bimodularity removes off-diagonal blocks,
 and a choice of normalized hidden state $\chi_R^{(a)}$ specifies a slice
 expectation in each block. The preserved states are
 \[
  \rho=\bigoplus_a p_a\,\rho_B^{(a)}\otimes\chi_R^{(a)}.
 \]
 With the label stored in a classical register $Z$, these have $I(B:R|Z)=0$.
 They maximize entropy at fixed retained state only when each hidden reference
 is maximally mixed; for a general reference, they are fixed states of that expectation.
\end{enumerate}
The subspace construction resolves alternatives among microstates, while a tensor
subsystem retains simultaneously specified degrees of freedom. The same distinction
organizes the classical constructions in Section~\ref{sec:classical}.


\section{Relative modular flow and relative entropy}
For faithful normal states $\varphi,\psi$ on a von Neumann algebra $\mathcal C$, represented by vectors $\Omega_\varphi,\Omega_\psi$ in standard form, the relative Tomita operator $S_{\psi|\varphi;C}$ is the closure of $c\,\Omega_\varphi\mapsto c^*\Omega_\psi$, $c\in\mathcal C$, and the relative modular operator is its positive square,
\begin{align}
 \Delta_{\psi|\varphi;C}=S_{\psi|\varphi;C}^\dagger S_{\psi|\varphi;C},
 \qquad
 \langle\Delta_{\psi|\varphi;C}^{1/2}a\Omega_\varphi|\Delta_{\psi|\varphi;C}^{1/2}b\Omega_\varphi\rangle=\psi(ba^*),\qquad a,b\in\mathcal C,
 \label{eq:relative-modular}
\end{align}
so that $\Delta_{\varphi|\varphi;C}=\Delta_{\varphi;C}$. Araki's relative entropy is
\begin{align}
 S_C(\varphi\|\psi)
 =-\langle\Omega_\varphi|\log\Delta_{\psi|\varphi;C}|\Omega_\varphi\rangle
 =i\left.\frac{d}{dt}\langle\Omega_\varphi|\Delta_{\psi|\varphi;C}^{it}|\Omega_\varphi\rangle\right|_{t=0},
 \label{eq:relative-entropy}
\end{align}
with the derivative understood when the logarithmic spectral moment is integrable~\cite{Araki1976,Araki1977}. In the matrix representation $\Delta_{\psi|\varphi}=L_{\rho_\psi}R_{\rho_\varphi}^{-1}$, and the classical counterpart is the Kullback--Leibler divergence \eqref{eq:KL-classical}. Relative entropy is nonnegative, vanishes exactly when $\varphi=\psi$, and is monotone under restriction to a subalgebra, which is what makes the gaps of Proposition~\ref{prop:sufficiency} nonnegative. For the coarse-grained state $\varphi\circ\mE$ of a faithful normal expectation $\mE$, $S_M(\varphi\|\varphi\circ\mE)$ is the entropy of the coarse-graining: it vanishes exactly when $\varphi=\varphi\circ\mE$, so an expectation preserving $\varphi$ has zero coarse-graining entropy. The classical identity $S(\rho\|\rho_0)=H(\rho_0)-H(\rho)$ of \eqref{eq:KL-classical} has a quantum counterpart only for a trace-preserving expectation.
%\section{The compressed subsystem flow}\label{app:compressed-comparison}

\bibliographystyle{unsrt}
\bibliography{subsystemflow}

\end{document}